% Generated mechanically from frozen input canon/ko-kr/integration-20260908-r10-p03-register/inputs/p03/ko/dga.tex.
% Do not edit this generated file; edit the bound locale source instead.

% OK, start here.
%

\setcounter{chapter}{21}
\chapter{미분 등급 대수}



\phantomsection
\label{dga-section-phantom}


\section{도입}
\label{dga-section-introduction}

\noindent
이 장에서는 미분 등급 대수, 가군, 범주 등을 다룬다. 기본 참고문헌은
\cite{Keller-Deriving}이고, 개관 논문은 \cite{Keller-survey}이다.

\medskip\noindent
Stacks 프로젝트에서는 서술의 길이를 염려하지 않으므로, 먼저 미분 등급
가군들의 범주라는 배경에서 내용을 전개한다. 그다음 미분 등급 범주 위의
미분 등급 가군이라는 배경에서 이 구성들을 다시 수행한다.



\section{규약}
\label{dga-section-conventions}

\noindent
이 장에서는 {\it 환}이 $1$을 갖는 가환환을 뜻한다는 규약을 유지한다.
$R$이 환이면, {\it $R$-대수 $A$}란 곱셈이라 불리는 $R$-쌍선형 사상
$A \times A \to A$가 주어져서 곱셈이 결합적이고 단위원을 갖는
$R$-가군 $A$를 뜻한다. 다시 말해, 이는 구조 사상 $R \to A$의 상이
$A$의 중심에 들어가는 단위원을 갖는 결합 $R$-대수이다.

\medskip\noindent
{\bf 부호 규칙.} 이 장에서는 흔히 미분을 갖춘 등급 대수와 등급 가군을
다룬다. 바탕 복합체의 부호 규칙은 항상 More on Algebra, Section
\ref{more-algebra-section-sign-rules}에서 도입한 규칙과 양립한다.
이 때문에 특히 왼쪽 가군이나 왼쪽 가군과 오른쪽 가군 사이의 관계를
고려할 때 곱셈 구조가 때때로 예상하지 못한 방식으로 꼬인다.





\section{미분 등급 대수}
\label{dga-section-dga}


\noindent
정의만 제시한다.

\begin{definition}
\label{dga-definition-dga}
$R$을 가환환이라 하자. {\it $R$ 위의 미분 등급 대수}란 다음 둘 중
하나이다.
\begin{enumerate}
\item $R$-쌍선형 사상 $A_n \times A_m \to A_{n + m}$,
$(a, b) \mapsto ab$가 주어진 $R$-가군의 사슬 복합체 $A_\bullet$로서
$$
\text{d}_{n + m}(ab) = \text{d}_n(a)b + (-1)^n a\text{d}_m(b)
$$
를 만족하고 $\bigoplus A_n$이 단위원을 갖는 결합 $R$-대수가 되는 것,
또는
\item $R$-쌍선형 사상 $A^n \times A^m \to A^{n + m}$,
$(a, b) \mapsto ab$가 주어진 $R$-가군의 쌍대사슬 복합체 $A^\bullet$로서
$$
\text{d}^{n + m}(ab) = \text{d}^n(a)b + (-1)^n a\text{d}^m(b)
$$
를 만족하고 $\bigoplus A^n$이 단위원을 갖는 결합 $R$-대수가 되는 것.
\end{enumerate}
\end{definition}

\noindent
흔히 $A = \bigoplus A_n$ 또는 $A = \bigoplus A^n$이라고만 쓰고, 이를
$\mathbf{Z}$-등급과 제곱이 영이며 위에서 설명한 라이프니츠 법칙을
만족하는 $R$-선형 연산자 $\text{d}$를 갖춘 단위원을 갖는 결합
$R$-대수로 생각한다. 이 경우 흔히 ``$(A, \text{d})$를 미분 등급
대수라 하자''라고 말한다.

\medskip\noindent
미분 등급 $R$-대수 $A$에서 미분과 곱셈을 잇는 라이프니츠 법칙은
곱셈 사상이 쌍대사슬 복합체의 사상
$$
\text{Tot}(A^\bullet \otimes_R A^\bullet) \to A^\bullet
$$
을 정의한다는 것과 정확히 같다. 여기서 $A^\bullet$은 $A$의 바탕
쌍대사슬 복합체를 나타낸다.

\begin{definition}
\label{dga-definition-homomorphism-dga}
{\it 미분 등급 대수의 준동형}
$f : (A, \text{d}) \to (B, \text{d})$란 등급 및 $\text{d}$와 양립하는
대수 준동형 $f : A \to B$이다.
\end{definition}

\begin{definition}
\label{dga-definition-cdga}
미분 등급 대수 $(A, \text{d})$는 차수가 $n$인 $a$와 차수가 $m$인
$b$에 대하여 $ab = (-1)^{nm}ba$이면 {\it 가환}이라고 한다.
이에 더하여 $\deg(a)$가 홀수일 때 $a^2 = 0$이면 $A$를
{\it 엄밀히 가환}이라고 한다.
\end{definition}

\noindent
다음 정의는 일반적으로도 의미가 있지만, 가환 미분 등급 대수들의
텐서곱을 취할 때에만 아마 ``올바른'' 정의일 것이다.

\begin{definition}
\label{dga-definition-tensor-product}
$R$을 환이라 하고, $(A, \text{d})$, $(B, \text{d})$를 $R$ 위의
미분 등급 대수라 하자. $A$와 $B$의 {\it 텐서곱 미분 등급 대수}란
곱셈이
$$
(a \otimes b)(a' \otimes b') = (-1)^{\deg(a')\deg(b)} aa' \otimes bb'
$$
로 정의된 대수 $A \otimes_R B$에 다음 규칙으로 정의되는 미분
$\text{d}$를 갖춘 것이다.
$\text{d}(a \otimes b) = \text{d}(a) \otimes b + (-1)^m a \otimes \text{d}(b)$
여기서 $m = \deg(a)$이다.
\end{definition}

\begin{lemma}
\label{dga-lemma-total-complex-tensor-product}
$R$을 환이라 하고, $(A, \text{d})$, $(B, \text{d})$를 $R$ 위의
미분 등급 대수라 하자. 그 바탕 쌍대사슬 복합체를 각각
$A^\bullet$, $B^\bullet$이라 쓰자. $R$-가군의 쌍대사슬 복합체로서
$$
(A \otimes_R B)^\bullet = \text{Tot}(A^\bullet \otimes_R B^\bullet).
$$
이다.
\end{lemma}

\begin{proof}
전체 복합체의 미분은 $A^p \otimes_R B^q$ 위에서
$\text{d}_1^{p, q} + (-1)^p \text{d}_2^{p, q}$로 주어진다는 것을
상기하자. 이는 정의 \ref{dga-definition-tensor-product}에서
$A \otimes_R B$ 위의 미분에 주어진 규칙과 정확히 같다.
\end{proof}






\section{미분 등급 가군}
\label{dga-section-modules}

\noindent
이 장에서는 오른쪽 가군을 기본으로 한다. 왼쪽 가군은 제
\ref{dga-section-left-modules}절에서 다룬다.

\begin{definition}
\label{dga-definition-dgm}
$R$을 환이라 하고, $(A, \text{d})$를 $R$ 위의 미분 등급 대수라 하자.
$A$ 위의 (오른쪽) {\it 미분 등급 가군} $M$이란 등급
$M = \bigoplus M^n$과 미분 $\text{d}$를 갖춘 오른쪽 $A$-가군
$M$으로서 $M^n A^m \subset M^{n + m}$ 및
$\text{d}(M^n) \subset M^{n + 1}$을 만족하고
$$
\text{d}(ma) = \text{d}(m)a + (-1)^n m\text{d}(a)
$$
가 $a \in A$와 $m \in M^n$에 대해 성립하는 것을 말한다.
{\it 미분 등급 가군의 준동형} $f : M \to N$이란 등급과 미분에
양립하는 $A$-가군 사상이다. (오른쪽) 미분 등급 $A$-가군의 범주를
$\text{Mod}_{(A, \text{d})}$로 나타낸다.
\end{definition}

\noindent
$M$을 (오른쪽) $R$-가군의 쌍대사슬 복합체 $M^\bullet$로 생각할 수
있음에 유의하자. 실제로 $r \in R$에 대해 $\text{d}(r) = 0$이고
$r$은 $A$의 차수 $0$인 원소로 사상되므로
$\text{d}(mr) = \text{d}(m)r$이다.

\medskip\noindent
미분 등급 $R$-대수 $A$ 위의 미분 등급 $R$-가군 $M$에서 미분과
곱셈을 잇는 라이프니츠 법칙은 곱셈 사상이 쌍대사슬 복합체의 사상
$$
\text{Tot}(M^\bullet \otimes_R A^\bullet) \to M^\bullet
$$
을 정의한다는 것과 정확히 같다. 여기서 $A^\bullet$과 $M^\bullet$은
각각 $A$와 $M$의 바탕 쌍대사슬 복합체를 나타낸다.

\begin{lemma}
\label{dga-lemma-dgm-abelian}
$(A, d)$를 미분 등급 대수라 하자. 범주
$\text{Mod}_{(A, \text{d})}$는 아벨 범주이고 임의의 극한과 여극한을 가진다.
\end{lemma}

\begin{proof}
핵과 여핵을 취하는 것은 바탕 $A$-가군을 취하는 것과 가환한다.
직합과 여극한도 마찬가지이다. 다시 말해 $\text{Mod}_{(A, \text{d})}$에서
이 연산들은 $A$-가군의 범주로 가는 망각 함자와 가환한다. 곱과 극한은
그렇지 않다. 즉, $N_i$, $i \in I$가 미분 등급 $A$-가군의 족이면
$\text{Mod}_{(A, \text{d})}$에서 곱 $\prod N_i$는
$(\prod N_i)^n = \prod N_i^n$ 및
$\prod N_i = \bigoplus_n (\prod N_i)^n$으로 주어진다. 따라서 곱은
등급 $A$-가군의 범주로 가는 망각 함자와 가환한다. 곱과 등화자를
갖는 범주는 극한을 가진다. Categories, Lemma
\ref{categories-lemma-limits-products-equalizers}를 보라.
\end{proof}

\noindent
따라서 $(A, \text{d})$가 $R$ 위의 미분 등급 대수이면 아벨 범주들의
완전 함자
$$
\text{Mod}_{(A, \text{d})} \longrightarrow \text{Comp}(R)
$$
가 존재한다. 미분 등급 가군 $M$의 코호몰로지군 $H^n(M)$은 이에
대응하는 $R$-가군 복합체의 코호몰로지로 정의한다. 그러므로 미분 등급
가군의 짧은 완전열 $0 \to K \to L \to M \to 0$은 코호몰로지 가군들의
긴 완전열
\begin{equation}
\label{dga-equation-les}
H^n(K) \to H^n(L) \to H^n(M) \to H^{n + 1}(K)
\end{equation}
을 유도한다. Homology, Lemma
\ref{homology-lemma-long-exact-sequence-cochain}를 보라.

\medskip\noindent
이제부터는 가군 복합체에 쓰는 모든 용어도 빌려 쓴다. 예를 들어 모든
$k \in \mathbf{Z}$에 대해 $H^k(M) = 0$이면 미분 등급 $A$-가군 $M$을
{\it 비순환}이라 한다. 미분 등급 $A$-가군의 준동형 $M \to N$이 모든
$k \in \mathbf{Z}$에 대해 동형 $H^k(M) \to H^k(N)$을 유도하면 이를
{\it 유사동형}이라 한다. 그 밖의 용어도 같은 방식으로 쓴다.

\begin{definition}
\label{dga-definition-shift}
$(A, \text{d})$를 미분 등급 대수라 하자. $M$을 바탕 $R$-가군 복합체가
$M^\bullet$인 미분 등급 가군이라 하자. 임의의 $k \in \mathbf{Z}$에 대해
{\it $k$-시프트 가군} $M[k]$를 다음과 같이 정의한다.
\begin{enumerate}
\item $M[k]$의 바탕 $R$-가군 복합체는 $M^\bullet[k]$이다. 즉,
$M[k]^n = M^{n + k}$이고 $\text{d}_{M[k]} = (-1)^k\text{d}_M$이다.
\item $A$-가군으로서의 곱셈
$$
(M[k])^n \times A^m \longrightarrow (M[k])^{n + m}
$$
은 주어진 곱셈 $M^{n + k} \times A^m \to M^{n + k + m}$과 같다.
\end{enumerate}
미분 등급 $A$-가군의 사상 $f : M \to N$에 대해
$f[k] : M[k] \to N[k]$를 바탕 $A$-가군 위에서 $f$와 같은 사상으로
정한다. 이는 함자
$[k] : \text{Mod}_{(A, \text{d})} \to \text{Mod}_{(A, \text{d})}$를
정의한다.
\end{definition}

\noindent
이 선택으로 라이프니츠 법칙이 성립함을 확인하자.
$x \in M[k]^n = M^{n + k}$, $a \in A^m$라 하고 $M[k]$의 곱을
$\cdot_{M[k]}$로 나타내면
\begin{align*}
\text{d}_{M[k]}(x \cdot_{M[k]} a)
& =
(-1)^k \text{d}_M(xa) \\
& =
(-1)^k \text{d}_M(x) a + (-1)^{k + n + k} x \text{d}(a) \\
& =
\text{d}_{M[k]}(x) a + (-1)^n x \text{d}(a) \\
& =
\text{d}_{M[k]}(x) \cdot_{M[k]} a + (-1)^n x \cdot_{M[k]} \text{d}(a)
\end{align*}
을 얻는다. $x$는 $M[k]$의 동차원소로서 차수가 $n$이므로 이것이 바로
원하는 식이다. 또한 이 선택 아래에서 곱셈 사상을 복합체의 사상
$$
\text{Tot}(M^\bullet[k] \otimes _R A^\bullet) \to
\text{Tot}(M^\bullet \otimes _R A^\bullet)[k] \to
M^\bullet[k]
$$
으로 생각할 수도 있다. 여기서 첫 번째 화살표는 More on Algebra, Section
\ref{more-algebra-section-sign-rules}의
(\ref{more-algebra-item-shift-tensor})이고, 이 경우에는 부호를 수반하지
않는다. (사실 이 관찰로부터 라이프니츠 법칙이 성립함을 유도할 수도 있었다.)

\medskip\noindent
Homology, Section \ref{homology-section-homotopy-shift}의 설명을 적용할 수
있다. 특히 모든 시프트 $M[k]$의 코호몰로지군을 부호의 개입 없이
동일시할 것이다.

\medskip\noindent
이제 {\it 삼각형}을 말할 수 있을 만큼 충분한 구조가 갖추어졌다.
Derived Categories, Definition \ref{derived-definition-triangle}을 보라.
다음 목표는 미분 등급 가군의 호모토피 범주가 삼각범주임을 진술하고
증명할 수 있을 만큼 이론을 전개하는 것이다. 먼저 호모토피 범주를 정의한다.






\section{호모토피 범주}
\label{dga-section-homotopy}

\noindent
우리의 호모토피는 $A$-가군 구조와 등급을 고려하지만, (물론) 미분은
고려하지 않는다.

\begin{definition}
\label{dga-definition-homotopy}
$(A, \text{d})$를 미분 등급 대수라 하자. $f, g : M \to N$을 미분 등급
$A$-가군의 준동형이라 하자. {\it $f$와 $g$ 사이의 호모토피}란 다음을
만족하는 $A$-가군 사상 $h : M \to N$이다.
\begin{enumerate}
\item 모든 $n$에 대해 $h(M^n) \subset N^{n - 1}$이고,
\item 모든 $x \in M$에 대해
$f(x) - g(x) = \text{d}_N(h(x)) + h(\text{d}_M(x))$이다.
\end{enumerate}
호모토피가 존재하면 $f$와 $g$가 {\it 호모토픽}이라고 한다.
\end{definition}

\noindent
따라서 $h$는 $A$-가군 구조 및 등급과 양립하지만 미분과는 양립하지 않는다.
$f = g$이고 $h$가 정의에서와 같은 호모토피이면, $h$는
$\text{Mod}_{(A, \text{d})}$ 안의 사상 $h : M \to N[-1]$을 정의한다.

\begin{lemma}
\label{dga-lemma-compose-homotopy}
$(A, \text{d})$를 미분 등급 대수라 하자. $f, g : L \to M$을 미분 등급
$A$-가군의 준동형이라 하고, 준동형 $a : K \to L$ 및 $c : M \to N$도
주어졌다고 하자. $h : L \to M$이 $f$와 $g$ 사이의 호모토피를 정의하는
$A$-가군 사상이면, $c \circ h \circ a$는 $c \circ f \circ a$와
$c \circ g \circ a$ 사이의 호모토피를 정의한다.
\end{lemma}

\begin{proof}
Homology, Lemma \ref{homology-lemma-compose-homotopy-cochain}에서 바로 따른다.
\end{proof}

\noindent
이 보조정리로 호모토피 범주를 다음과 같이 정의할 수 있다.

\begin{definition}
\label{dga-definition-complexes-notation}
$(A, \text{d})$를 미분 등급 대수라 하자. {\it 호모토피 범주}
$K(\text{Mod}_{(A, \text{d})})$란 그 대상이
$\text{Mod}_{(A, \text{d})}$의 대상이고, 그 사상이 미분 등급 $A$-가군
준동형의 호모토피류인 범주이다.
\end{definition}

\noindent
표기 $K(\text{Mod}_{(A, \text{d})})$는 표준적이지 않지만, 적어도 Stacks
프로젝트의 다른 곳에서 쓰는 $K(-)$와는 양립한다.

\begin{lemma}
\label{dga-lemma-homotopy-direct-sums}
$(A, \text{d})$를 미분 등급 대수라 하자. 호모토피 범주
$K(\text{Mod}_{(A, \text{d})})$는 직합과 곱을 가진다.
\end{lemma}

\begin{proof}
생략한다. 힌트: 보조정리 \ref{dga-lemma-dgm-abelian}에서와 같은 직합과 곱을
쓰면 된다. 이 함자들이 등급 $A$-가군의 범주로 가는 망각 함자와 가환함을
보았고, $\prod$가 아벨군의 족들로 이루어진 범주 위의 완전 함자이기 때문이다.
\end{proof}







\section{원뿔}
\label{dga-section-cones}

\noindent
미분 등급 가군의 범주에 원뿔을 도입한다.

\begin{definition}
\label{dga-definition-cone}
$(A, \text{d})$를 미분 등급 대수라 하자. $f : K \to L$을 미분 등급
$A$-가군의 준동형이라 하자. $f$의 {\it 원뿔}은
$C(f) = L \oplus K$, 등급 $C(f)^n = L^n \oplus K^{n + 1}$, 미분
$$
d_{C(f)} =
\left(
\begin{matrix}
\text{d}_L & f \\
0 & -\text{d}_K
\end{matrix}
\right)
$$
로 주어지는 미분 등급 $A$-가군 $C(f)$이다. 여기에는 자명한 사상
$L \to C(f)$와 $C(f) \to K$가 유도하는 복합체의 표준 사상
$i : L \to C(f)$와 $p : C(f) \to K[1]$이 갖추어져 있다.
\end{definition}

\noindent
원뿔 삼각형의 구성은 다음 의미에서 함자적이다.

\begin{lemma}
\label{dga-lemma-functorial-cone}
$(A, \text{d})$를 미분 등급 대수라 하자. 다음 도식이
$$
\xymatrix{
K_1 \ar[r]_{f_1} \ar[d]_a & L_1 \ar[d]^b \\
K_2 \ar[r]^{f_2} & L_2
}
$$
미분 등급 $A$-가군의 준동형들로 이루어지고 호모토피까지 가환한다고 하자.
그러면 $K(\text{Mod}_{(A, \text{d})})$ 안에서 삼각형의 사상
$$
(a, b, c) : (K_1, L_1, C(f_1), f_1, i_1, p_1) \to
(K_1, L_1, C(f_1), f_2, i_2, p_2)
$$
을 유도하는 사상 $c : C(f_1) \to C(f_2)$가 존재한다.
\end{lemma}

\begin{proof}
$h : K_1 \to L_2$를 $f_2 \circ a$와 $b \circ f_1$ 사이의 호모토피라
하자. $c$를 행렬
$$
c =
\left(
\begin{matrix}
b & h \\
0 & a
\end{matrix}
\right) :
L_1 \oplus K_1 \to L_2 \oplus K_2
$$
로 정의한다. 행렬 계산에 의해 $c$가 미분 등급 가군의 사상임을 알 수 있다.
$c \circ i_1 = i_2 \circ b$는 자명하고,
$p_2 \circ c = a \circ p_1$도 바로 확인된다.
\end{proof}











\section{허용 짧은 완전열}
\label{dga-section-admissible}

\noindent
허용 짧은 완전열은 미분 등급 가군의 상황에서 항별 분할 완전열에 대응하는
개념이다.

\begin{definition}
\label{dga-definition-admissible-ses}
$(A, \text{d})$를 미분 등급 대수라 하자.
\begin{enumerate}
\item 미분 등급 $A$-가군의 준동형 $K \to L$에 대해, $K \to L$의 왼쪽
역인 등급 $A$-가군 사상 $L \to K$가 존재하면 이를 {\it 허용 단사상}이라 한다.
\item 미분 등급 $A$-가군의 준동형 $L \to M$에 대해, $L \to M$의 오른쪽
역인 등급 $A$-가군 사상 $M \to L$이 존재하면 이를 {\it 허용 전사상}이라 한다.
\item 미분 등급 $A$-가군의 짧은 완전열 $0 \to K \to L \to M \to 0$이
등급 $A$-가군의 열로서 분할되면 이를 {\it 허용 짧은 완전열}이라 한다.
\end{enumerate}
\end{definition}

\noindent
따라서 이 분할들은 미분을 제외한 모든 자료와 양립한다. 허용 짧은 완전열이
주어지면 삼각형을 얻는다. 이것이 분할에 $A$-가군 구조와의 양립성을 요구하는
이유이다.

\begin{lemma}
\label{dga-lemma-admissible-ses}
$(A, \text{d})$를 미분 등급 대수라 하자.
$0 \to K \to L \to M \to 0$을 미분 등급 $A$-가군의 허용 짧은 완전열이라
하자. $s : M \to L$과 $\pi : L \to K$를
$\Ker(\pi) = \Im(s)$를 만족하는 분할이라 하자. 그러면
$\text{Mod}_{(A, \text{d})}$ 안의 사상
$$
\delta = \pi \circ \text{d}_L \circ s : M \to K[1]
$$
을 얻으며, 이는 코호몰로지 긴 완전열 (\ref{dga-equation-les})의 경계 사상들을
유도한다.
\end{lemma}

\begin{proof}
사상 $\pi \circ \text{d}_L \circ s$는 구성에 의해 $A$-가군 구조 및 등급과
양립한다. Homology, Lemma
\ref{homology-lemma-ses-termwise-split-cochain}에 의해 미분과도 양립한다.
$R$을 $A$가 그 위의 미분 등급 대수인 환이라 하자. 사상들의 등식은
$R$-가군에 관한 명제이므로, Homology, Lemmas
\ref{homology-lemma-ses-termwise-split-cochain} 및
\ref{homology-lemma-ses-termwise-split-long-cochain}에서 결론이 따른다.
\end{proof}

\begin{lemma}
\label{dga-lemma-make-commute-map}
$(A, \text{d})$를 미분 등급 대수라 하자. 다음 도식을 생각하자.
$$
\xymatrix{
K \ar[r]_f \ar[d]_a & L \ar[d]^b \\
M \ar[r]^g & N
}
$$
이는 미분 등급 $A$-가군의 준동형들로 이루어지고 호모토피까지 가환한다고 하자.
\begin{enumerate}
\item $f$가 허용 단사상이면, $b$는 도식을 가환하게 만드는 준동형과
호모토픽이다.
\item $g$가 허용 전사상이면, $a$는 도식을 가환하게 만드는 사상과
호모토픽이다.
\end{enumerate}
\end{lemma}

\begin{proof}
$h : K \to N$을 $bf$와 $ga$ 사이의 호모토피라 하자. 즉,
$bf - ga = \text{d}h + h\text{d}$라 하자. $\pi : L \to K$가 $f$의 왼쪽
역인 등급 $A$-가군 사상이라고 가정하고
$b' = b - \text{d}h\pi - h\pi \text{d}$로 둔다.
$s : N \to M$이 $g$의 오른쪽 역인 등급 $A$-가군 사상이라고 가정하고
$a' = a + \text{d}sh + sh\text{d}$로 둔다. 계산은 생략한다.
\end{proof}

\begin{lemma}
\label{dga-lemma-make-injective}
$(A, \text{d})$를 미분 등급 대수라 하자. $\alpha : K \to L$을 미분 등급
$A$-가군의 준동형이라 하자. 다음 분해가 존재한다.
$$
\xymatrix{
K \ar[r]^{\tilde \alpha} \ar@/_1pc/[rr]_\alpha &
\tilde L \ar[r]^\pi & L
}
$$
이는 $\text{Mod}_{(A, \text{d})}$ 안에서 다음을 만족한다.
\begin{enumerate}
\item $\tilde \alpha$는 허용 단사상이고(정의
\ref{dga-definition-admissible-ses} 참조),
\item $\pi \circ s = \text{id}_L$이고 $s \circ \pi$가
$\text{id}_{\tilde L}$과 호모토픽인 사상 $s : L \to \tilde L$이 존재한다.
\end{enumerate}
\end{lemma}

\begin{proof}
증명은 Derived Categories, Lemma \ref{derived-lemma-make-injective}의
증명과 같다. 즉, $\tilde L = L \oplus C(1_K)$로 두고 원뿔 구성의 기본
성질들을 사용한다.
\end{proof}

\begin{lemma}
\label{dga-lemma-sequence-maps-split}
$(A, \text{d})$를 미분 등급 대수라 하자.
$L_1 \to L_2 \to \ldots \to L_n$을 미분 등급 $A$-가군의 합성 가능한
준동형들의 열이라 하자. $\text{Mod}_{(A, \text{d})}$ 안에 가환 도식
$$
\xymatrix{
L_1 \ar[r] &
L_2 \ar[r] &
\ldots \ar[r] &
L_n \\
M_1 \ar[r] \ar[u] &
M_2 \ar[r] \ar[u] &
\ldots \ar[r] &
M_n \ar[u]
}
$$
이 존재하여, 각 $M_i \to M_{i + 1}$은 허용 단사상이고 각
$M_i \to L_i$는 호모토피 동치이다.
\end{lemma}

\begin{proof}
$n = 1$인 경우에는 주장할 것이 없다. 보조정리
\ref{dga-lemma-make-injective}가 $n = 2$인 경우이다. $M_n$을 제외한 도식이
구성되었다고 가정하자. 합성 $M_{n - 1} \to L_{n - 1} \to L_n$에
보조정리 \ref{dga-lemma-make-injective}를 적용하면 원하는 분해
$M_{n - 1} \to M_n \to L_n$을 얻는다.
\end{proof}



\begin{lemma}
\label{dga-lemma-nilpotent}
$(A, \text{d})$를 미분 등급 대수라 하자.
$0 \to K_i \to L_i \to M_i \to 0$, $i = 1, 2, 3$을 미분 등급 $A$-가군의
허용 짧은 완전열이라 하자. $b : L_1 \to L_2$와 $b' : L_2 \to L_3$를
다음 두 도식이 호모토피까지 가환하도록 하는 미분 등급 가군의 준동형이라 하자.
$$
\vcenter{
\xymatrix{
K_1 \ar[d]_0 \ar[r] &
L_1 \ar[r] \ar[d]_b &
M_1 \ar[d]_0 \\
K_2 \ar[r] & L_2 \ar[r] & M_2
}
}
\quad\text{and}\quad
\vcenter{
\xymatrix{
K_2 \ar[d]^0 \ar[r] &
L_2 \ar[r] \ar[d]^{b'} &
M_2 \ar[d]^0 \\
K_3 \ar[r] & L_3 \ar[r] & M_3
}
}
$$
그러면 $b' \circ b$는 $0$과 호모토픽이다.
\end{lemma}

\begin{proof}
보조정리 \ref{dga-lemma-make-commute-map}에 의해 $b$와 $b'$를 각각 호모토픽인
사상으로 바꾸어 왼쪽 도식의 오른쪽 정사각형과 오른쪽 도식의 왼쪽
정사각형을 가환하게 할 수 있다. 다시 말해
$\Im(b) \subset \Im(K_2 \to L_2)$이고
$\Ker(b') \supset \Im(K_2 \to L_2)$이다. 따라서 가군의 사상으로서
$b' \circ b = 0$이다.
\end{proof}
















\section{특수 삼각형}
\label{dga-section-distinguished}

\noindent
다음 보조정리로 특수 삼각형을 얻는다.

\begin{lemma}
\label{dga-lemma-triangle-independent-splittings}
$(A, \text{d})$를 미분 등급 대수라 하자.
$0 \to K \to L \to M \to 0$을 미분 등급 $A$-가군의 허용 짧은 완전열이라
하자. 삼각형
\begin{equation}
\label{dga-equation-triangle-associated-to-admissible-ses}
K \to L \to M \xrightarrow{\delta} K[1]
\end{equation}
에서 $\delta$를 보조정리 \ref{dga-lemma-admissible-ses}와 같이 택하자. 이
삼각형은 $K(\text{Mod}_{(A, \text{d})})$ 안의 표준 동형까지 보조정리
\ref{dga-lemma-admissible-ses}에서 한 선택들에 무관하다.
\end{lemma}

\begin{proof}
$(s', \pi')$를 보조정리 \ref{dga-lemma-admissible-ses}에서와 같은 또 다른
분할의 선택이라 하자. 그러면 $\delta$와 $\delta'$가 호모토픽임을 보이겠다.
% P03-STACKS-E883: frozen degree claim retained; see sidecar.
유일한 차수 $-1$의 $A$-가군 준동형 $h : M \to K$와 $g : M \to K$에 대해
$s' = s + \alpha \circ h$ 및 $\pi' = \pi + g \circ \beta$로 쓸 수 있다.
그러면 $g = -h$이고 $g$는 $\delta$와 $\delta'$ 사이의 호모토피이다.
계산은 Homology, Lemma
\ref{homology-lemma-ses-termwise-split-homotopy-cochain}의 증명에 들어 있다.
\end{proof}

\begin{definition}
\label{dga-definition-distinguished-triangle}
$(A, \text{d})$를 미분 등급 대수라 하자.
\begin{enumerate}
\item $0 \to K \to L \to M \to 0$이 미분 등급 $A$-가군의 허용 짧은
완전열이면, {\it $0 \to K \to L \to M \to 0$에 딸린 삼각형}이란
$K(\text{Mod}_{(A, \text{d})})$의 삼각형
(\ref{dga-equation-triangle-associated-to-admissible-ses})을 말한다.
\item $K(\text{Mod}_{(A, \text{d})})$의 삼각형이 미분 등급 $A$-가군의 허용
짧은 완전열에 딸린 삼각형과 동형이면 이를 {\it 특수 삼각형}이라 한다.
\end{enumerate}
\end{definition}









\section{원뿔과 특수 삼각형}
\label{dga-section-cones-and-triangles}

\noindent
$(A, \text{d})$를 미분 등급 대수라 하자.
$f : K \to L$을 미분 등급 $A$-가군의 준동형이라 하자.
그러면 $(K, L, C(f), f, i, p)$는 삼각형
$$
K \to L \to C(f) \to K[1]
$$
을 $\text{Mod}_{(A, \text{d})}$에서, 따라서
$K(\text{Mod}_{(A, \text{d})})$에서도 이룬다. 일반적으로 원뿔은 특수 삼각형이
{\bf 아니지만}, 둘의 차이는 부호 또는 회전이다(어느 쪽으로 보아도 된다).
다음은 이를 정밀하게 나타내는 두 명제이다.

\begin{lemma}
\label{dga-lemma-rotate-cone}
$(A, \text{d})$를 미분 등급 대수라 하고,
$f : K \to L$을 미분 등급 가군의 준동형이라 하자.
삼각형 $(L, C(f), K[1], i, p, f[1])$는 허용 짧은 완전열
$$
0 \to L \to C(f) \to K[1] \to 0
$$
에 결부된 삼각형이며, 이 완전열은 $f$의 원뿔의 정의에서 나온다.
\end{lemma}

\begin{proof}
정의에서 바로 따른다.
\end{proof}

\begin{lemma}
\label{dga-lemma-rotate-triangle}
$(A, \text{d})$를 미분 등급 대수라 하자.
$\alpha : K \to L$과 $\beta : L \to M$이 미분 등급 $A$-가군의
허용 짧은 완전열
$$
0 \to K \to L \to M \to 0
$$
을 정의한다고 하자. $(K, L, M, \alpha, \beta, \delta)$를 이에 결부된
삼각형이라 하자. 그러면 삼각형
$$
(M[-1], K, L, \delta[-1], \alpha, \beta)
\quad\text{and}\quad
(M[-1], K, C(\delta[-1]), \delta[-1], i, p)
$$
은 서로 동형이다.
\end{lemma}

\begin{proof}
분할을 하나 선택하여 $L = K \oplus M$이라 쓰고, $\alpha$와 $\beta$를 각각
자연 포함 사상과 사영 사상으로 동일시한다. $\delta$의 구성에 의해
% P03-STACKS-E523: frozen differential name retained; see sidecar.
$$
d_B =
\left(
\begin{matrix}
d_K & \delta \\
0 & d_M
\end{matrix}
\right)
$$
를 얻는다. 한편 $\delta[-1] : M[-1] \to K$의 원뿔은
$C(\delta[-1]) = K \oplus M$으로 주어지며 그 미분은 위 행렬과 동일하다!
따라서 보조정리가 성립한다.
\end{proof}

\begin{lemma}
\label{dga-lemma-third-isomorphism}
$(A, \text{d})$를 미분 등급 대수라 하자.
$f_1 : K_1 \to L_1$과 $f_2 : K_2 \to L_2$를 미분 등급 $A$-가군의
준동형이라 하자. 다음을 생각하자.
% P03-STACKS-E524: frozen codomain tuple retained; see sidecar.
$$
(a, b, c) :
(K_1, L_1, C(f_1), f_1, i_1, p_1)
\longrightarrow
(K_1, L_1, C(f_1), f_2, i_2, p_2)
$$
이것을 $K(\text{Mod}_{(A, \text{d})})$에서 삼각형들의 임의의 사상이라 하자.
$a$와 $b$가 호모토피 동치이면 $c$도 호모토피 동치이다.
\end{lemma}

\begin{proof}
$a^{-1} : K_2 \to K_1$을 $K(\text{Mod}_{(A, \text{d})})$에서 $a$의 역인
미분 등급 $A$-가군의 준동형이라 하자. 마찬가지로
$b^{-1} : L_2 \to L_1$을 $K(\text{Mod}_{(A, \text{d})})$에서 $b$의 역인
미분 등급 $A$-가군의
준동형이라 하자. 등식 $f_1 \circ a^{-1} = b^{-1} \circ f_2$에
보조정리 \ref{dga-lemma-functorial-cone}를 적용하여 얻는 사상을
$c' : C(f_2) \to C(f_1)$이라 하자. $c \circ c'$과 $c' \circ c$이
$K(\text{Mod}_{(A, \text{d})})$에서 동형임을 보이면 충분하다.
따라서 다음을 보이면 된다. $K(\text{Mod}_{(A, \text{d})})$에서
$(1, 1, c) : (K, L, C(f), f, i, p)$와 같은 삼각형의 사상이 주어지면,
$c$는 $K(\text{Mod}_{(A, \text{d})})$에서 동형이다. 가정에 의해 도식
$$
\xymatrix{
L \ar[r] \ar[d]_1 &
C(f) \ar[r] \ar[d]_c &
K[1] \ar[d]_1 \\
L \ar[r] &
C(f) \ar[r] &
K[1]
}
$$
의 두 사각형은 호모토피까지 가환한다. $C(f)$의 구성에 의해 두 행은
허용 짧은 완전열을 이룬다. 따라서 보조정리 \ref{dga-lemma-nilpotent}에 의해
$K(\text{Mod}_{(A, \text{d})})$에서 $(c - 1)^2 = 0$이다.
그러므로 $c$는 $K(\text{Mod}_{(A, \text{d})})$에서 동형이고 그 역은 $2 - c$이다.
\end{proof}

\noindent
다음 보조정리는 호모토피 범주에서 원뿔로 주어지는 삼각형들의 모임과
특수 삼각형들의 모임이, 적어도 부호를 무시하면 동형까지 서로 같음을 보인다!

\begin{lemma}
\label{dga-lemma-the-same-up-to-isomorphisms}
$(A, \text{d})$를 미분 등급 대수라 하자.
\begin{enumerate}
\item 미분 등급 $A$-가군의 허용 짧은 완전열
$0 \to K \xrightarrow{\alpha} L \to M \to 0$
이 주어지면, 도식
$$
\xymatrix{
K \ar[r] \ar[d] & L \ar[d] \ar[r] &
C(\alpha) \ar[r]_{-p} \ar[d] & K[1] \ar[d] \\
K \ar[r]^\alpha & L \ar[r]^\beta &
M \ar[r]^\delta & K[1]
}
$$
이 $K(\text{Mod}_{(A, \text{d})})$에서 삼각형들의 동형을 정의하도록 하는
호모토피 동치 $C(\alpha) \to M$이 존재한다.
\item 복합체의 사상 $f : K \to L$이 주어지면, 삼각형들의 동형
$$
\xymatrix{
K \ar[r] \ar[d] & \tilde L \ar[d] \ar[r] &
M \ar[r]_{\delta} \ar[d] & K[1] \ar[d] \\
K \ar[r] & L \ar[r] &
C(f) \ar[r]^{-p} & K[1]
}
$$
이 존재한다. 여기서 위쪽 삼각형은 허용 짧은 완전열
$K \to \tilde L \to M$에 결부된 삼각형이다.
\end{enumerate}
\end{lemma}

\begin{proof}
(1)의 증명. $C(\alpha) = L \oplus K$이다. $L$로의 사영 다음에 $\beta$를
합성하여 $C(\alpha) \to M$을 정의한다. 합성
$K^{n + 1} \to L^{n + 1} \to M^{n + 1}$이 영이므로 이것은 미분 등급
가군의 사상을 정의한다. $\Ker(\pi) = \Im(s)$를 만족하는 분할
$s : M \to L$과 $\pi : L \to K$를 선택하고, 평소와 같이
$\delta = \pi \circ \text{d}_L \circ s$로 둔다. 호모토피 역으로
$(s , -\delta)$로 주어지는 $M \to C(\alpha)$를 취한다. $\delta^n$은
$\text{d} \circ s^n - s^{n + 1} \circ \text{d} = \alpha \circ \delta^n$,
을 만족하는 유일한 사상 $M^n \to K^{n + 1}$로 특징지어지므로, 이 사상은
미분과 양립한다. Homology, Lemma
\ref{homology-lemma-ses-termwise-split-cochain}의 증명을 보라.
% P03-STACKS-E525: frozen C(f) notation retained; see sidecar.
합성 $M \to C(f) \to M$은 항등 사상이다.
합성 $C(f) \to M \to C(f)$은 다음 사상과 같다.
$$
\left(
\begin{matrix}
s \circ \beta & 0 \\
-\delta \circ \beta & 0
\end{matrix}
\right)
$$
이 사상이 항등 사상과 호모토픽임을 보이기 위해 행렬
$$
\left(
\begin{matrix}
0 & 0 \\
\pi & 0
\end{matrix}
\right) :
C(\alpha) = L \oplus K
\to
L \oplus K = C(\alpha)
$$
로 주어지는 호모토피 $h : C(\alpha) \to C(\alpha)$를 사용한다.
다음 등식은 바로 확인할 수 있다.
$$
\left(
\begin{matrix}
1 & 0 \\
0 & 1
\end{matrix}
\right)
-
\left(
\begin{matrix}
s \\
-\delta
\end{matrix}
\right)
\left(
\begin{matrix}
\beta & 0
\end{matrix}
\right)
=
\left(
\begin{matrix}
\text{d} & \alpha \\
0 & -\text{d}
\end{matrix}
\right)
\left(
\begin{matrix}
0 & 0 \\
\pi & 0
\end{matrix}
\right)
+
\left(
\begin{matrix}
0 & 0 \\
\pi & 0
\end{matrix}
\right)
\left(
\begin{matrix}
\text{d} & \alpha \\
0 & -\text{d}
\end{matrix}
\right)
$$
(1)의 증명을 끝내려면 사상 $-p : C(\alpha) \to K[1]$(정의
\ref{dga-definition-cone} 참조)과 합성 $C(\alpha) \to M \to K[1]$이
호모토피까지 일치함을 보여야 한다. 이는 위에서 명백하다. 즉, 호모토피 역
$(s, -\delta) : M \to C(\alpha)$를 사용하여 대신 두 사상
$M \to K[1]$이 일치함을 확인할 수 있다. 그리고 원하는 대로
$p \circ (s, -\delta) = -\delta$임을 주목하라.

\medskip\noindent
(2)의 증명. $\tilde f : K \to \tilde L$, $s : L \to \tilde L$, 그리고
% P03-STACKS-E526: frozen retraction type retained; see sidecar.
$\pi : L \to L$을 보조정리 \ref{dga-lemma-make-injective}의 것이라 하자.
보조정리 \ref{dga-lemma-functorial-cone}와 \ref{dga-lemma-third-isomorphism}에 의해
삼각형 $(K, L, C(f), i, p)$와
$(K, \tilde L, C(\tilde f), \tilde i, \tilde p)$는 동형이다.
삼각형들의 동형은 합성할 수 있음을 주목하라. 따라서 $L$을 $\tilde L$로,
$f$를 $\tilde f$로 바꾸어도 된다. 다시 말해 $f$가 허용 단사상이라고
가정해도 된다. 이 경우 결과는 (1)에서 따른다.
\end{proof}







\section{호모토피 범주는 삼각범주이다}
\label{dga-section-homotopy-triangulated}

\noindent
먼저 이것이 전삼각범주임을 증명한다.

\begin{lemma}
\label{dga-lemma-homotopy-category-pre-triangulated}
$(A, \text{d})$를 미분 등급 대수라 하자. 자연 이동 함자들과 특수 삼각형들을 갖춘
호모토피 범주 $K(\text{Mod}_{(A, \text{d})})$는 전삼각범주이다.
\end{lemma}

\begin{proof}
TR1의 증명. 정의에 따라 특수 삼각형과 동형인 모든 삼각형은 특수 삼각형이다.
또한 $0 \to K \to K \to 0 \to 0$이 허용 짧은 완전열이므로 임의의 삼각형
$(K, K, 0, 1, 0, 0)$은 특수 삼각형이다. 마지막으로 미분 등급 $A$-가군의
임의의 준동형 $f : K \to L$이 주어지면, 보조정리
\ref{dga-lemma-the-same-up-to-isomorphisms}에 의해 삼각형
$(K, L, C(f), f, i, -p)$는 특수 삼각형이다.

\medskip\noindent
TR2의 증명. $(X, Y, Z, f, g, h)$를 삼각형이라 하고
$(Y, Z, X[1], g, h, -f[1])$가 특수 삼각형이라고 가정하자.
그러면 허용 짧은 완전열 $0 \to K \to L \to M \to 0$이 존재하여, 이에 결부된
삼각형 $(K, L, M, \alpha, \beta, \delta)$는
$(Y, Z, X[1], g, h, -f[1])$와 동형이다. 역으로 회전하면
$(X, Y, Z, f, g, h)$는
$(M[-1], K, L, -\delta[-1], \alpha, \beta)$와 동형이다.
보조정리 \ref{dga-lemma-rotate-triangle}에 의해 삼각형
$(M[-1], K, L, \delta[-1], \alpha, \beta)$는
$(M[-1], K, C(\delta[-1]), \delta[-1], i, p)$와 동형이다.
앞의 삼각형 동형에 $Y$ 위의 $-1$을 선합성하면,
$(X, Y, Z, f, g, h)$가
$(M[-1], K, C(\delta[-1]), \delta[-1], i, -p)$와 동형임을 얻는다.
따라서 보조정리 \ref{dga-lemma-the-same-up-to-isomorphisms}에 의해 이것은
특수 삼각형이다. 반대로 $(X, Y, Z, f, g, h)$가 특수 삼각형이라고 하자.
보조정리 \ref{dga-lemma-the-same-up-to-isomorphisms}에 의해 이는
$\text{Mod}_{(A, \text{d})}$의 어떤 사상 $f$에 대한 삼각형
$(K, L, C(f), f, i, -p)$와 동형이다. 그러면 회전한 삼각형
$(Y, Z, X[1], g, h, -f[1])$는
$(L, C(f), K[1], i, -p, -f[1])$와 동형이고, 이는 다시 삼각형
$(L, C(f), K[1], i, p, f[1])$와 동형이다.
보조정리 \ref{dga-lemma-rotate-cone}에 의해 이 삼각형은 특수 삼각형이다.
따라서 원하는 대로 $(Y, Z, X[1], g, h, -f[1])$는 특수 삼각형이다.

\medskip\noindent
TR3의 증명. $(X, Y, Z, f, g, h)$와 $(X', Y', Z', f', g', h')$를
% P03-STACKS-E527: frozen undefined category name retained; see sidecar.
$K(\mathcal{A})$의 특수 삼각형이라 하고, $a : X \to X'$와
$b : Y \to Y'$를 $f' \circ a = b \circ f$를 만족하는 사상이라 하자.
보조정리 \ref{dga-lemma-the-same-up-to-isomorphisms}에 의해
$(X, Y, Z, f, g, h) = (X, Y, C(f), f, i, -p)$이고
$(X', Y', Z', f', g', h') = (X', Y', C(f'), f', i', -p')$라고
가정해도 된다. 이제 $f, f', a, b$가 주는 가환 도식에 보조정리
\ref{dga-lemma-functorial-cone}를 적용하면 된다.
\end{proof}

\noindent
TR4를 일반적으로 증명하기 전에 특별한 경우를 먼저 증명한다.

\begin{lemma}
\label{dga-lemma-two-split-injections}
$(A, \text{d})$를 미분 등급 대수라 하자. $\alpha : K \to L$과
$\beta : L \to M$이 미분 등급 $A$-가군의 허용 단사상이라고 하자.
그러면 TR4가 성립하는 특수 삼각형
$(K, L, Q_1, \alpha, p_1, d_1)$, $(K, M, Q_2, \beta \circ \alpha, p_2, d_2)$,
그리고 $(L, M, Q_3, \beta, p_3, d_3)$가 존재한다.
\end{lemma}

\begin{proof}
$\pi_1 : L \to K$와 $\pi_3 : M \to L$을 각각 $\alpha$와 $\beta$의
왼쪽 역인 등급 $A$-가군의 준동형이라 하자. 그러면 $K \to M$도
왼쪽 역 $\pi_2 = \pi_1 \circ \pi_3$을 갖는 허용 단사상이다.
$K \to L$, $K \to M$, $L \to M$의 여핵을 각각 $Q_1$, $Q_2$,
$Q_3$라 쓰자. 그러면 등급 $A$-가군으로서
$Q_1 = \Ker(\pi_1)$, $Q_3 = \Ker(\pi_3)$, 그리고
$Q_2 = \Ker(\pi_2)$라는 동일시를 얻는다. 또한 등급 $A$-가군으로서
$L = K \oplus Q_1$이고 $M = L \oplus Q_3$이므로
$M = K \oplus Q_1 \oplus Q_3$이다. $\pi_2 = \pi_1 \circ \pi_3$이
$Q_1$과 $Q_3$ 모두에서 영임을 주목하라. 따라서 $Q_2 = Q_1 \oplus Q_3$이다.
다음 가환 도식을 생각하자.
$$
\begin{matrix}
0 & \to & K & \to & L & \to & Q_1 & \to & 0 \\
  &     & \downarrow &     & \downarrow &     & \downarrow  & \\
0 & \to & K & \to & M & \to & Q_2 & \to & 0 \\
  &     & \downarrow &     & \downarrow &     & \downarrow  & \\
0 & \to & L & \to & M & \to & Q_3 & \to & 0
\end{matrix}
$$
이 도식의 각 행은 허용 짧은 완전열이므로 정의에 따라 특수 삼각형을 정한다.
더욱이 위 도식의 아래쪽 화살표들은 선택한 분할들과 양립하므로 삼각형들의 사상
$$
(K \to L \to Q_1 \to K[1])
\longrightarrow
(K \to M \to Q_2 \to K[1])
$$
and
$$
(K \to M \to Q_2 \to K[1])
\longrightarrow
(L \to M \to Q_3 \to L[1]).
$$
을 정한다. 도식의 맨 아래 열의 분할 $Q_3 \to M$을 $M \to Q_2$와 합성하면
분할 완전열 $0 \to Q_1 \to Q_2 \to Q_3 \to 0$의 분할을 얻음을 주목하라.
이로부터 대응하는 특수 삼각형
$$
(Q_1 \to Q_2 \to Q_3 \to Q_1[1])
$$
의 사상 $\delta : Q_3 \to Q_1[1]$가 합성 $Q_3 \to L[1] \to Q_1[1]$와
같음을 쉽게 얻는다. 따라서 공리 TR4의 결론과 같은 구조를 얻는다.
\end{proof}

\noindent
이제 최종 결과를 제시한다.

\begin{proposition}
\label{dga-proposition-homotopy-category-triangulated}
$(A, \text{d})$를 미분 등급 대수라 하자. 자연 이동 함자들과 특수 삼각형들을 갖춘
미분 등급 $A$-가군의 호모토피 범주 $K(\text{Mod}_{(A, \text{d})})$는
삼각범주이다.
\end{proposition}

\begin{proof}
$K(\text{Mod}_{(A, \text{d})})$가 전삼각범주임을 알고 있다. 따라서 TR4만
증명하면 충분하며, 그 증명에는 Derived Categories, Lemma
\ref{derived-lemma-easier-axiom-four}를 사용할 수 있다. $K \to L$과
$L \to M$을 $K(\text{Mod}_{(A, \text{d})})$의 합성 가능한 사상이라 하자.
보조정리 \ref{dga-lemma-sequence-maps-split}에 의해 $K \to L$과 $L \to M$이
허용 단사상이라고 가정해도 된다. 이 경우 결과는 보조정리
\ref{dga-lemma-two-split-injections}에서 따른다.
\end{proof}











\section{왼쪽 가군}
\label{dga-section-left-modules}

\noindent
지금까지 말한 모든 내용에는 왼쪽 미분 등급 가군의 설정에서도 대응물이 있지만,
몇 가지 부호 규칙에 주의해야 한다.

\medskip\noindent
$(A, \text{d})$를 미분 등급 $R$-대수라 하자. 오른쪽 가군의 경우와 완전히
유사하게, {\it 왼쪽 미분 등급 $A$-가군} $M$을 등급
$M = \bigoplus M^n$과 미분 $\text{d}$를 갖는 왼쪽 $A$-가군 $M$으로 정의하되,
$A^n M^m \subset M^{n + m}$ 및 $\text{d}(M^n) \subset M^{n + 1}$을 만족하고
$$
\text{d}(am) = \text{d}(a) m + (-1)^{\deg(a)}a \text{d}(m)
$$
이 모든 동차 원소 $a \in A$와 $m \in M$에 대해 성립한다고 한다. 앞에서와
마찬가지로 이 라이프니츠 법칙은 곱셈이 복합체의 사상
$$
\text{Tot}(A^\bullet \otimes_R M^\bullet) \to M^\bullet
$$
을 정의한다는 뜻과 정확히 같다. 여기서 $A^\bullet$과 $M^\bullet$은 각각
$A$와 $M$의 밑에 놓인 $R$-가군 복합체를 나타낸다.

\begin{definition}
\label{dga-definition-opposite-dga}
$R$을 환이라 하고, $(A, \text{d})$를 $R$ 위의 미분 등급 대수라 하자.
{\it 반대 미분 등급 대수}란 등급 $R$-가군으로서 $A^{opp} = A$이고
$\text{d} = \text{d}$이며 곱셈이
$$
a \cdot_{opp} b = (-1)^{\deg(a)\deg(b)} b a
$$
로 주어지는 $R$ 위의 미분 등급 대수 $(A^{opp}, \text{d})$를 말한다.
여기서 $a, b \in A$는 동차 원소이다.
\end{definition}

\noindent
이는 원하는 대로 다음이 성립하므로 잘 정의된다.
\begin{align*}
\text{d}(a \cdot_{opp} b)
& =
(-1)^{\deg(a)\deg(b)} \text{d}(b a) \\
& =
(-1)^{\deg(a)\deg(b)} \text{d}(b) a +
(-1)^{\deg(a)\deg(b) + \deg(b)}b\text{d}(a) \\
& =
(-1)^{\deg(a)}a \cdot_{opp} \text{d}(b) + \text{d}(a) \cdot_{opp} b
\end{align*}
밑의 $R$-가군 복합체로 표현하면 이는 도식
$$
\xymatrix{
\text{Tot}(A^\bullet \otimes_R A^\bullet)
\ar[rrr]_-{A^{opp}\text{의 곱셈}}
\ar[d]_{\text{가환성 제약}} & & &
A^\bullet \ar[d]^{\text{id}} \\
\text{Tot}(A^\bullet \otimes_R A^\bullet)
\ar[rrr]^-{A\text{의 곱셈}} & & &
A^\bullet
}
$$
이 가환한다는 뜻이다. 여기서 $R$-가군 복합체의 대칭 모노이드 범주에 대한
가환성 제약은 More on Algebra, Section
\ref{more-algebra-section-sign-rules}에 주어져 있다.

\medskip\noindent
$(A, \text{d})$를 $R$ 위의 미분 등급 대수라 하고, $M$을 왼쪽 미분 등급
$A$-가군이라 하자. $M$을 곱셈 $\cdot_{opp}$이 규칙
$$
m \cdot_{opp} a = (-1)^{\deg(a)\deg(m)} a m
$$
으로 정의된 오른쪽 $A^{opp}$-가군으로 보았을 때 $M^{opp}$라 쓰겠다.
여기서 $a$와 $m$은 동차 원소이다. 이 곱셈은 미분과 양립한다. 실제로
$M^{opp}$ 위의 곱셈 $\cdot_{opp}$을 정의할 때 도식
$$
\xymatrix{
\text{Tot}(M^\bullet \otimes_R A^\bullet)
\ar[rrr]_-{M^{opp}\text{ 위의 곱셈}}
\ar[d]_{\text{가환성 제약}} & & &
M^\bullet \ar[d]^{\text{id}} \\
\text{Tot}(A^\bullet \otimes_R M^\bullet)
\ar[rrr]^-{M\text{ 위의 곱셈}} & & &
M^\bullet
}
$$
을 사용할 수도 있었다. 이것이 결합적인 곱셈임을 보이기 위해 동차 원소
$a, b \in A$와 $m \in M$에 대해 계산하면
\begin{align*}
m \cdot_{opp} (a \cdot_{opp} b)
& =
(-1)^{\deg(a)\deg(b)} m \cdot_{opp} (ba) \\
& =
(-1)^{\deg(a)\deg(b) + \deg(ab)\deg(m)} bam \\
& =
(-1)^{\deg(a)\deg(b) + \deg(ab)\deg(m) + \deg(b)\deg(am)}
(am) \cdot_{opp} b \\
& =
\begin{aligned}
&(-1)^{\deg(a)\deg(b) + \deg(ab)\deg(m) + \deg(b)\deg(am) + \deg(a)\deg(m)} \\
&(m \cdot_{opp} a) \cdot_{opp} b
\end{aligned} \\
& =
(m \cdot_{opp} a) \cdot_{opp} b
\end{align*}
를 얻는다. 물론 $R$-가군 복합체의 대칭 모노이드 범주에서 결합성 제약과
가환성 제약 사이의 양립성을 사용해서도 이를 보일 수 있었다.

\begin{lemma}
\label{dga-lemma-left-right}
$(A, \text{d})$를 미분 등급 $R$-대수라 하자. 왼쪽 미분 등급 $A$-가군의
범주에서 오른쪽 미분 등급 $A^{opp}$-가군의 범주로 가는 함자
$M \mapsto M^{opp}$는 동치이다.
\end{lemma}

\begin{proof}
생략한다.
\end{proof}

\noindent
다음으로 시프트를 다룬다. $(A, \text{d})$를 미분 등급 대수라 하고,
$M$을 그 밑의 $R$-가군 복합체를 $M^\bullet$이라 쓰는 왼쪽 미분 등급
$A$-가군이라 하자. 임의의 $k \in \mathbf{Z}$에 대해
{\it $k$-시프트 가군} $M[k]$를 다음과 같이 정의한다.
\begin{enumerate}
\item $M[k]$의 밑 $R$-가군 복합체는 $M^\bullet[k]$이다.
\item $A$-가군으로서의 곱셈
$$
A^n \times (M[k])^m \longrightarrow (M[k])^{n + m}
$$
은 주어진 곱셈 $A^n \times M^{m + k} \to M^{n + m + k}$의
$(-1)^{nk}$배이다.
\end{enumerate}
이 선택에서 라이프니츠 법칙이 성립함을 확인하자. $a \in A^n$이고
$x \in M[k]^m = M^{m + k}$라 하며, $M[k]$의 곱을 $\cdot_{M[k]}$로 쓰면
\begin{align*}
\text{d}_{M[k]}(a \cdot_{M[k]} x)
& =
(-1)^{k + nk} \text{d}_M(ax) \\
& =
(-1)^{k + nk} \text{d}(a) x + (-1)^{k + nk + n} a \text{d}_M(x) \\
& =
\text{d}(a) \cdot_{M[k]} x + (-1)^{nk + n} a \text{d}_{M[k]}(x) \\
& =
\text{d}(a) \cdot_{M[k]} x + (-1)^n a \cdot_{M[k]} \text{d}_{M[k]}(x)
\end{align*}
를 얻는다. $a$는 $A$의 동차 원소로서 차수 $n$을 가지므로 이것이 원하는
식이다. 또한 이 선택 아래에서 곱셈 사상을 복합체의 사상
$$
\text{Tot}(A^\bullet \otimes_R M^\bullet[k]) \to
\text{Tot}(A^\bullet \otimes_R M^\bullet)[k] \to
M^\bullet[k]
$$
으로 볼 수 있다. 여기서 첫째 화살표는 More on Algebra, Section
\ref{more-algebra-section-sign-rules}의
(\ref{more-algebra-item-shift-tensor})이며, 이 경우에는 위에서 선택한 부호가
정확히 나타난다. (사실 이 관찰에서 라이프니츠 법칙이 성립함을 도출할 수도 있었다.)

\medskip\noindent
위 규칙에 의해 오른쪽 미분 등급 $A^{opp}$-가군들의 표준적 동일시
$$
(M[k])^{opp} = M^{opp}[k]
$$
를 부호의 개입 없이 얻는다. 다시 말해 보조정리 \ref{dga-lemma-left-right}의
동치는 시프트 함자들과 양립한다.

\medskip\noindent
위 선택 때문에 다음 정의가 필요하다.

\begin{definition}
\label{dga-definition-shift-graded-module}
$R$을 환이라 하고, $A$를 $\mathbf{Z}$-등급 $R$-대수라 하자.
\begin{enumerate}
\item 오른쪽 등급 $A$-가군 $M$이 주어지면, {\it 제 $k$ 시프트
$A$-가군} $M[k]$를 오른쪽 $A$-가군으로서는 같고 등급이
$(M[k])^n = M^{n + k}$인 것으로 정의한다.
\item 왼쪽 등급 $A$-가군 $M$이 주어지면, {\it 제 $k$ 시프트
$A$-가군} $M[k]$를 등급 $(M[k])^n = M^{n + k}$을 가지며 곱셈
$A^n \times (M[k])^m \to (M[k])^{n + m}$
이 주어진 곱셈 $A^n \times M^{m + k} \to M^{n + m + k}$의
$(-1)^{nk}$배인 가군으로 정의한다.
\end{enumerate}
\end{definition}

\noindent
$(A, \text{d})$를 미분 등급 대수라 하고, $f, g : M \to N$을 왼쪽
미분 등급 $A$-가군의 준동형이라 하자. {\it $f$와 $g$ 사이의 호모토피}란
다음을 만족하는 등급 $A$-가군 사상 $h : M \to N[-1]$(시프트에 주의하라!)이다.
$$
f(x) - g(x) = \text{d}_N(h(x)) + h(\text{d}_M(x))
$$
이 식은 모든 $x \in M$에 대해 성립한다. 호모토피가 존재하면 $f$와 $g$가
{\it 호모토픽}이라고 한다. 따라서 $h$는 $A$-가군 구조($N$에서는 시프트된
구조) 및 등급($N$에서는 시프트된 등급)과 양립하지만 미분과는 양립하지 않는다.
$f = g$이고 $h$가 호모토피이면, $h$는 왼쪽 미분 등급 $A$-가군의 사상
$h : M \to N[-1]$을 정의한다.

\medskip\noindent
위 규칙에 의해 $f, g : M \to N$이 호모토픽인 것은 유도된 사상
$f^{opp}, g^{opp} : M^{opp} \to N^{opp}$가 오른쪽 미분 등급
$A^{opp}$-가군의 준동형으로서 (같은 호모토피에 의해) 호모토픽인 것과 동치이다.

\medskip\noindent
호모토피 범주, 원뿔, 허용 짧은 완전열, 특수 삼각형은 모두 오른쪽 미분 등급
가군의 경우와 정확히 같은 방식으로 정의한다(그리고 밑 $R$-가군 복합체 위에서
모든 것은 $R$-가군 복합체에 대한 구성과 일치한다). 이 방식으로 왼쪽 가군에
대해서도 명제 \ref{dga-proposition-homotopy-category-triangulated}의 대응물을 얻거나,
반대대수 위의 오른쪽 가군을 다루어 이를 도출할 수 있다.


\section{텐서곱}
\label{dga-section-tensor-product}

\noindent
$R$을 환이라 하자. $A$를 $R$-대수라 하자
(절 \ref{dga-section-conventions} 참조). 오른쪽 $A$-가군 $M$과
왼쪽 $A$-가군 $N$이 주어지면 {\it 텐서곱}
$$
M \otimes_A N
$$
이 있다. 이 텐서곱은 $R$-가군이다. $R$-가군으로서
$M \otimes_A N$은 $x \in M$ 및 $y \in N$에 대한 기호
$x \otimes y$들로 생성되며 다음 관계들을 만족한다.
$$
\begin{matrix}
(x_1 + x_2) \otimes y - x_1 \otimes y - x_2 \otimes y, \\
x \otimes (y_1 + y_2) - x \otimes y_1 - x \otimes y_2, \\
xa \otimes y - x \otimes ay
\end{matrix}
$$
여기서 $a \in A$, $x, x_1, x_2 \in M$이고
$y, y_1, y_2 \in N$이다. 텐서곱의 몇 가지 성질을 열거하자.

\medskip\noindent
텐서곱은 각 변수에 대해 우완전이며, 실제로 직합 및 임의의
여극한과 가환한다.

\medskip\noindent
텐서곱 $M \otimes_A N$은 보편 $A$-쌍선형 사상
$M \times N \to M \otimes_A N$, $(x, y) \mapsto x \otimes y$의
보편 대상이다. 공식으로 쓰면 임의의 $R$-가군 $Q$에 대해
$$
\text{쌍선형}_A(M \times N, Q) = \Hom_R(M \otimes_A N, Q)
$$
이다.

\medskip\noindent
$A$가 $\mathbf{Z}$-등급 대수이고 $M$, $N$이 등급
$A$-가군이면, $M \otimes_A N$은 등급 $R$-가군이다.
이때 $M \otimes_A N$의 $n$번째 등급 성분
$(M \otimes_A N)^n$은
$$
\Coker\left(
\bigoplus\nolimits_{r + t + s = n}
M^r \otimes_R A^t \otimes_R N^s \to
\bigoplus\nolimits_{p + q = n} M^p \otimes_R N^q
\right)
$$
와 같다. 여기서 사상은 $r + s + t = n$인
$x \in M^r$, $y \in N^s$, $a \in A^t$에 대해
$x \otimes a \otimes y$를 $x \otimes ay - xa \otimes y$로 보낸다.
이 경우 사상 $M \times N \to M \otimes_A N$은
$A$-쌍선형이고 등급과 양립하며, 임의의 등급 $R$-가군 $Q$에 대해
$$
\text{등급 쌍선형}_A(M \times N, Q) =
\Hom_{\text{등급 }R\text{-가군}}(M \otimes_A N, Q)
$$
라는 의미에서 보편적이다. 여기서는 자명한 의미의 등급 쌍선형
사상을 사용한다.

\medskip\noindent
% P03-STACKS-E533: source handedness is preserved; see the sidecar ledger.
$(A, \text{d})$가 미분 등급 대수이고 $M$과 $N$이 각각 왼쪽 및
오른쪽 미분 등급 $A$-가군이면, $M \otimes_A N$은 다음 미분을 갖는
미분 등급 $R$-가군이다.
$$
\text{d}(x \otimes y) =
\text{d}(x) \otimes y + (-1)^{\deg(x)}x \otimes \text{d}(y)
$$
여기서 $x \in M$과 $y \in N$은 동차 원소이다. 이 경우 사상
$M \times N \to M \otimes_A N$은 $A$-쌍선형이고 등급 및 미분과
양립하며, 임의의 미분 등급 $R$-가군 $Q$에 대해
$$
\text{미분 등급 쌍선형}_A(M \times N, Q) =
\Hom_{\text{Comp}(R)}(M \otimes_A N, Q)
$$
라는 의미에서 보편적이다. 여기서는 자명한 의미의 미분 등급
쌍선형 사상을 사용한다.





\section{Hom 복합체와 미분 등급 가군}
\label{dga-section-hom-complexes}

\noindent
독자에게 이 절을 건너뛸 것을 권한다.

\medskip\noindent
$R$을 환이라 하고 $M^\bullet$을 $R$-가군의 복합체라 하자.
More on Algebra, Section \ref{more-algebra-section-hom-complexes}에서
도입한 $R$-가군의 복합체
$$
E^\bullet = \Hom^\bullet(M^\bullet, M^\bullet)
$$
를 생각하자. More on Algebra, Lemma
\ref{more-algebra-lemma-composition}에 의해 복합체 사상인 표준 합성 법칙
$$
\text{Tot}(E^\bullet \otimes_R E^\bullet) \to E^\bullet
$$
이 있다. 따라서 이 곱셈을 갖춘 $E^\bullet$은 미분 등급
$R$-대수이며, 이를 $(E, \text{d})$로 나타내기로 한다. 또한
$M^\bullet$을 $\Hom^\bullet(R, M^\bullet)$로 보면 합성이 곱셈
$$
\text{Tot}(E^\bullet \otimes_R M^\bullet) \to M^\bullet
$$
을 정의함을 알 수 있다. 이 곱셈은 $M^\bullet$을 {\bf 왼쪽}
미분 등급 $E$-가군으로 만들며, 이를 $M$으로 나타내기로 한다.

\begin{lemma}
\label{dga-lemma-left-module-structure}
위의 상황에서 $A$를 미분 등급 $R$-대수라 하자. $M$ 위에 왼쪽
$A$-가군 구조를 주는 것은 미분 등급 $R$-대수의 준동형
$A \to E$를 주는 것과 같다.
\end{lemma}

\begin{proof}
증명은 생략한다. 이 대응에는 아무 부호도 개입하지 않음에 유의하라.
\end{proof}

\noindent
위의 논의를 계속하고 $R$-가군의 또 다른 복합체 $N^\bullet$이
주어졌다고 하자. More on Algebra, Section
\ref{more-algebra-section-hom-complexes}에서 도입한 $R$-가군의 복합체
$\Hom^\bullet(M^\bullet, N^\bullet)$를 생각하자. 위와 마찬가지로 합성
$$
\text{Tot}(\Hom^\bullet(M^\bullet, N^\bullet) \otimes_R E^\bullet)
\to \Hom^\bullet(M^\bullet, N^\bullet)
$$
은 $\Hom^\bullet(M^\bullet, N^\bullet)$을 {\bf 오른쪽} 미분 등급
$E$-가군으로 만드는 곱셈을 정의한다. 보조정리
\ref{dga-lemma-left-module-structure}를 사용하면, 왼쪽 미분 등급
$A$-가군 $M$과 $R$-가군의 복합체 $N^\bullet$이 주어졌을 때
바탕 $R$-가군 복합체가 $\Hom^\bullet(M^\bullet, N^\bullet)$인
표준 오른쪽 미분 등급 $A$-가군이 있음을 얻는다. 그 곱셈
$$
\Hom^n(M^\bullet, N^\bullet) \times A^m \longrightarrow
\Hom^{n + m}(M^\bullet, N^\bullet)
$$
은 $f_{p, q} \in \Hom(M^{-q}, N^p)$인
$f = (f_{p, q})_{p + q = n}$과 $a \in A^m$을 원소
$f \cdot a = (f_{p, q} \circ a)$로 보낸다. 여기서
$f_{p, q} \circ a$는 사상
$$
M^{-q - m} \xrightarrow{a} M^{-q} \xrightarrow{f_{p, q}} N^p, \quad
x \longmapsto f_{p, q}(ax)
$$
이며 부호는 개입하지 않는다. 이 오른쪽 미분 등급 $A$-가군을
$\Hom(M, N^\bullet)$로 나타내자.

\begin{lemma}
\label{dga-lemma-characterize-hom}
$R$을 환이라 하자. $(A, \text{d})$를 미분 등급 $R$-대수라 하자.
$M'$을 오른쪽 미분 등급 $A$-가군, $M$을 왼쪽 미분 등급
$A$-가군이라 하고, $N^\bullet$을 $R$-가군의 복합체라 하자. 그러면
$$
\Hom_{\text{Mod}_{(A, d)}}(M', \Hom(M, N^\bullet)) =
\Hom_{\text{Comp}(R)}(M' \otimes_A M, N^\bullet)
$$
이다. 여기서 $M \otimes_A M$은 절 \ref{dga-section-tensor-product}에서와
같이 $R$-가군의 복합체로 본다.
% P03-STACKS-E535: the frozen missing prime is preserved; see the sidecar ledger.
\end{lemma}

\begin{proof}
양변이 미분과 양립하는 등급 $A$-쌍선형 사상
$$
M' \times M \longrightarrow N^\bullet
$$
에 대응함을 보이자. 우변에 대해서는 절 \ref{dga-section-tensor-product}에서
이미 이를 보았다. 좌변의 원소 $g$가 주어지면 More on Algebra, Lemma
\ref{more-algebra-lemma-compose}의 등식이 $R$-가군 복합체의 사상
$g' : \text{Tot}(M' \otimes_R M) \to N^\bullet$을 결정한다.
달리 말해, 미분과 양립하는 등급 $R$-쌍선형 사상
$g'' : M' \times M \to N^\bullet$을 얻는다. $g$의 $A$-선형성은 곧
$g''$의 $A$-쌍선형성으로 옮겨진다.
\end{proof}

\noindent
$R$, $M^\bullet$, $E^\bullet$, $E$, $M$을 위와 같이 두자. 다만 이제
미분 등급 $R$-대수 $A$와 $M$ 위의 {\bf 오른쪽} 미분 등급
$A$-가군 구조가 주어졌다고 하자. 그러면 사상
$$
\text{Tot}(A^\bullet \otimes_R M^\bullet)
\xrightarrow{\psi}
\text{Tot}(A^\bullet \otimes_R M^\bullet)
\to
M^\bullet
$$
을 생각할 수 있다. 여기서 첫째 화살표는 $R$-가군 복합체의 미분 등급
범주에 주어진 가환성 제약이다. 이는 $R$-가군 복합체의 사상
$$
\tau : A^\bullet \longrightarrow E^\bullet
$$
에 대응한다. $E^n = \prod_{p + q = n} \Hom_R(M^{-q}, M^p)$임을
상기하고, $a \in A^n$에 대해
$\tau(a) = (\tau_{p, q}(a))_{p + q = n}$으로 쓰자. 그러면
$$
\tau_{p, q}(a) : M^{-q} \longrightarrow M^p,\quad
x \longmapsto (-1)^{\deg(a)\deg(x)}x a = (-1)^{-nq}xa
$$
임을 알 수 있다. 절 \ref{dga-section-left-modules}의 논의에서 독자가
예상하듯이, 이것은 $A$의 곱과 양립하지 않는다. 즉,
$$
\tau(a a') = (-1)^{\deg(a)\deg(a')}\tau(a') \tau(a)
$$
이다. 따라서 다음 보조정리를 얻는다.

\begin{lemma}
\label{dga-lemma-right-module-structure}
위의 상황에서 $A$를 미분 등급 $R$-대수라 하자. $M$ 위에 오른쪽
$A$-가군 구조를 주는 것은 미분 등급 $R$-대수의 준동형
$\tau : A \to E^{opp}$를 주는 것과 같다.
\end{lemma}

\begin{proof}
위의 논의를 보라. 곱셈 사상 $M^n \times A^m \to M^{n + m}$으로부터
$\tau$를 구성할 때에는 부호를 사용함에 유의하라.
\end{proof}

\noindent
$R$, $M^\bullet$, $E^\bullet$, $E$, $A$, $M$을 위와 같이 두고,
보조정리에서와 같이 $M$ 위에 오른쪽 미분 등급 $A$-가군 구조가
주어졌다고 하자. 이 경우 바탕 $R$-가군 복합체가
$\Hom^\bullet(M^\bullet, N^\bullet)$인 표준 왼쪽 미분 등급
$A$-가군이 있다. 구체적으로 곱셈에는 다음을 사용할 수 있다.
% P03-STACKS-E536: the frozen unclosed final expression is preserved.
\begin{align*}
\text{Tot}(A^\bullet \otimes_R \Hom^\bullet(M^\bullet, N^\bullet))
& \xrightarrow{\psi}
\text{Tot}(\Hom^\bullet(M^\bullet, N^\bullet) \otimes_R A^\bullet) \\
& \xrightarrow{\tau}
\text{Tot}(\Hom^\bullet(M^\bullet, N^\bullet) \otimes_R
\Hom^\bullet(M^\bullet, M^\bullet)) \\
& \to
\text{Tot}(\Hom^\bullet(M^\bullet, N^\bullet)
\end{align*}
첫째 화살표는 $R$-가군 복합체 범주의 가환성 제약을 사용하고,
둘째 화살표는 위에서 설명했으며, 셋째 화살표는 Hom 복합체의 합성
법칙이다. 각 사상이 복합체의 사상이므로 결과도 복합체의 사상이다.
실제로 이 구성은 $\Hom^\bullet(M^\bullet, N^\bullet)$을 왼쪽 미분
등급 $A$-가군으로 만든다. 곱셈의 결합성은 대칭 모노이드 구조를
사용하거나 아래 공식으로 직접 계산하여 보일 수 있다. 곱셈을 명시하면
$$
A^n \times \Hom^m(M^\bullet, N^\bullet) \longrightarrow
\Hom^{n + m}(M^\bullet, N^\bullet)
$$
이다. 이는 $a \in A^n$과 $f_{p, q} \in \Hom(M^{-q}, N^p)$인
$f = (f_{p, q})_{p + q = m}$을 다음 성분들을 갖는 원소
$a \cdot f$로 보낸다.
$$
(-1)^{nm}f_{p, q} \circ \tau_{-q, q + n}(a) =
(-1)^{nm - n(q + n)}f_{p, q} \circ a =
(-1)^{np + n} f_{p, q} \circ a
$$
이 성분은 $\Hom_R(M^{-q - n}, N^p)$에 속하며,
$f_{p, q} \circ a$는 사상
$$
M^{-q - n} \xrightarrow{a} M^{-q} \xrightarrow{f_{p, q}} N^p,\quad
x \longmapsto f_{p, q}(xa)
$$
이다. 여기서는 $(-1)^{np + n}$이라는 부호가 실제로 개입한다.
이 왼쪽 미분 등급 $A$-가군을 $\Hom(M, N^\bullet)$로 나타내자.

\begin{lemma}
\label{dga-lemma-characterize-hom-other-side}
$R$을 환이라 하자. $(A, \text{d})$를 미분 등급 $R$-대수라 하자.
$M$을 오른쪽 미분 등급 $A$-가군, $M'$을 왼쪽 미분 등급
$A$-가군이라 하고, $N^\bullet$을 $R$-가군의 복합체라 하자. 그러면
$$
\Hom_{\text{왼쪽 미분 등급 }A\text{-가군}}(M', \Hom(M, N^\bullet)) =
\Hom_{\text{Comp}(R)}(M \otimes_A M', N^\bullet)
$$
이다. 여기서 $M \otimes_A M'$은 절 \ref{dga-section-tensor-product}에서와
같이 $R$-가군의 복합체로 본다.
\end{lemma}

\begin{proof}
양변이 미분과 양립하는 등급 $A$-쌍선형 사상
$$
M \times M' \longrightarrow N^\bullet
$$
에 대응함을 보이자. 우변에 대해서는 절 \ref{dga-section-tensor-product}에서
이미 이를 보았다. 좌변의 원소 $g$가 주어지면 More on Algebra, Lemma
\ref{more-algebra-lemma-compose}의 등식이 복합체의 사상
$g' : \text{Tot}(M' \otimes_R M) \to N^\bullet$을 결정한다.
가환성 제약을 앞에 합성하여
$$
\text{Tot}(M \otimes_R M') \xrightarrow{\psi}
\text{Tot}(M' \otimes_R M) \xrightarrow{g'}
N^\bullet
$$
을 얻는다. 이는 미분과 양립하는 등급 $R$-쌍선형 사상
$g'' : M \times M' \to N^\bullet$에 대응한다. $g$의 $A$-선형성은
곧 $g''$의 $A$-쌍선형성으로 옮겨진다. 구체적으로
$x \in M^e$, $x' \in (M')^{e'}$, $a \in A^n$이라 하자. 한편으로
\begin{align*}
g''(x, ax')
& =
(-1)^{e(n + e')} g'(ax' \otimes x) \\
& =
(-1)^{e(n + e')} g(ax')(x) \\
& =
(-1)^{e(n + e')} (a \cdot g(x'))(x) \\
& =
(-1)^{e(n + e') + n(n + e + e') + n} g(x')(xa)
\end{align*}
이고, 다른 한편으로
$$
g''(xa, x') = (-1)^{(e + n)e'} g'(x' \otimes xa) =
(-1)^{(e + n)e'} g(x')(xa) 
$$
이다. 지수들에 대한 자명한 모듈로 $2$ 계산에 의해 이 둘은 같다.
\end{proof}

\begin{remark}
\label{dga-remark-evaluation-map-left}
$R$을 환이라 하자. $A$를 미분 등급 $R$-대수라 하자.
$M$을 왼쪽 미분 등급 $A$-가군이라 하자. $N^\bullet$을
$R$-가군의 복합체라 하자. 위의 구성들은 오른쪽 미분 등급
$A$-가군 $\Hom(M, N^\bullet)$을 만들고, 이어서 왼쪽 미분 등급
$A$-가군 $\Hom(\Hom(M, N^\bullet), N^\bullet)$을 만든다.
다음 평가 사상
$$
ev : M \longrightarrow \Hom(\Hom(M, N^\bullet), N^\bullet)
$$
이 왼쪽 미분 등급 $A$-가군의 범주에 존재한다고 주장한다.
이를 정의하기 위해 보조정리 \ref{dga-lemma-characterize-hom}에 따라
등급 및 미분과 양립하는 $A$-쌍선형 짝짓기
$$
\Hom(M, N^\bullet) \times M \longrightarrow N^\bullet
$$
를 구성하면 충분하다. 이를 위해
$$
(f, x) \longmapsto f(x)
$$
로 둔다. 이것이 등급 및 미분과 양립하고 $A$-쌍선형이라는 검증은
독자에게 맡긴다. 바탕 $R$-가군 복합체 위에서 사상 $ev$는
More on Algebra, Item (\ref{more-algebra-item-evaluation})이다.
\end{remark}

\begin{remark}
\label{dga-remark-evaluation-map-right}
$R$을 환이라 하자. $A$를 미분 등급 $R$-대수라 하자.
$M$을 오른쪽 미분 등급 $A$-가군이라 하자. $N^\bullet$을
$R$-가군의 복합체라 하자. 위의 구성들은 왼쪽 미분 등급
$A$-가군 $\Hom(M, N^\bullet)$을 만들고, 이어서 오른쪽 미분 등급
$A$-가군 $\Hom(\Hom(M, N^\bullet), N^\bullet)$을 만든다.
다음 평가 사상
$$
ev : M \longrightarrow \Hom(\Hom(M, N^\bullet), N^\bullet)
$$
이 오른쪽 미분 등급 $A$-가군의 범주에 존재한다고 주장한다.
이를 정의하기 위해 보조정리 \ref{dga-lemma-characterize-hom}에 따라
등급 및 미분과 양립하는 $A$-쌍선형 짝짓기
$$
M \times \Hom(M, N^\bullet) \longrightarrow N^\bullet
$$
를 구성하면 충분하다. 이를 위해
$$
(x, f) \longmapsto (-1)^{\deg(x)\deg(f)}f(x)
$$
로 둔다. 이것이 등급 및 미분과 양립하고 $A$-쌍선형이라는 검증은
독자에게 맡긴다. 바탕 $R$-가군 복합체 위에서 사상 $ev$는
More on Algebra, Item (\ref{more-algebra-item-evaluation})이다.
\end{remark}

\begin{remark}
\label{dga-remark-shift-dual}
$R$을 환이라 하자. $A$를 미분 등급 $R$-대수라 하자.
$M^\bullet$과 $N^\bullet$을 $R$-가군의 복합체라 하자.
$k \in \mathbf{Z}$라 하고, More on Algebra, Item
(\ref{more-algebra-item-shift-hom})에서 정의한 $R$-가군 복합체의 동형사상
$$
\Hom^\bullet(M^\bullet, N^\bullet)[-k]
\longrightarrow
\Hom^\bullet(M^\bullet[k], N^\bullet)
$$
을 생각하자. $M^\bullet$이 왼쪽, 각각 오른쪽 미분 등급 $A$-가군
구조를 가지면, 이는 오른쪽, 각각 왼쪽 미분 등급 $A$-가군의 사상이다.
여기서 가군 구조는 이 절에서 정의한 것을 사용한다. 검증은 생략한다.
왼쪽 등급 $A$-가군의 시프트 위의 $A$-가군 구조는 부호를 사용하여
정의된다는 점을 독자에게 경고한다. 정의
\ref{dga-definition-shift-graded-module}를 보라.
\end{remark}





\section{대수 위의 사영 가군}
\label{dga-section-projectives-over-algebras}

\noindent
이 절에서는 Algebra, Section \ref{algebra-section-projective}와 유사하게
대수 위의 사영 가군을 논의한다. 이 절은 아마 다른 곳으로 옮겨야 할 것이다.

\medskip\noindent
$R$을 환이라 하고 $A$를 $R$-대수라 하자. 우리의 규약은 절
\ref{dga-section-conventions}를 보라. 임의의 오른쪽 $A$-가군 $M$에 대해
$\Hom_A(A, M) = M$이므로 $A$는 사영 오른쪽 $A$-가군임이 분명하다
(따라서 $\Hom_A(A, -)$는 완전하다). 역으로, $P$를 사영 오른쪽
$A$-가군이라 하자. $P$의 $A$ 위의 생성원 집합
$\{p_i\}_{i \in I}$를 택하여 전사
$\bigoplus_{i \in I} A \to P$를 택할 수 있다. $P$가 사영이므로
이 전사에는 왼쪽 역이 있다. 따라서 가환인 경우와 정확히 마찬가지로
$P$는 자유 가군의 직합인자와 동형이다(Algebra, Lemma
\ref{algebra-lemma-characterize-projective}).

\medskip\noindent
다음을 얻는다.
\begin{enumerate}
\item $A$-가군의 범주는 충분한 사영 대상을 가진다.
\item $A$는 사영 $A$-가군이다.
\item 모든 $A$-가군은 $A$의 복사본들의 직합의 몫이다.
\item 모든 사영 $A$-가군은 $A$의 복사본들의 직합의 직합인자이다.
\end{enumerate}



\section{등급 대수 위의 사영 가군}
\label{dga-section-projectives-over-graded-algebras}

\noindent
이 절에서는 Algebra, Section \ref{algebra-section-projective}와 유사하게
등급 대수 위의 사영 등급 가군을 논의한다.

\medskip\noindent
$R$을 환이라 하자. $A$를 $R$ 위의 $\mathbf{Z}$-등급 대수라 하자.
우리의 규약은 절 \ref{dga-section-conventions}를 보라. 등급 오른쪽
$A$-가군의 범주를 $\text{Mod}_A$로 나타내자. 정수 $k$에 대해
$A$의 시프트를 $A[k]$로 나타내자. 등급 오른쪽 $A$-가군에 대해
$$
\Hom_{\text{Mod}_A}(A[k], M) = M^{-k}
$$
이다. 함자 $M \mapsto M^{-k}$가 $\text{Mod}_A$ 위에서 완전하므로
$A[k]$는 $\text{Mod}_A$의 사영 대상이다. 역으로, $P$가
$\text{Mod}_A$의 사영 대상이라고 하자. $P$를 $A$-가군으로 보아
동차 생성원들의 집합을 택하면 전사
$$
\bigoplus\nolimits_{i \in I} A[k_i] \longrightarrow P
$$
를 얻을 수 있다. 따라서 $\text{Mod}_A$의 사영 대상은 시프트
$A[k]$들의 직합의 직합인자이다.

\medskip\noindent
다음을 얻는다.
\begin{enumerate}
\item 등급 $A$-가군의 범주는 충분한 사영 대상을 가진다.
\item 모든 $k \in \mathbf{Z}$에 대해 $A[k]$는 사영 $A$-가군이다.
\item 모든 등급 $A$-가군은 $k$가 변할 때의 가군 $A[k]$들의
복사본들의 직합의 몫이다.
\item 모든 사영 $A$-가군은 $k$가 변할 때의 가군 $A[k]$들의
복사본들의 직합의 직합인자이다.
\end{enumerate}




\section{사영 가군과 미분 등급 대수}
\label{dga-section-projective-over-differential-graded}

\noindent
$(A, \text{d})$가 미분 등급 대수이고 $P$가
$\text{Mod}_{(A, \text{d})}$의 대상이면, $P$가 등급 $A$-가군들의
아벨 범주 $\text{Mod}_A$에서 사영 대상이라는 뜻으로
{\it $P$는 등급 $A$-가군으로서 사영이다}라고 하거나, 때로
{\it $P$는 등급 사영이다}라고 한다. 여기서 이 아벨 범주는 절
\ref{dga-section-projectives-over-graded-algebras}에서와 같다.

\begin{lemma}
\label{dga-lemma-target-graded-projective}
$(A, \text{d})$를 미분 등급 대수라 하자. $M \to P$를 미분 등급
$A$-가군의 전사 준동형이라 하자. $P$가 등급 $A$-가군으로서
사영이면 $M \to P$는 허용 전사상이다.
\end{lemma}

\begin{proof}
이는 정의에서 바로 따른다.
\end{proof}

\begin{lemma}
\label{dga-lemma-hom-from-shift-free}
$(A, d)$를 미분 등급 대수라 하자. 그러면 임의의 미분 등급
$A$-가군 $M$에 대해
$$
\Hom_{\text{Mod}_{(A, \text{d})}}(A[k], M) =
\Ker(\text{d} : M^{-k} \to M^{-k + 1})
$$
이고
$$
\Hom_{K(\text{Mod}_{(A, \text{d})})}(A[k], M) = H^{-k}(M)
$$
이다.
\end{lemma}

\begin{proof}
정의에서 바로 따른다.
\end{proof}







\section{대수 위의 단사 가군}
\label{dga-section-modules-noncommutative}

\noindent
이 절에서는 More on Algebra, Section
\ref{more-algebra-section-injectives-modules}와 유사하게 대수 위의
단사 가군을 논의한다. 이 절은 아마 다른 곳으로 옮겨야 할 것이다.

\medskip\noindent
$R$을 환이라 하고 $A$를 $R$-대수라 하자. 우리의 규약은 절
\ref{dga-section-conventions}를 보라. 오른쪽 $A$-가군 $M$에 대해
$$
M^\vee = \Hom_\mathbf{Z}(M, \mathbf{Q}/\mathbf{Z})
$$
로 둔다. 곱셈 $(a f)(x) = f(xa)$에 의해 이를 왼쪽 $A$-가군으로
생각한다. 즉,
$((ab)f)(x) = f(xab) = (bf)(xa) = (a(bf))(x)$이다. 역으로,
$M$이 왼쪽 $A$-가군이면 $M^\vee$는 오른쪽 $A$-가군이다.
$\mathbf{Q}/\mathbf{Z}$가 단사 아벨 군이므로(More on Algebra, Lemma
\ref{more-algebra-lemma-injective-abelian}) 함자
$M \mapsto M^\vee$는 완전하다(More on Algebra, Lemma
\ref{more-algebra-lemma-vee-exact}). 또한 모든 가군 $M$에 대해
평가 사상 $M \to (M^\vee)^\vee$는 단사이다(More on Algebra, Lemma
\ref{more-algebra-lemma-ev-injective}).

\medskip\noindent
$A^\vee$가 단사 오른쪽 $A$-가군이라고 주장한다. 실제로 오른쪽
$A$-가군 $N$이 주어지면
$$
\Hom_A(N, A^\vee) =
\Hom_A(N, \Hom_\mathbf{Z}(A, \mathbf{Q}/\mathbf{Z})) = N^\vee
$$
이고, 함자 $N \mapsto N^\vee$가 완전하므로 결론을 얻는다.
둘째 등식은 Algebra, Lemma
\ref{algebra-lemma-hom-from-tensor-product}에 의해
$$
\Hom_\mathbf{Z}(N, \Hom_\mathbf{Z}(A, \mathbf{Q}/\mathbf{Z})) =
\Hom_\mathbf{Z}(N \otimes_\mathbf{Z} A, \mathbf{Q}/\mathbf{Z})
$$
이므로 성립한다. 이 가군 안에서 $A$-선형성은
$\varphi : N \times A \to \mathbf{Q}/\mathbf{Z}$ 중
$x \otimes a$와 $xa \otimes 1$에서 같은 값을 갖는 쌍선형 사상들,
즉 $N^\vee$의 원소에서 오는 것들을 정확히 골라낸다.

\medskip\noindent
마지막으로 모든 오른쪽 $A$-가군 $M$에 대해 전사
$\bigoplus_{i \in I} A \to M^\vee$를 택하여 단사
$M \to (M^\vee)^\vee \to \prod_{i \in I} A^\vee$를 얻을 수 있다.

\medskip\noindent
다음을 얻는다.
\begin{enumerate}
\item $A$-가군의 범주는 충분한 단사 대상을 가진다.
\item $A^\vee$는 단사 $A$-가군이다.
\item 모든 $A$-가군은 $A^\vee$의 복사본들의 곱 안으로 단사된다.
\end{enumerate}





\section{등급 대수 위의 단사 가군}
\label{dga-section-modules-noncommutative-graded}

\noindent
이 절에서는 More on Algebra, Section
\ref{more-algebra-section-injectives-modules}와 유사하게 등급 대수 위의
단사 등급 가군을 논의한다.

\medskip\noindent
$R$을 환이라 하자. $A$를 $R$ 위의 $\mathbf{Z}$-등급 대수라 하자.
우리의 규약은 절 \ref{dga-section-conventions}를 보라. $M$이 등급
$R$-가군이면 등급 $R$-가군으로서
$$
M^\vee =
\bigoplus\nolimits_{n \in \mathbf{Z}}
\Hom_\mathbf{Z}(M^{-n}, \mathbf{Q}/\mathbf{Z}) =
\bigoplus\nolimits_{n \in \mathbf{Z}} (M^{-n})^\vee
$$
로 둔다($R$이 동차 성분들에 작용할 때에는 부호를 쓰지 않는다).
$M$이 왼쪽 등급 $A$-가군 구조를 가지면 $a \in A^m$의 작용을
$$
(M^{-n})^\vee \to (M^{-n - m})^\vee, \quad
f \mapsto f \circ a
$$
로 두어 $M^\vee$ 위에 오른쪽 등급 $A$-가군 구조를 정의한다.
이는 절 \ref{dga-section-hom-complexes}에서와 같다. $M$이 오른쪽 등급
$A$-가군 구조를 가지면 $a \in A^n$의 작용을
$$
(M^{-m})^\vee \to (M^{-m - n})^\vee, \quad
f \mapsto (-1)^{nm}f \circ a
$$
로 두어 $M^\vee$ 위에 왼쪽 등급 $A$-가군 구조를 정의한다.
이 역시 절 \ref{dga-section-hom-complexes}에서와 같다. 부호를 쓰는 이유는
미분 등급 가군에서도 같은 공식을 사용하려 하기 때문이다. 등급 가군에만
관심이 있다면 여기서 부호를 생략할 수 있다. (왼쪽 또는 오른쪽) 등급
$A$-가군의 범주에서 함자 $M \mapsto M^\vee$는 완전하다
(각 등급 성분에서 확인하라). 또한 등급 $R$-가군의 단사 평가 사상
$$
ev : M \longrightarrow (M^\vee)^\vee, \quad
ev^n = (-1)^n \text{ 배의 평가 사상 }M^n \to ((M^n)^\vee)^\vee
$$
이 있다. More on Algebra, Item
(\ref{more-algebra-item-evaluation})을 보라. $M$이 왼쪽, 각각 오른쪽
$A$-가군이면 이 평가 사상은 왼쪽, 각각 오른쪽 $A$-가군 준동형이다.
평주 \ref{dga-remark-evaluation-map-left}와
\ref{dga-remark-evaluation-map-right}를 보라. 마지막으로
$k \in \mathbf{Z}$가 주어지면 부호를 사용하는 등급 $R$-가군의
표준 동형사상
$$
M^\vee[-k] \longrightarrow (M[k])^\vee
$$
이 있다. $M$이 왼쪽, 각각 오른쪽 $A$-가군이면 이는 오른쪽, 각각
왼쪽 $A$-가군의 동형사상이다. 평주 \ref{dga-remark-shift-dual}를 보라.

\medskip\noindent
% P03-STACKS-E886: the frozen extra closing parenthesis is preserved.
$A^\vee$가 등급 오른쪽 $A$-가군의 범주 $\text{Mod}_A$의 단사 대상이라고
주장한다. 실제로 등급 오른쪽 $A$-가군 $N$이 주어지면
$$
\Hom_{\text{Mod}_A}(N, A^\vee) =
\Hom_{\text{Comp}(\mathbf{Z})}(N \otimes_A A, \mathbf{Q}/\mathbf{Z})) =
(N^0)^\vee
$$
이다. 이는 보조정리 \ref{dga-lemma-characterize-hom}을 모든 미분이
영인 경우에 적용한 것이다. 함자
$N \mapsto (N^0)^\vee = (N^\vee)^0$가 완전하므로 결론을 얻는다.

\medskip\noindent
마지막으로 모든 등급 오른쪽 $A$-가군 $M$에 대해 등급 왼쪽
$A$-가군의 전사
$$
\bigoplus\nolimits_{i \in I} A[k_i] \to M^\vee
$$
를 택할 수 있다. 여기서 $A[k_i]$는 $k_i \in \mathbf{Z}$만큼의
$A$의 시프트이다. 이는 $M^\vee$의 동차 생성원들을 택하여 얻는다.
이 방법으로 단사
$$
M \to (M^\vee)^\vee \to \prod A[k_i]^\vee = \prod A^\vee[-k_i]
$$
를 얻는다. 위 공식의 곱들은 등급 가군의 범주에서의 곱들이다.
다시 말해, 각 차수에서 곱을 취한 뒤 성분들의 직합을 취한다.

\medskip\noindent
다음을 얻는다.
\begin{enumerate}
\item 등급 $A$-가군의 범주는 충분한 단사 대상을 가진다.
\item 모든 $k \in \mathbf{Z}$에 대해 가군 $A^\vee[k]$는 단사이다.
\item 모든 $A$-가군은 등급 가군의 범주에서 $A^\vee[k]$의
시프트 복사본들의 곱 안으로 단사된다.
\end{enumerate}





\section{단사 가군과 미분 등급 대수}
\label{dga-section-modules-noncommutative-differential-graded}

\noindent
$(A, \text{d})$가 미분 등급 대수이고 $I$가
$\text{Mod}_{(A, \text{d})}$의 대상이면, $I$가 등급 $A$-가군들의
아벨 범주 $\text{Mod}_A$에서 단사 대상이라는 뜻으로
{\it $I$는 등급 $A$-가군으로서 단사이다}라고 하거나, 때로
{\it $I$는 등급 단사이다}라고 한다.

\begin{lemma}
\label{dga-lemma-source-graded-injective}
$(A, \text{d})$를 미분 등급 대수라 하자. $I \to M$을 미분 등급
$A$-가군의 단사 준동형이라 하자. $I$가 등급 단사이면
$I \to M$은 허용 단사상이다.
\end{lemma}

\begin{proof}
이는 정의에서 바로 따른다.
\end{proof}

\noindent
$(A, \text{d})$를 미분 등급 대수라 하자. $M$이 왼쪽, 각각 오른쪽
미분 등급 $A$-가군이면
$$
M^\vee = \Hom^\bullet(M^\bullet, \mathbf{Q}/\mathbf{Z})
$$
와 절 \ref{dga-section-modules-noncommutative-graded}에서 구성한
$A$-가군 구조는 절 \ref{dga-section-hom-complexes}의 논의에 의해
오른쪽, 각각 왼쪽 미분 등급 $A$-가군을 이룬다. 평주
\ref{dga-remark-evaluation-map-left}와 \ref{dga-remark-evaluation-map-right}에 의해
절 \ref{dga-section-modules-noncommutative-graded}의 평가 사상
$$
M \longrightarrow (M^\vee)^\vee
$$
은 왼쪽, 각각 오른쪽 미분 등급 $A$-가군의 준동형이다.

\begin{lemma}
\label{dga-lemma-map-into-dual}
$(A, \text{d})$를 미분 등급 대수라 하자. $M$이 왼쪽 미분 등급
$A$-가군이고 $N$이 오른쪽 미분 등급 $A$-가군이면
\begin{align*}
\Hom_{\text{Mod}_{(A, \text{d})}}(N, M^\vee)
& =
\Hom_{\text{Comp}(\mathbf{Z})}(N \otimes_A M, \mathbf{Q}/\mathbf{Z}) \\
& =
\text{미분 등급 쌍선형}_A(N \times M, \mathbf{Q}/\mathbf{Z})
\end{align*}
이다.
\end{lemma}

\begin{proof}
첫째 등식은 보조정리 \ref{dga-lemma-characterize-hom}이고, 둘째 등식은
절 \ref{dga-section-tensor-product}에서 보였다.
\end{proof}

\begin{lemma}
\label{dga-lemma-hom-into-shift-dual-free}
$(A, \text{d})$를 미분 등급 대수라 하자. 그러면
$$
\Hom_{\text{Mod}_{(A, \text{d})}}(M, A^\vee[k]) =
\Ker(\text{d} : (M^\vee)^k \to (M^\vee)^{k + 1})
$$
이고
$$
\Hom_{K(\text{Mod}_{(A, \text{d})})}(M, A^\vee[k]) = H^k(M^\vee)
$$
이다. 이 등식들은 미분 등급 $A$-가군 $M$에 대한 함자 등식이다.
\end{lemma}

\begin{proof}
이는 위의 논의에서 분명하다.
\end{proof}















\section{P-분해}
\label{dga-section-P-resolutions}

\noindent
이 절은 Derived Categories, Section \ref{derived-section-unbounded}의
대응물이다.

\medskip\noindent
$(A, \text{d})$를 미분 등급 대수라 하자. $P$를 미분 등급
$A$-가군이라 하자. 다음 조건을 만족하는 미분 등급 부분가군들에 의한 여과
$$
0 = F_{-1}P \subset F_0P \subset F_1P \subset \ldots \subset P
$$
가 존재하면 $P$가 {\it 성질 (P)를 갖는다}고 한다.
\begin{enumerate}
\item $P = \bigcup F_pP$이다.
\item 포함 $F_iP \to F_{i + 1}P$는 허용 단사상이다.
\item 몫 $F_{i + 1}P/F_iP$는 미분 등급 $A$-가군으로서
$A[k]$들의 직합과 동형이다.
\end{enumerate}
실제로 조건 (2)는 조건 (3)의 귀결이다. 보조정리
\ref{dga-lemma-target-graded-projective}를 보라. 또한 독자는 등급
$A$-가군으로서 $P$가 $A$의 시프트들의 직합과 동형임을 확인할 수 있다.

\begin{lemma}
\label{dga-lemma-property-P-sequence}
$(A, \text{d})$를 미분 등급 대수라 하자. $P$를 미분 등급
$A$-가군이라 하자. $F_\bullet$이 성질 (P)에서와 같은 여과이면
미분 등급 $A$-가군의 허용 짧은 완전열
$$
0 \to
\bigoplus\nolimits F_iP \to
\bigoplus\nolimits F_iP \to P \to 0
$$
을 얻는다.
\end{lemma}

\begin{proof}
둘째 사상은 포함 사상들의 직합이다. 첫째 사상의 정의역에서 직합인자
$F_iP$ 위의 사상은 항등사상 $F_iP \to F_iP$와 포함 사상
$F_iP \to F_{i + 1}P$의 음의 합이다. 포함 사상들의 왼쪽 역인
등급 $A$-가군의 준동형 $s_i : F_{i + 1}P \to F_iP$들을 택하자.
이들을 합성하여 모든 $j > i$에 대한 사상
$s_{j, i} : F_jP \to F_iP$를 얻는다. 그러면 첫째 화살표의 왼쪽 역은
$x \in F_jP$를 $\bigoplus F_iP$ 안의
$(s_{j, 0}(x), s_{j, 1}(x), \ldots, s_{j, j - 1}(x), 0, \ldots)$로 보낸다.
\end{proof}

\noindent
다음 보조정리는 성질 (P)를 갖는 미분 등급 가군이 K-단사 가군의
쌍대 개념임을 보인다(즉, 어떤 의미에서는 K-사영이다). Derived Categories,
Definition \ref{derived-definition-K-injective}을 보라.

\begin{lemma}
\label{dga-lemma-property-P-K-projective}
$(A, \text{d})$를 미분 등급 대수라 하자. $P$를 성질 (P)를 갖는
미분 등급 $A$-가군이라 하자. 그러면 모든 비순환 미분 등급
$A$-가군 $N$에 대해
$$
\Hom_{K(\text{Mod}_{(A, \text{d})})}(P, N) = 0
$$
이다.
\end{lemma}

\begin{proof}
$K(\text{Mod}_{(A, \text{d})})$가 삼각범주라는 사실을 사용할 것이다
(명제 \ref{dga-proposition-homotopy-category-triangulated}). $F_\bullet$을
성질 (P)에서와 같은 $P$ 위의 여과라 하자. 보조정리
\ref{dga-lemma-property-P-sequence}의 짧은 완전열은 특수 삼각형을 만든다.
따라서 Derived Categories, Lemma
\ref{derived-lemma-representable-homological}에 의해 모든 비순환
미분 등급 $A$-가군 $N$과 모든 $i$에 대해
$$
\Hom_{K(\text{Mod}_{(A, \text{d})})}(F_iP, N) = 0
$$
임을 보이면 충분하다. 각 미분 등급 가군 $F_iP$에는 허용 단사상들에
의한 유한 여과가 있고, 그 등급 조각들은 시프트 $A[k]$들의 직합이다.
따라서 모든 비순환 미분 등급 $A$-가군 $N$과 모든 $k$에 대해
$$
\Hom_{K(\text{Mod}_{(A, \text{d})})}(A[k], N) = 0
$$
임을 증명하면 충분하다. 이는 보조정리
\ref{dga-lemma-hom-from-shift-free}에서 따른다.
\end{proof}

\begin{lemma}
\label{dga-lemma-good-quotient}
$(A, \text{d})$를 미분 등급 대수라 하자. $M$을 미분 등급
$A$-가군이라 하자. 다음 성질들을 갖는 미분 등급 $A$-가군의 준동형
$P \to M$이 존재한다.
\begin{enumerate}
\item $P \to M$은 전사이다.
\item $\Ker(\text{d}_P) \to \Ker(\text{d}_M)$은 전사이다.
\item $P$는 허용 짧은 완전열 $0 \to P' \to P \to P'' \to 0$에
들어가며, 여기서 $P'$, $P''$는 $A$의 시프트들의 직합이다.
\end{enumerate}
\end{lemma}

\begin{proof}
$P_k$를 차수 $k$와 $k + 1$의 생성원 $x, y$를 갖는 자유
$A$-가군이라 하자. $\text{d}(x) = y$와 $\text{d}(y) = 0$으로 두어
$P_k$ 위에 미분 등급 $A$-가군 구조를 정의한다. 모든 원소
$m \in M^k$에 대해 $x$를 $m$으로, $y$를 $\text{d}(m)$으로 보내는
준동형 $P_k \to M$이 있다. 따라서 $P_k$의 복사본들의 직합에서
$M$으로 가는 전사가 있음을 알 수 있다. 이는 성질 (1)과 (3)을 갖는
$P \to M$을 분명히 만든다. 성질 (2)를 얻기 위해
$m \in \Ker(\text{d}_M)$의 차수가 $k$이면 $1$을 $m$으로 보내는 사상
$A[k] \to M$이 있음에 유의하라. 따라서 $A$의 시프트들의 복사본들의
직합을 더하여 (2)를 얻을 수 있다.
\end{proof}

\begin{lemma}
\label{dga-lemma-resolve}
$(A, \text{d})$를 미분 등급 대수라 하자. $M$을 미분 등급
$A$-가군이라 하자. 다음을 만족하는 미분 등급 $A$-가군의 준동형
$P \to M$이 존재한다.
\begin{enumerate}
\item $P \to M$은 유사동형이다.
\item $P$는 성질 (P)를 갖는다.
\end{enumerate}
\end{lemma}

\begin{proof}
$M = M_0$으로 둔다. 보조정리 \ref{dga-lemma-good-quotient}에서와 같이
사상 $P_i \to M_i$를 택하여 짧은 완전열
$$
0 \to M_{i + 1} \to P_i \to M_i \to 0
$$
들을 귀납적으로 택한다. 이로부터 “분해”
$$
\ldots \to P_2 \xrightarrow{f_2} P_1 \xrightarrow{f_1} P_0 \to M \to 0
$$
를 얻는다. 이제 $A$-가군으로서
$$
P = \bigoplus\nolimits_{i \geq 0} P_i
$$
로 둔다. 등급은 $P^n = \bigoplus_{a + b = n} P_{-a}^b$로 주고,
미분은 이중복합체에 결부된 전체 복합체의 구성에서와 같이
$x \in P_{-a}^b$에 대해
$$
\text{d}_P(x) = f_{-a}(x) + (-1)^a \text{d}_{P_{-a}}(x)
$$
로 준다. 이 규약에 따라 $P$는 실제로 미분 등급 $A$-가군이다.
각 $P_i$가 두 단계 여과
$0 \to P_i' \to P_i \to P_i'' \to 0$를 가짐을 상기하고,
$$
F_{2i}P = \bigoplus\nolimits_{i \geq j \geq 0} P_j
\subset
\bigoplus\nolimits_{i \geq 0} P_i = P
$$
로 둔다. 그리고 $P'_{i + 1}$을 $F_{2i}P$에 더하여
$F_{2i + 1}$을 얻는다. 이들은 미분 등급 부분가군들이고 연속하는
몫들은 $A$의 시프트들의 직합이다. 보조정리
\ref{dga-lemma-target-graded-projective}에 의해 포함
$F_iP \to F_{i + 1}P$가 허용 단사상임을 알 수 있다. 마지막으로
증대 $P_0 \to M$으로 주어진 사상 $P \to M$이 유사동형임을 보여야
한다. 이는 Homology, Lemma
\ref{homology-lemma-good-resolution-gives-qis}에서 따른다.
\end{proof}





\section{I-분해}
\label{dga-section-I-resolutions}

\noindent
이 절은 P-분해에 관한 절의 쌍대이다.

\medskip\noindent
$(A, \text{d})$를 미분 등급 대수라 하자. $I$를 미분 등급
$A$-가군이라 하자. 다음 조건을 만족하는 미분 등급 부분가군들에 의한 여과
$$
I = F_0I \supset F_1I \supset F_2I \supset \ldots \supset 0
$$
가 존재하면 $I$가 {\it 성질 (I)를 갖는다}고 한다.
\begin{enumerate}
\item $I = \lim I/F_pI$이다.
\item 사상 $I/F_{i + 1}I \to I/F_iI$는 허용 전사상이다.
\item 몫 $F_iI/F_{i + 1}I$는 미분 등급 $A$-가군으로서 절
\ref{dga-section-modules-noncommutative-differential-graded}에서 구성한
가군 $A^\vee[k]$들의 곱과 동형이다.
\end{enumerate}
실제로 조건 (2)는 조건 (3)의 귀결이다. 보조정리
\ref{dga-lemma-source-graded-injective}를 보라. 독자는 등급 가군으로서
$I$가 $A^\vee[k]$들의 곱과 동형임을 확인할 수 있다.

\begin{lemma}
\label{dga-lemma-property-I-sequence}
$(A, \text{d})$를 미분 등급 대수라 하자. $I$를 미분 등급
$A$-가군이라 하자. $F_\bullet$이 성질 (I)에서와 같은 여과이면
미분 등급 $A$-가군의 허용 짧은 완전열
$$
0 \to I \to
\prod\nolimits I/F_iI \to
\prod\nolimits I/F_iI \to 0
$$
을 얻는다.
\end{lemma}

\begin{proof}
생략한다. 힌트: 이는 보조정리 \ref{dga-lemma-property-P-sequence}의 쌍대이다.
\end{proof}

\noindent
다음 보조정리는 성질 (I)를 갖는 미분 등급 가군이 K-단사 가군의
대응물임을 보인다. Derived Categories, Definition
\ref{derived-definition-K-injective}을 보라.

\begin{lemma}
\label{dga-lemma-property-I-K-injective}
$(A, \text{d})$를 미분 등급 대수라 하자. $I$를 성질 (I)를 갖는
미분 등급 $A$-가군이라 하자. 그러면 모든 비순환 미분 등급
$A$-가군 $N$에 대해
$$
\Hom_{K(\text{Mod}_{(A, \text{d})})}(N, I) = 0
$$
이다.
\end{lemma}

\begin{proof}
$K(\text{Mod}_{(A, \text{d})})$가 삼각범주라는 사실을 사용할 것이다
(명제 \ref{dga-proposition-homotopy-category-triangulated}). $F_\bullet$을
성질 (I)에서와 같은 $I$ 위의 여과라 하자. 보조정리
\ref{dga-lemma-property-I-sequence}의 짧은 완전열은 특수 삼각형을 만든다.
따라서 Derived Categories, Lemma
\ref{derived-lemma-representable-homological}에 의해 모든 비순환
미분 등급 $A$-가군 $N$과 모든 $i$에 대해
$$
\Hom_{K(\text{Mod}_{(A, \text{d})})}(N, I/F_iI) = 0
$$
임을 보이면 충분하다. 각 미분 등급 가군 $I/F_iI$에는 허용
단사상들에 의한 유한 여과가 있고, 그 등급 조각들은
$A^\vee[k]$들의 곱이다. 따라서 모든 비순환 미분 등급
$A$-가군 $N$과 모든 $k$에 대해
$$
\Hom_{K(\text{Mod}_{(A, \text{d})})}(N, A^\vee[k]) = 0
$$
임을 증명하면 충분하다. 이는 보조정리
\ref{dga-lemma-hom-into-shift-dual-free}와 $(-)^\vee$가 완전 함자라는
사실에서 따른다.
\end{proof}

\begin{lemma}
\label{dga-lemma-good-sub}
$(A, \text{d})$를 미분 등급 대수라 하자. $M$을 미분 등급
$A$-가군이라 하자. 다음 성질들을 갖는 미분 등급 $A$-가군의 준동형
$M \to I$가 존재한다.
\begin{enumerate}
\item $M \to I$는 단사이다.
\item $\Coker(\text{d}_M) \to \Coker(\text{d}_I)$는 단사이다.
\item $I$는 허용 짧은 완전열 $0 \to I' \to I \to I'' \to 0$에
들어가며, 여기서 $I'$, $I''$는 $A^\vee$의 시프트들의 곱이다.
\end{enumerate}
\end{lemma}

\begin{proof}
절 \ref{dga-section-modules-noncommutative-differential-graded}에서 구성한
함자 $N \mapsto N^\vee$(왼쪽 미분 등급 가군에서 오른쪽 미분 등급
가군으로, 또 오른쪽에서 왼쪽으로 가는 함자)와 그 모든 성질을 사용할
것이다. 모든 $k \in \mathbf{Z}$에 대해 $Q_k$를 차수 $k$와 $k + 1$의
생성원 $x, y$를 갖는 자유 왼쪽 $A$-가군이라 하자.
$\text{d}(x) = y$와 $\text{d}(y) = 0$으로 두어 $Q_k$ 위에 왼쪽
미분 등급 $A$-가군 구조를 정의한다. 보조정리
\ref{dga-lemma-good-quotient}의 증명과 정확히 같은 논증으로 왼쪽 미분 등급
$A$-가군의 전사
$$
\bigoplus\nolimits_{i \in I} Q_{k_i} \longrightarrow M^\vee
$$
를 얻는다. 그러면 단사
$$
M \to (M^\vee)^\vee \to (\bigoplus\nolimits_{i \in I} Q_{k_i})^\vee =
\prod\nolimits_{i \in I} I_{k_i}
$$
를 생각할 수 있다. 여기서
$I_k = Q_{-k}^\vee$는 “쌍대” 오른쪽 미분 등급 $A$-가군이다.
또 짧은 완전열 $0 \to A[-k - 1] \to Q_k \to A[-k] \to 0$은
짧은 완전열 $0 \to A^\vee[k] \to I_k \to A^\vee[k + 1] \to 0$을
만든다.

\medskip\noindent
앞 단락의 결과는 성질 (1)과 (3)을 갖는 $M \to I$를 만든다.
성질 (2)를 얻기 위해 차수 $k$의 영이 아닌 원소
$\overline{m} \in \Coker(\text{d}_M)$이 주어졌다고 하자. 사상
$\Im(M^{k - 1} \to M^k)$ 위에서는 영이지만 $m$에서는 영이 아닌
$\lambda : M^k \to \mathbf{Q}/\mathbf{Z}$를 택하자. 보조정리
\ref{dga-lemma-hom-into-shift-dual-free}에 의해 이는 $m$에서 영이 아닌
미분 등급 $A$-가군의 준동형 $M \to A^\vee[k]$에 대응한다.
따라서 $A^\vee$의 시프트들의 복사본들의 곱을 더하여 (2)를 얻을 수 있다.
\end{proof}

\begin{lemma}
\label{dga-lemma-right-resolution}
$(A, \text{d})$를 미분 등급 대수라 하자. $M$을 미분 등급
$A$-가군이라 하자. 다음을 만족하는 미분 등급 $A$-가군의 준동형
$M \to I$가 존재한다.
\begin{enumerate}
\item $M \to I$는 유사동형이다.
\item $I$는 성질 (I)를 갖는다.
\end{enumerate}
\end{lemma}

\begin{proof}
$M = M_0$으로 둔다. 보조정리 \ref{dga-lemma-good-sub}에서와 같이
사상 $M_i \to I_i$를 택하여 짧은 완전열
$$
0 \to M_i \to I_i \to M_{i + 1} \to 0
$$
들을 귀납적으로 택한다. 이로부터 “분해”
% P03-STACKS-E541: the frozen repeated I_1 is preserved.
$$
0 \to M \to I_0 \xrightarrow{f_0} I_1 \xrightarrow{f_1} I_1 \to \ldots
$$
를 얻는다. $I$를 등급 성분
$$
I^n = \prod\nolimits_{i \geq 0} I^{n - i}_i
$$
과 $x \in I_i^{n - i}$에 대해
$$
\text{d}_I(x) = f_i(x) + (-1)^i \text{d}_{I_i}(x)
$$
로 정의된 미분을 갖는 미분 등급 $A$-가군이라 하자. 이 규약에 따라
$I$는 실제로 미분 등급 $A$-가군이다. 각 $I_i$가 두 단계 여과
$0 \to I_i' \to I_i \to I_i'' \to 0$를 가짐을 상기하고,
$$
F_{2i}I^n = \prod\nolimits_{j \geq i} I^{n - j}_j
\subset
\prod\nolimits_{i \geq 0} I^{n - i}_i = I^n
$$
로 둔다. 그리고 인자 $I'_{i + 1}$을 $F_{2i}I$에 더하여
$F_{2i + 1}I$를 얻는다. 이들은 미분 등급 부분가군들이고 연속하는
몫들은 $A^\vee$의 시프트들의 곱이다. 보조정리
\ref{dga-lemma-source-graded-injective}에 의해 포함
$F_{i + 1}I \to F_iI$가 허용 단사상임을 알 수 있다. 마지막으로
증대 $M \to I_0$으로 주어진 사상 $M \to I$가 유사동형임을 보여야
한다. 이는 Homology, Lemma
\ref{homology-lemma-good-right-resolution-gives-qis}에서 따른다.
\end{proof}





\section{유도범주}
\label{dga-section-derived}

\noindent
비순환 미분 등급 가군과 미분 등급 가군의 유사동형이라는 개념이
정의되어 있음을 상기하자(\ref{dga-section-modules}절 참조).

\begin{lemma}
\label{dga-lemma-acyclic}
$(A, \text{d})$를 미분 등급 대수라 하자.
비순환 가군들로 이루어진 $K(\text{Mod}_{(A, \text{d})})$의
충만한 부분범주 $\text{Ac}$는 $K(\text{Mod}_{(A, \text{d})})$의
엄밀히 충만한 포화 삼각 부분범주이다.
이에 대응하는 $K(\text{Mod}_{(A, \text{d})})$의 포화 곱셈계
(Derived Categories, Lemma \ref{derived-lemma-operations} 참조)는
유사동형들의 모임 $\text{Qis}$이다. 특히 국소화 함자
$$
Q : K(\text{Mod}_{(A, \text{d})}) \to
\text{Qis}^{-1}K(\text{Mod}_{(A, \text{d})})
$$
의 핵은 $\text{Ac}$이다. 또한 함자 $H^0$는 $Q$를 통해 인수분해된다.
\end{lemma}

\begin{proof}
호몰로지 긴 완전열 (\ref{dga-equation-les})에 의해 $H^0$가
호몰로지적 함자임을 안다. $H^0$의 핵은 비순환 대상들의
부분범주이고, 모든 $H^i$에서 동형을 유도하는 화살표들은
유사동형이다. 따라서 이 보조정리는 Derived Categories,
Lemma \ref{derived-lemma-acyclic-general}의 특수한 경우이다.

\medskip\noindent
집합론적 주석. Derived Categories, Proposition
\ref{derived-proposition-construct-localization}의 국소화 구성은 주어진
삼각범주가 ``작다'', 즉 그 대상들의 기저 모임이 집합을 이룬다고
가정한다. $(A, \text{d})$를 포함하고 $\alpha$의 공종성이
$\aleph_0$보다 큰 부분 우주 $V_\alpha$를 택하자
(Sets, Section \ref{sets-section-sets-hierarchy} 및 Sets,
Proposition \ref{sets-proposition-exist-ordinals-large-cofinality} 참조).
그러면 $V_\alpha$에 들어 있는 미분 등급 $A$-가군들의 범주
$\text{Mod}_{(A, \text{d}), \alpha}$를 생각할 수 있다.
직접 확인하면 Proposition
\ref{dga-proposition-homotopy-category-triangulated}의 증명에 쓰인 모든 구성이
$\text{Mod}_{(A, \text{d}), \alpha}$ 안에서 작동함을 알 수 있다
(최악의 경우에도 미분 등급 가군들의 유한 직합만 취하기 때문이다).
따라서 삼각범주
$\text{Qis}_\alpha^{-1}K(\text{Mod}_{(A, \text{d}), \alpha})$를 얻는다.
아래에서 $\beta > \alpha$이면 전이 함자
$$
\text{Qis}_\alpha^{-1}K(\text{Mod}_{(A, \text{d}), \alpha})
\longrightarrow
\text{Qis}_\beta^{-1}K(\text{Mod}_{(A, \text{d}), \beta})
$$
들이 충실충만함을 보게 될 것이다. 실제로 몫범주의 사상 집합은
P-분해에서 나오는 호모토피 범주의 사상들로 계산된다
(Lemma \ref{dga-lemma-resolve}의 증명에서 P-분해를 구성할 때는 가산 직합뿐
아니라 주어진 가군의 부분집합들로 지표화된 직합도 취한다).
따라서 독자는 보조정리의 범주를 이 부분범주들의 합집합으로
생각해야 한다.
\end{proof}

\noindent
앞 보조정리의 증명 끝에 있는 집합론적 주석을 고려하여 유도범주를
다음과 같이 정의한다.

\begin{definition}
\label{dga-definition-unbounded-derived-category}
$(A, \text{d})$를 미분 등급 대수라 하자. $\text{Ac}$와
$\text{Qis}$는 Lemma \ref{dga-lemma-acyclic}에서와 같이 둔다.
{\it $(A, \text{d})$의 유도범주}는 삼각범주
$$
D(A, \text{d}) =
K(\text{Mod}_{(A, \text{d})})/\text{Ac} =
\text{Qis}^{-1}K(\text{Mod}_{(A, \text{d})}).
$$
몫함자와 합성하면 위에서 정의한 함자 $H^0$가 되는 유일한 함자를
$H^0 : D(A, \text{d}) \to \text{Mod}_R$로 나타낸다.
\end{definition}

\noindent
다음은 앞서 예고한, 유도범주의 사상 집합을 계산하는 보조정리이다.

\begin{lemma}
\label{dga-lemma-hom-derived}
$(A, \text{d})$를 미분 등급 대수라 하고, $M$과 $N$을 미분 등급
$A$-가군이라 하자.
\begin{enumerate}
\item $P \to M$을 Lemma \ref{dga-lemma-resolve}에서와 같은 P-분해라 하자.
그러면
$$
\Hom_{D(A, \text{d})}(M, N) =
\Hom_{K(\text{Mod}_{(A, \text{d})})}(P, N)
$$
\item $N \to I$를 Lemma \ref{dga-lemma-right-resolution}에서와 같은
I-분해라 하자. 그러면
$$
\Hom_{D(A, \text{d})}(M, N) =
\Hom_{K(\text{Mod}_{(A, \text{d})})}(M, I)
$$
\end{enumerate}
\end{lemma}

\begin{proof}
(1)에서와 같이 $P \to M$을 택하자. $P \to M$은 유사동형이므로
$$
\Hom_{D(A, \text{d})}(P, N) = \Hom_{D(A, \text{d})}(M, N)
$$
임을 유도범주의 정의에서 알 수 있다. $D(A, \text{d})$의 사상
$f : P \to N$은 $s^{-1}f'$와 같다. 여기서 $f' : P \to N'$은
사상이고 $s : N \to N'$은 유사동형이다. 특수 삼각형
$$
N \to N' \to Q \to N[1]
$$
을 택하자. $s$가 유사동형이므로 $Q$는 비순환이다. 따라서 Lemma
\ref{dga-lemma-property-P-K-projective}에 의해 모든 $k$에 대하여
$\Hom_{K(\text{Mod}_{(A, \text{d})})}(P, Q[k]) = 0$이다. 또한
$\Hom_{K(\text{Mod}_{(A, \text{d})})}(P, -)$
가 코호몰로지적 함자이므로 $f' : P \to N'$을 사상
$f : P \to N$으로 유일하게 올릴 수 있다. 이로써 증명이 끝난다.

\medskip\noindent
(2)의 증명은 Lemma \ref{dga-lemma-property-I-K-injective}을 쓰는
(1)의 쌍대이며, 여기서는 Lemma \ref{dga-lemma-property-P-K-projective}을
대신한다.
\end{proof}

\begin{lemma}
\label{dga-lemma-derived-products}
$(A, \text{d})$를 미분 등급 대수라 하자. 그러면
\begin{enumerate}
\item $D(A, \text{d})$에는 직합과 곱이 모두 존재한다.
\item 직합은 미분 등급 가군들의 직합을 취하여 얻는다.
\item 곱은 미분 등급 가군들의 곱을 취하여 얻는다.
\end{enumerate}
\end{lemma}

\begin{proof}
$\text{Mod}_{(A, \text{d})}$가 임의의 직합과 곱을 갖는 아벨 범주이고,
이들이 $K(\text{Mod}_{(A, \text{d})})$의 직합과 곱을 준다는 사실을
사용한다. Lemmas \ref{dga-lemma-dgm-abelian} 및
\ref{dga-lemma-homotopy-direct-sums}를 보라.

\medskip\noindent
$M_j$를 미분 등급 $A$-가군들의 족이라 하자. 자명한 구조를 갖는
미분 등급 $A$-가군인 등급 직합 $M = \bigoplus M_j$를 생각하자.
미분 등급 $A$-가군 $N$에 대하여, $I$가 성질 (I)를 갖는 미분 등급
$A$-가군이 되도록 유사동형 $N \to I$를 택하자. Lemma
\ref{dga-lemma-right-resolution}를 보라. Lemma \ref{dga-lemma-hom-derived}를 쓰면
\begin{align*}
\Hom_{D(A, \text{d})}(M, N)
& =
\Hom_{K(A, \text{d})}(M, I) \\
& =
\prod \Hom_{K(A, \text{d})}(M_j, I) \\
& =
\prod \Hom_{D(A, \text{d})}(M_j, N)
\end{align*}
을 얻는다. 따라서 보조정리의 (2)에 적힌 대로 $D(A, \text{d})$에
직합이 존재한다.

\medskip\noindent
$M_j$를 미분 등급 $A$-가군들의 족이라 하자. 미분 등급 $A$-가군들의
곱 $M = \prod M_j$를 생각하자. 미분 등급 $A$-가군 $N$에 대하여,
$P$가 성질 (P)를 갖는 미분 등급 $A$-가군이 되도록 유사동형
$P \to N$을 택하자. Lemma \ref{dga-lemma-resolve}를 보라. Lemma
\ref{dga-lemma-hom-derived}를 쓰면
\begin{align*}
\Hom_{D(A, \text{d})}(N, M)
& =
\Hom_{K(A, \text{d})}(P, M) \\
& =
\prod \Hom_{K(A, \text{d})}(P, M_j) \\
& =
\prod \Hom_{D(A, \text{d})}(N, M_j)
\end{align*}
을 얻는다. 따라서 보조정리의 (3)에 적힌 대로 $D(A, \text{d})$에
직합이 존재한다.
\end{proof}

\begin{remark}
\label{dga-remark-P-resolution}
$R$을 환이라 하고 $(A, \text{d})$를 미분 등급 $R$-대수라 하자.
P-분해를 사용하면 때때로 $D(A, \text{d})$의 일반적인 대상에 관한
명제를 $A[k]$에 관한 명제로 환원할 수 있다. 구체적으로 $T$를
$D(A, \text{d})$의 대상들의 성질이라 하고 다음을 가정하자.
\begin{enumerate}
\item $K_i$, $i \in I$가 $D(A, \text{d})$의 대상들의 족이고 모든
$i \in I$에 대하여 $T(K_i)$가 성립하면 $T(\bigoplus K_i)$가 성립한다.
\item $K \to L \to M \to K[1]$이 $D(A, \text{d})$의 특수
삼각형이고 두 대상에 대하여 $T$가 성립하면 셋째 대상에도 $T$가
성립한다.
\item 모든 $k \in \mathbf{Z}$에 대하여 $T(A[k])$가 성립한다.
\end{enumerate}
그러면 $D(A, \text{d})$의 모든 대상에 대하여 $T$가 성립한다. 이는
Lemmas \ref{dga-lemma-property-P-sequence} 및 \ref{dga-lemma-resolve}에서
명백하다.
\end{remark}







\section{표준 델타 함자}
\label{dga-section-canonical-delta-functor}

\noindent
$(A, \text{d})$를 미분 등급 대수라 하자. 함자
$\text{Mod}_{(A, \text{d})} \to K(\text{Mod}_{(A, \text{d})})$를
생각하자. 일반적으로 이 함자는 $\delta$-함자가 {\bf 아니다}.
그러나 함자 $\text{Mod}_{(A, \text{d})} \to D(A, \text{d})$는
$\delta$-함자임이 밝혀진다. 이를 보이려면 아벨 범주
$\text{Mod}_{(A, \text{d})}$의 짧은 완전열
$$
0 \to K \xrightarrow{a} L \xrightarrow{b} M \to 0
$$
에 결부되는 사상 $\delta$를 정의해야 한다. 사상 $a$의 원뿔
$C(a)$를 생각하자. $C(a) = L \oplus K$이고, $L$로의 사영 다음에
$b$를 합성하여 $q : C(a) \to M$을 정의한다. 이로써 미분 등급
$A$-가군의 준동형
$$
q : C(a) \longrightarrow M.
$$
을 얻는다. 여기서 $i$는 Definition \ref{dga-definition-cone}에서와 같으며,
$q \circ i = b$임은 명백하다. $a$가 단사이므로 $q$의 핵은 비순환인
$\text{id}_K$의 원뿔과 동일시된다. 따라서 $q$가 유사동형임을 알 수
있다. Lemma \ref{dga-lemma-the-same-up-to-isomorphisms}에 따르면 삼각형
$$
(K, L, C(a), a, i, -p)
$$
은 $K(\text{Mod}_{(A, \text{d})})$의 특수 삼각형이다. 국소화 함자
$K(\text{Mod}_{(A, \text{d})}) \to D(A, \text{d})$가 완전하므로
$(K, L, C(a), a, i, -p)$는 $D(A, \text{d})$의 특수 삼각형이다.
$q$가 유사동형이므로 $q$는 $D(A, \text{d})$에서 동형이다. 따라서
$$
(K, L, M, a, b, -p \circ q^{-1})
$$
이 $D(A, \text{d})$의 특수 삼각형임을 얻는다. 이는 다음 보조정리를
제시한다.

\begin{lemma}
\label{dga-lemma-derived-canonical-delta-functor}
$(A, \text{d})$를 미분 등급 대수라 하자. 위에서 정의한 함자
$\text{Mod}_{(A, \text{d})} \to D(A, \text{d})$에는 다음으로 주어지는
$\delta$-함자의 자연스러운 구조가 있다.
$$
\delta_{K \to L \to M} = - p \circ q^{-1}
$$
여기서 $p$와 $q$는 위에서 설명한 것과 같다.
\end{lemma}

\begin{proof}
이 선택이 복합체의 짧은 완전열마다 특수 삼각형을 준다는 것은 이미
보았다. 이 구성의 함자성을 보여야 한다. Derived Categories,
Definition \ref{derived-definition-delta-functor}을 보라. 이는 약간의
계산과 함께 Lemma \ref{dga-lemma-functorial-cone}에서 따른다. Derived
Categories, Lemma \ref{derived-lemma-derived-canonical-delta-functor}과
비교하라.
\end{proof}

\begin{lemma}
\label{dga-lemma-homotopy-colimit}
$(A, \text{d})$를 미분 등급 대수라 하고 $M_n$을 미분 등급 가군들의
계라 하자. 그러면 $D(A, \text{d})$의 유도 여극한
$\text{hocolim} M_n$은 미분 등급 가군 $\colim M_n$으로 표현된다.
\end{lemma}

\begin{proof}
$M = \colim M_n$이라 두자. 미분 등급 가군의 완전열
$$
0 \to \bigoplus M_n \to \bigoplus M_n \to M \to 0
$$
을 얻는다. 이는 아벨 군의 기저 복합체에 적용한 Derived Categories,
Lemma \ref{derived-lemma-compute-colimit}에서 따른다.
% P03-STACKS-E905: frozen source uses undefined D(\mathcal A).
Lemma \ref{dga-lemma-derived-products}에 의해 이 직합들은
$D(\mathcal{A})$에서의 직합이다. 따라서 Derived Categories,
Definition \ref{derived-definition-derived-colimit}의 유도 여극한 정의와
복합체의 짧은 완전열이 특수 삼각형을 준다는 사실
(Lemma \ref{dga-lemma-derived-canonical-delta-functor})에서 결론이 따른다.
\end{proof}








\section{선형 범주}
\label{dga-section-linear}

\noindent
정의만 제시한다.

\begin{definition}
\label{dga-definition-linear-category}
$R$을 환이라 하자. {\it $R$-선형 범주 $\mathcal{A}$}란 각 사상 집합에
$R$-가군 구조가 주어져 있고, $x, y, z \in \Ob(\mathcal{A})$에 대한
합성법칙
$$
\Hom_\mathcal{A}(y, z) \times \Hom_\mathcal{A}(x, y)
\longrightarrow
\Hom_\mathcal{A}(x, z)
$$
이 $R$-쌍선형인 범주를 말한다.
\end{definition}

\noindent
따라서 합성은 $R$-가군의 $R$-선형 사상
$$
\Hom_\mathcal{A}(y, z) \otimes_R \Hom_\mathcal{A}(x, y)
\longrightarrow
\Hom_\mathcal{A}(x, z)
$$
을 결정한다. $R$-선형 범주가 가법적이라고는 가정하지 않음에
유의하라.

\begin{definition}
\label{dga-definition-functor-linear-categories}
$R$을 환이라 하자. {\it $R$-선형 범주들의 함자}, 또는
{\it $R$-선형 함자}란 $\mathcal{A}$의 모든 대상 $x, y$에 대하여 사상
$F : \Hom_\mathcal{A}(x, y) \to \Hom_\mathcal{B}(F(x), F(y))$가
$R$-가군 준동형인 함자 $F : \mathcal{A} \to \mathcal{B}$를 말한다.
\end{definition}







\section{등급 범주}
\label{dga-section-graded}

\noindent
몇 가지 정의만 제시한다.

\begin{definition}
\label{dga-definition-graded-category}
$R$을 환이라 하자. {\it $R$ 위의 등급 범주 $\mathcal{A}$}란 각 사상
집합에 등급 $R$-가군 구조가 주어져 있고,
$x, y, z \in \Ob(\mathcal{A})$에 대하여 합성이 $R$-쌍선형이며 등급
$R$-가군의 준동형
$$
\Hom_\mathcal{A}(y, z) \otimes_R \Hom_\mathcal{A}(x, y)
\longrightarrow
\Hom_\mathcal{A}(x, z)
$$
을 유도하는 범주를 말한다(즉, 차수를 보존한다).
\end{definition}

\noindent
이 상황에서 등급 대상 $\Hom_\mathcal{A}(x, y)$의 차수 $i$ 부분을
$\Hom_\mathcal{A}^i(x, y)$로 나타낸다. 따라서
$$
\Hom_\mathcal{A}(x, y) =
\bigoplus\nolimits_{i \in \mathbf{Z}} \Hom_\mathcal{A}^i(x, y)
$$
는 등급 부분들로의 직합 분해이다.

\begin{definition}
\label{dga-definition-functor-graded-categories}
$R$을 환이라 하자. {\it $R$ 위의 등급 범주들의 함자}, 또는
{\it 등급 함자}란 $\mathcal{A}$의 모든 대상 $x, y$에 대하여 사상
% P03-STACKS-E908: frozen source retains Hom_A as the codomain.
$F : \Hom_\mathcal{A}(x, y) \to \Hom_\mathcal{A}(F(x), F(y))$
이 등급 $R$-가군의 준동형인 함자 $F : \mathcal{A} \to \mathcal{B}$를
말한다.
\end{definition}

\noindent
등급 범주가 주어지면 흔히 차수 $0$인 사상들의 대응하는 ``보통''
범주에 관심을 둔다. 이를 형식적으로 정의하자.

\begin{definition}
\label{dga-definition-H0-of-graded-category}
$R$을 환이라 하고 $\mathcal{A}$를 $R$ 위의 등급 범주라 하자.
{\it $\mathcal{A}^0$}를 $\mathcal{A}$와 같은 대상들을 가지며
$$
\Hom_{\mathcal{A}^0}(x, y) = \Hom^0_\mathcal{A}(x, y)
$$
를 사상 집합으로 갖는 범주라 하자. 우변은 $\mathcal{A}$의 사상들로
이루어진 등급 가군의 차수 $0$ 등급 부분이다.
\end{definition}

\begin{definition}
\label{dga-definition-graded-direct-sum}
$R$을 환이라 하고 $\mathcal{A}$를 $R$ 위의 등급 범주라 하자.
$\mathcal{A}$의 직합 $(x, y, z, i, j, p, q)$가
(표기는 Homology, Remark \ref{homology-remark-direct-sum}에서와 같다)
{\it 등급 직합}이라는 것은 $i, j, p, q$가 차수 $0$인 동차 사상임을
뜻한다.
\end{definition}

\begin{example}[등급 대상들의 등급 범주]
\label{dga-example-graded-category-graded-objects}
$\mathcal{B}$를 가법 범주라 하자. Homology, Definition
\ref{homology-definition-graded}에서 $\mathcal{B}$의 등급 대상들의 범주
$\text{Gr}(\mathcal{B})$를 정의했음을 상기하자. 이 예에서는
$R = \mathbf{Z}$ 위의 등급 범주 $\text{Gr}^{gr}(\mathcal{B})$를
구성하여, 그에 결부된 범주 $\text{Gr}^{gr}(\mathcal{B})^0$가
$\text{Gr}(\mathcal{B})$를 복원하도록 한다.
$\text{Gr}^{gr}(\mathcal{B})$의 대상으로는 $\mathcal{B}$의 등급 대상들을
취한다. 이제 $\mathcal{B}$의 등급 대상 $A = (A^i)$와 $B = (B^i)$가
주어지면
$$
\Hom_{\text{Gr}^{gr}(\mathcal{B})}(A, B) =
\bigoplus\nolimits_{n \in \mathbf{Z}} \Hom^n(A, B)
$$
로 둔다. 여기서 차수 $n$인 등급 부분은 $A$에서 $B$로 가는 차수
$n$의 동차 사상들로 이루어진 아벨 군이다. 구체적으로
$$
\Hom^n(A, B) = \prod\nolimits_{p + q = n} \Hom_\mathcal{B}(A^{-q}, B^p)
$$
이다(지표가 뒤집혔음과 여기에는 직합이 아니라 곱이 쓰였음에
유의하라). 다시 말해 $A$에서 $B$로 가는 차수 $n$의 사상 $f$는
$p, q \in \mathbf{Z}$이고 $p + q = n$이며
$f_{p, q} : A^{-q} \to B^p$가 $\mathcal{B}$의 사상이 되도록 하는 계
$f = (f_{p, q})$로 볼 수 있다. $\mathcal{B}$의 등급 대상 $A$, $B$,
$C$가 주어지면 $\text{Gr}^{gr}(\mathcal{B})$에서 사상의 합성은 사상
$$
\Hom^m(B, C) \times \Hom^n(A, B) \longrightarrow \Hom^{n + m}(A, C)
$$
을 등급 대상의 동차 사상들의 보통 합성 $(g, f) \mapsto g \circ f$로
정의한다. 성분으로 쓰면
$$
(g \circ f)_{p, r} = g_{p, q} \circ f_{-q, r}
$$
이며, 여기서 $q$는 $p + q = m$이고 $-q + r = n$이 되도록 택한다.
\end{example}

\begin{example}[등급 가군들의 등급 범주]
\label{dga-example-gm-gr-cat}
$A$를 환 $R$ 위의 $\mathbf{Z}$-등급 대수라 하자. 그에 결부된 범주
$(\text{Mod}^{gr}_A)^0$가 등급 $A$-가군들의 범주가 되도록 $R$ 위의
등급 범주 $\text{Mod}^{gr}_A$를 구성한다. $\text{Mod}^{gr}_A$의
대상으로 오른쪽 등급 $A$-가군들을 취한다
(Section \ref{dga-section-projectives-over-algebras} 참조).
% P03-STACKS-E909: frozen source retains the ungraded-section reference.
등급 $A$-가군 $L$과 $M$이 주어지면
$$
\Hom_{\text{Mod}^{gr}_A}(L, M) =
\bigoplus\nolimits_{n \in \mathbf{Z}} \Hom^n(L, M)
$$
로 둔다. 여기서 $\Hom^n(L, M)$은 차수 $n$의 동차 사상인 오른쪽
$A$-가군 사상 $L \to M$들의 집합이다. 즉 모든
$i \in \mathbf{Z}$에 대하여 $f(L^i) \subset M^{i + n}$이다. 성분으로
쓰면
$$
\Hom^n(L, M) \subset \prod\nolimits_{p + q = n} \Hom_R(L^{-q}, M^p)
$$
는 (지표가 뒤집혔음에 유의하라) 다음을 만족하는
$f = (f_{p, q})$들로 이루어진 부분집합이다.
$$
f_{p, q}(m a) = f_{p - i, q + i}(m)a
$$
여기서 $a \in A^i$이고 $m \in L^{-q - i}$이다. 등급 $A$-가군
$K$, $L$, $M$에 대하여 $\text{Mod}^{gr}_A$의 합성은 사상
$$
\Hom^m(L, M) \times \Hom^n(K, L) \longrightarrow \Hom^{n + m}(K, M)
$$
을 오른쪽 $A$-가군 사상들의 보통 합성
$(g, f) \mapsto g \circ f$로 정의한다.
\end{example}

\begin{remark}
\label{dga-remark-graded-shift-functors}
$R$을 환이라 하자. $\mathcal{D}$를 $R$-선형 범주라 하고,
$n \in \mathbf{Z}$로 지표화된 $R$-선형 함자들의 모임
$[n] : \mathcal{D} \to \mathcal{D}$, $x \mapsto x[n]$이 주어졌다고
하자. 이들은 $[n] \circ [m] = [n + m]$ 및
$[0] = \text{id}_\mathcal{D}$를 만족한다(함자로서의 등식이다).
그러면 $\mathcal{D}$와 같은 대상들을 갖는 $R$ 위의 등급 범주
$\mathcal{D}^{gr}$를 다음과 같이 구성할 수 있다.
$$
\Hom_{\mathcal{D}^{gr}}(x, y) =
\bigoplus\nolimits_{n \in \mathbf{Z}} \Hom_\mathcal{D}(x, y[n])
$$
이는 $\mathcal{D}$의 $x, y$에 대하여 정의한다.
$(\mathcal{D}^{gr})^0 = \mathcal{D}$임에 유의하라
(Definition \ref{dga-definition-H0-of-graded-category} 참조). 또한 등급 범주
$\mathcal{D}^{gr}$는 $[n] \circ [m] = [n + m]$ 및
$[0] = \text{id}_{\mathcal{D}^{gr}}$를 만족하는 $R$-선형 등급 함자
$[n]$들을 이어받으며, 이들은
$$
\Hom_{\mathcal{D}^{gr}}(x, y[n]) = \Hom_{\mathcal{D}^{gr}}(x, y)[n]
$$
이라는 성질을 갖는다. 이는 사상의 합성과 양립하는 등급
$R$-가군으로서의 등식이다.

\medskip\noindent
역으로, $R$ 위의 등급 범주 $\mathcal{A}$와 $R$-선형 등급 함자들의
모임 $[n]$이 주어졌다고 하자. 이들은
$[n] \circ [m] = [n + m]$ 및 $[0] = \text{id}_\mathcal{A}$를
만족하고, 더 나아가 동형
$$
\Hom_\mathcal{A}(x, y[n]) = \Hom_\mathcal{A}(x, y)[n]
$$
을 갖춘다고 하자. 이 동형은 사상의 합성과 양립하는 등급
$R$-가군의 동형이다. 그러면 독자는
$\mathcal{A} = (\mathcal{A}^0)^{gr}$임을 쉽게 보일 수 있다.

\medskip\noindent
위에서 확립한 대응 $\mathcal{D} \leftrightarrow \mathcal{A}$의 두 예를
들자.
\begin{enumerate}
\item $\mathcal{B}$를 가법 범주라 하자.
$\mathcal{D} = \text{Gr}(\mathcal{B})$이면 Example
\ref{dga-example-graded-category-graded-objects}에서와 같이
$\mathcal{A} = \text{Gr}^{gr}(\mathcal{B})$이다.
\item $A$가 등급 환이고 $\mathcal{D} = \text{Mod}_A$가 등급 오른쪽
$A$-가군들의 범주이면 $\mathcal{A} = \text{Mod}^{gr}_A$이다. Example
\ref{dga-example-gm-gr-cat}을 보라.
\end{enumerate}
\end{remark}






\section{미분 등급 범주}
\label{dga-section-dga-categories}

\noindent
$R$이 환이면 $R$ 자체는 (물론 $R = R^0$로 두어) 자기 자신 위의
미분 등급 대수임에 유의하라. 이 경우 미분 등급 $R$-가군은
$R$-가군들의 복합체와 같은 것이다. 특히 두 미분 등급 $R$-가군
$M$과 $N$이 주어졌을 때, $M$과 $N$에 결부된 $R$-가군 복합체들의
텐서곱으로 얻는 이중 복합체에 결부된 전체 복합체에 대응하는 미분
등급 $R$-가군을 $M \otimes_R N$으로 나타낸다.

\begin{definition}
\label{dga-definition-dga-category}
$R$을 환이라 하자. {\it $R$ 위의 미분 등급 범주 $\mathcal{A}$}란
각 사상 집합에 미분 등급 $R$-가군 구조가 주어져 있고,
$x, y, z \in \Ob(\mathcal{A})$에 대하여 합성이 $R$-쌍선형이며
미분 등급 $R$-가군의 준동형
$$
\Hom_\mathcal{A}(y, z) \otimes_R \Hom_\mathcal{A}(x, y)
\longrightarrow
\Hom_\mathcal{A}(x, z)
$$
을 유도하는 범주를 말한다.
\end{definition}

\noindent
정의의 마지막 조건은 다음을 뜻한다. 차수 $n$과 $m$의 동차 원소
$f \in \Hom_\mathcal{A}^n(x, y)$와
$g \in \Hom_\mathcal{A}^m(y, z)$에 대하여
$$
\text{d}(g \circ f) = \text{d}(g) \circ f + (-1)^mg \circ \text{d}(f)
$$
가 $\Hom_\mathcal{A}^{n + m + 1}(x, z)$에서 성립한다. 이는 이중
복합체의 전체 복합체 위의 미분에 대한 부호 규칙에서 따른다.
Homology, Definition \ref{homology-definition-associated-simple-complex}을 보라.

\begin{definition}
\label{dga-definition-functor-dga-categories}
$R$을 환이라 하자. {\it $R$ 위의 미분 등급 범주들의 함자}란
$\mathcal{A}$의 모든 대상 $x, y$에 대하여 사상
% P03-STACKS-E912: frozen source retains Hom_A as the codomain.
$F : \Hom_\mathcal{A}(x, y) \to \Hom_\mathcal{A}(F(x), F(y))$
이 미분 등급 $R$-가군의 준동형인 함자
$F : \mathcal{A} \to \mathcal{B}$를 말한다.
\end{definition}

\noindent
미분 등급 범주가 주어지면 흔히 이에 대응하는 복합체 범주와 호모토피
범주에 관심을 둔다. 이를 형식적으로 정의하자.

\begin{definition}
\label{dga-definition-homotopy-category-of-dga-category}
$R$을 환이라 하고 $\mathcal{A}$를 $R$ 위의 미분 등급 범주라 하자.
다음과 같이 둔다.
\begin{enumerate}
\item {\it $\mathcal{A}$의 복합체 범주}\footnote{이는 비표준 용어일 수
있다.}를 $\mathcal{A}$와 같은 대상들을 가지며
$$
\Hom_{\text{Comp}(\mathcal{A})}(x, y) =
\Ker(d : \Hom^0_\mathcal{A}(x, y) \to \Hom^1_\mathcal{A}(x, y))
$$
를 사상 집합으로 갖는 범주 $\text{Comp}(\mathcal{A})$라 한다.
\item {\it $\mathcal{A}$의 호모토피 범주}를 $\mathcal{A}$와 같은
대상들을 가지며
$$
\Hom_{K(\mathcal{A})}(x, y) = H^0(\Hom_\mathcal{A}(x, y))
$$
를 사상 집합으로 갖는 범주 $K(\mathcal{A})$라 한다.
\end{enumerate}
\end{definition}

\noindent
기호 $K(\mathcal{A})$의 이러한 사용은 비표준이지만, 적어도 Stacks
project의 다른 장에서 $K(-)$를 사용하는 방식과 양립한다.

\begin{definition}
\label{dga-definition-dg-direct-sum}
$R$을 환이라 하고 $\mathcal{A}$를 $R$ 위의 미분 등급 범주라 하자.
$\mathcal{A}$의 직합 $(x, y, z, i, j, p, q)$가
(표기는 Homology, Remark \ref{homology-remark-direct-sum}에서와 같다)
{\it 미분 등급 직합}이라는 것은 $i, j, p, q$가 차수 $0$인 닫힌 동차
사상, 즉 $\text{d}(i) = 0$ 등을 만족하는 사상임을 뜻한다.
\end{definition}

\begin{lemma}
\label{dga-lemma-functorial}
$R$을 환이라 하자. $R$ 위의 미분 등급 범주들의 함자
$F : \mathcal{A} \to \mathcal{B}$는 함자
$\text{Comp}(\mathcal{A}) \to \text{Comp}(\mathcal{B})$
와 $K(\mathcal{A}) \to K(\mathcal{B})$를 유도한다.
\end{lemma}

\begin{proof}
생략한다.
\end{proof}

\begin{example}[복합체들의 미분 등급 범주]
\label{dga-example-category-complexes}
$\mathcal{B}$를 가법 범주라 하자. 이에 결부된 복합체 범주는
$\text{Comp}(\mathcal{B})$이고 이에 결부된 호모토피 범주는
$K(\mathcal{B})$인, $R = \mathbf{Z}$ 위의 미분 등급 범주
$\text{Comp}^{dg}(\mathcal{B})$를 구성하겠다.
$\text{Comp}^{dg}(\mathcal{B})$의 대상으로 $\mathcal{B}$의 복합체들을
취한다. $\mathcal{B}$의 복합체 $A^\bullet$과 $B^\bullet$이 주어지면,
때때로 $\mathcal{B}$의 대응하는 등급 대상들도 $A^\bullet$과
$B^\bullet$로 나타낸다(즉 미분을 잊는다). 이 표기 남용 아래
$$
\Hom_{\text{Comp}^{dg}(\mathcal{B})}(A^\bullet, B^\bullet) =
\Hom_{\text{Gr}^{gr}(\mathcal{B})}(A^\bullet, B^\bullet) =
\bigoplus\nolimits_{n \in \mathbf{Z}} \Hom^n(A, B)
$$
로 둔다. 이는 Example \ref{dga-example-graded-category-graded-objects}의
표기와 정의를 사용한 등급 $\mathbf{Z}$-가군이다. 다시 말해 차수
$n$인 등급 부분은 등급 대상들의 차수 $n$ 동차 사상들로 이루어진
아벨 군
$$
\Hom^n(A^\bullet, B^\bullet) =
\prod\nolimits_{p + q = n} \Hom_\mathcal{B}(A^{-q}, B^p)
$$
이다. 지표가 뒤집혔고 직합이 아니라 곱을 취한다는 점에 유의하라.
차수 $n$인 원소 $f \in \Hom^n(A^\bullet, B^\bullet)$에 대하여
$$
\text{d}(f) = \text{d}_B \circ f - (-1)^n f \circ \text{d}_A
$$
로 둔다. 부호는 More on Algebra, Section
\ref{more-algebra-section-sign-rules}에서와 정확히 같다. 이를 해석하기
위해 $\text{d}_B$와 $\text{d}_A$를 $\mathcal{B}$의 등급 대상들의 차수
$1$ 동차 사상으로 보고, 우변에서는 범주
$\text{Gr}^{gr}(\mathcal{B})$의 합성을 사용한다. 성분으로 쓰면,
$f = (f_{p, q})$이고 $f_{p, q} : A^{-q} \to B^p$일 때
\begin{equation}
\label{dga-equation-differential-hom-complex}
\text{d}(f_{p, q}) =
\text{d}_B \circ f_{p, q} - (-1)^{p + q} f_{p, q} \circ \text{d}_A 
\end{equation}
이다. 이 식의 첫째 항은 $\Hom_\mathcal{B}(A^{-q}, B^{p + 1})$에
속하고 둘째 항은 $\Hom_\mathcal{B}(A^{-q - 1}, B^p)$에 속함에
유의하라. 독자는 다음을 확인한다.
\begin{enumerate}
\item $\text{d}$의 제곱은 영이다.
\item $\Hom^n(A^\bullet, B^\bullet)$의 원소 $f$에 대하여
$\text{d}(f) = 0$일 필요충분조건은 $\mathcal{B}$의 등급 대상들의 사상
$f : A^\bullet \to B^\bullet[n]$가 실제로 복합체의 사상인 것이다.
\item 특히 $\text{Comp}^{dg}(\mathcal{B})$의 복합체 범주는
$\text{Comp}(\mathcal{B})$와 같다.
\item (2)에서와 같이 $f$가 정의하는 복합체의 사상이 영과 호모토피
동치일 필요충분조건은 어떤
$g \in \Hom^{n - 1}(A^\bullet, B^\bullet)$에 대하여
$f = \text{d}(g)$인 것이다.
\item 특히 표준 동형
$$
\Hom_{K(\mathcal{B})}(A^\bullet, B^\bullet)
\longrightarrow
H^0(\Hom_{\text{Comp}^{dg}(\mathcal{B})}(A^\bullet, B^\bullet))
$$
을 얻으며, $\text{Comp}^{dg}(\mathcal{B})$의 호모토피 범주는
$K(\mathcal{B})$와 같다.
\end{enumerate}
복합체 $A^\bullet$, $B^\bullet$, $C^\bullet$이 주어지면 합성
$$
\Hom^m(B^\bullet, C^\bullet) \times \Hom^n(A^\bullet, B^\bullet)
\longrightarrow
\Hom^{n + m}(A^\bullet, C^\bullet)
$$
을 등급 범주 $\text{Gr}^{gr}(\mathcal{B})$의 합성
$(g, f) \mapsto g \circ f$로 정의한다. Example
\ref{dga-example-graded-category-graded-objects}을 보라. 이는 미분 등급
가군의 사상
$$
\Hom_{\text{Comp}^{dg}(\mathcal{B})}(B^\bullet, C^\bullet)
\otimes_R
\Hom_{\text{Comp}^{dg}(\mathcal{B})}(A^\bullet, B^\bullet)
\longrightarrow
\Hom_{\text{Comp}^{dg}(\mathcal{B})}(A^\bullet, C^\bullet)
$$
을 정의하며, 이는 Definition \ref{dga-definition-dga-category}에서 요구하는
사상이다. 실제로
\begin{align*}
\text{d}(g \circ f) & =
\text{d}_C \circ g \circ f - (-1)^{n + m} g \circ f \circ \text{d}_A \\
& =
\left(\text{d}_C \circ g - (-1)^m g \circ \text{d}_B\right) \circ f +
(-1)^m g \circ \left(\text{d}_B \circ f - (-1)^n f \circ \text{d}_A\right) \\
& =
\text{d}(g) \circ f + (-1)^m g \circ \text{d}(f)
\end{align*}
이므로 원하는 결론을 얻는다.
\end{example}

\begin{lemma}
\label{dga-lemma-additive-functor-induces-dga-functor}
$F : \mathcal{B} \to \mathcal{B}'$를 가법 범주들 사이의 가법
함자라 하자. 그러면 $F$는 미분 등급 범주들의 함자
$$
F : \text{Comp}^{dg}(\mathcal{B}) \to \text{Comp}^{dg}(\mathcal{B}')
$$
를 유도한다. 이는 Example \ref{dga-example-category-complexes}의 미분 등급
범주들 사이의 함자이며 복합체 범주와 호모토피 범주 위의 통상적인
함자들을 유도한다.
\end{lemma}

\begin{proof}
생략한다.
\end{proof}

\begin{example}[미분 등급 가군들의 미분 등급 범주]
\label{dga-example-dgm-dg-cat}
$(A, \text{d})$를 환 $R$ 위의 미분 등급 대수라 하자. 복합체 범주는
$\text{Mod}_{(A, \text{d})}$이고 호모토피 범주는
$K(\text{Mod}_{(A, \text{d})})$인, $R$ 위의 미분 등급 범주
$\text{Mod}^{dg}_{(A, \text{d})}$를 구성하겠다.
$\text{Mod}^{dg}_{(A, \text{d})}$의 대상으로 미분 등급 $A$-가군들을
취한다. 미분 등급 $A$-가군 $L$과 $M$이 주어지면
% P03-STACKS-E916: frozen source retains the unindexed direct sum.
$$
\Hom_{\text{Mod}^{dg}_{(A, \text{d})}}(L, M) =
\Hom_{\text{Mod}^{gr}_A}(L, M) = \bigoplus \Hom^n(L, M)
$$
로 둔다. 이는 우변을 Example \ref{dga-example-gm-gr-cat}에서와 같이 정의한
등급 $R$-가군이다. 다시 말해 차수 $n$인 등급 부분
$\Hom^n(L, M)$은 차수 $n$의 동차 오른쪽 $A$-가군 사상들로 이루어진
$R$-가군이다. 원소 $f \in \Hom^n(L, M)$에 대하여
$$
\text{d}(f) = \text{d}_M \circ f - (-1)^n f \circ \text{d}_L
$$
로 둔다. 이를 해석하기 위해 $\text{d}_M$과 $\text{d}_L$을 등급
$R$-가군 사상으로 보고 등급 $R$-가군 사상들의 합성을 사용한다.
$\text{d}(f)$가 등급 $R$-가군 사상으로서 차수 $n + 1$의 동차 사상임은
명백하며, 다음 계산에 의해 $A$-선형이다.
\begin{align*}
\text{d}(f)(xa)
& =
\text{d}_M(f(x) a) - (-1)^n f (\text{d}_L(xa)) \\
& =
\begin{aligned}
&\text{d}_M(f(x)) a + (-1)^{\deg(x) + n} f(x) \text{d}(a) \\
&- (-1)^n f(\text{d}_L(x)) a - (-1)^{n + \deg(x)} f(x) \text{d}(a)
\end{aligned} \\
& = \text{d}(f)(x) a
\end{align*}
원하는 식을 얻었다($\Hom$ 위의 미분 정의에 들어 있는 부호가 없으면
이 계산은 성립하지 않음에 유의하라). $f$의 미분을 성분으로 쓴 식에
대해서는 Example \ref{dga-example-category-complexes}의 식들과 유사한
공식들이 성립한다. 독자는 Example \ref{dga-example-category-complexes}에서와
같은 방식으로 다음을 확인한다.
\begin{enumerate}
\item $\text{d}$의 제곱은 영이다.
\item $\Hom^n(L, M)$의 원소 $f$에 대하여 $\text{d}(f) = 0$일
필요충분조건은 $f : L \to M[n]$이 미분 등급 $A$-가군의 준동형인
것이다.
\item 특히 $\text{Mod}^{dg}_{(A, \text{d})}$의 복합체 범주는
$\text{Mod}_{(A, \text{d})}$이다.
\item (2)에서와 같이 $f$가 정의하는 준동형이 영과 호모토피 동치일
필요충분조건은 어떤 $g \in \Hom^{n - 1}(L, M)$에 대하여
$f = \text{d}(g)$인 것이다.
\item 특히 표준 동형
$$
\Hom_{K(\text{Mod}_{(A, \text{d})})}(L, M)
\longrightarrow
H^0(\Hom_{\text{Mod}^{dg}_{(A, \text{d})}}(L, M))
$$
을 얻으며, $\text{Mod}^{dg}_{(A, \text{d})}$의 호모토피 범주는
$K(\text{Mod}_{(A, \text{d})})$이다.
\end{enumerate}
미분 등급 $A$-가군 $K$, $L$, $M$에 대하여 합성
$$
\Hom^m(L, M) \times \Hom^n(K, L) \longrightarrow \Hom^{n + m}(K, M)
$$
을 동차 오른쪽 $A$-가군 사상들의 합성
$(g, f) \mapsto g \circ f$로 정의한다. 이는 미분 등급 가군의 사상
$$
\Hom_{\text{Mod}^{dg}_{(A, \text{d})}}(L, M) \otimes_R
\Hom_{\text{Mod}^{dg}_{(A, \text{d})}}(K, L) \longrightarrow
\Hom_{\text{Mod}^{dg}_{(A, \text{d})}}(K, M)
$$
을 정의하며, 이는 Definition \ref{dga-definition-dga-category}에서 요구하는
사상이다. 실제로
\begin{align*}
\text{d}(g \circ f) & =
\text{d}_M \circ g \circ f - (-1)^{n + m} g \circ f \circ \text{d}_K \\
& =
\left(\text{d}_M \circ g - (-1)^m g \circ \text{d}_L\right) \circ f +
(-1)^m g \circ \left(\text{d}_L \circ f - (-1)^n f \circ \text{d}_K\right) \\
& =
\text{d}(g) \circ f + (-1)^m g \circ \text{d}(f)
\end{align*}
이므로 원하는 결론을 얻는다.
\end{example}

\begin{lemma}
\label{dga-lemma-homomorphism-induces-dga-functor}
$\varphi : (A, \text{d}) \to (E, \text{d})$를 미분 등급 대수의
준동형이라 하자. 그러면 $\varphi$는 미분 등급 범주들의 함자
$$
F :
\text{Mod}^{dg}_{(E, \text{d})}
\longrightarrow
\text{Mod}^{dg}_{(A, \text{d})}
$$
를 유도한다. 이는 Example \ref{dga-example-dgm-dg-cat}의 미분 등급 범주들
사이의 함자이며, 미분 등급 가군의 범주들과 호모토피 범주들 위의
자명한 제한 함자들을 유도한다.
\end{lemma}

\begin{proof}
생략한다.
\end{proof}

\begin{lemma}
\label{dga-lemma-construction}
$R$을 환이라 하고 $\mathcal{A}$를 $R$ 위의 미분 등급 범주라 하자.
$x$를 $\mathcal{A}$의 대상이라 하고
$$
(E, \text{d}) = \Hom_\mathcal{A}(x, x)
$$
를 $x$의 자기 사상들로 이루어진 미분 등급 $R$-대수라 하자.
$\mathcal{A}$에서의 합성을 통해 $E$가 $\Hom_\mathcal{A}(x, y)$에
작용하도록 하면 미분 등급
범주들의 함자
$$
\mathcal{A} \longrightarrow \text{Mod}^{dg}_{(E, \text{d})},\quad
y \longmapsto \Hom_\mathcal{A}(x, y)
$$
를 얻는다. 이 함자는 Lemma \ref{dga-lemma-functorial}을 적용하여 함자
% P03-STACKS-E918: frozen source retains the undefined algebra A twice.
$$
\text{Comp}(\mathcal{A}) \to \text{Mod}_{(A, \text{d})}
\quad\text{and}\quad
K(\mathcal{A}) \to K(\text{Mod}_{(A, \text{d})})
$$
를 유도한다.
\end{lemma}

\begin{proof}
이 보조정리의 주장이 곧 증명이다.
\end{proof}









\section{삼각범주를 얻는 방법}
\label{dga-section-review}

\noindent
이 절에서는 Derived Categories, Proposition
\ref{derived-proposition-homotopy-category-triangulated}과 Proposition
\ref{dga-proposition-homotopy-category-triangulated}을 증명하는 논증이
적용되는 가장 일반적인 설정을 논의한다.

\medskip\noindent
$R$을 환이라 하고 $\mathcal{A}$를 $R$ 위의 미분 등급 범주라 하자.
논증이 성립하도록 $\mathcal{A}$에 다음 공리들을 부과한다.
\begin{enumerate}
\item[(A)] $\mathcal{A}$는 영대상과 두 대상의 미분 등급 직합을 갖는다
(Definition \ref{dga-definition-dg-direct-sum}에서와 같다).
\item[(B)] $[0] = \text{id}_\mathcal{A}$ 및
$[n + m] = [n] \circ [m]$을 만족하는 미분 등급 범주들의 함자
$[n] : \mathcal{A} \longrightarrow \mathcal{A}$가 있고, 합성과 양립하는
미분 등급 $R$-가군의 동형
$$
\Hom_\mathcal{A}(x, y[n]) = \Hom_\mathcal{A}(x, y)[n]
$$
이 주어져 있다.
\end{enumerate}

\noindent
이 미분 등급 범주 $\mathcal{A}$에 대하여 다음과 같이 말한다.
\begin{enumerate}
\item $\text{Comp}(\mathcal{A})$의 사상들의 열 $x \to y \to z$가
{\it 허용 짧은 완전열}이라는 것은, 바탕 등급 범주에서 동형
$y \cong x \oplus z$가 존재하여
% P03-STACKS-E551: frozen source retains x -> z as the first (co)projection.
$x \to z$와 $y \to z$가 (여)사영이 됨을 뜻한다.
\item $\text{Comp}(\mathcal{A})$의 사상 $x\to y$가
{\it 허용 단사상}이라는 것은 허용 짧은 완전열 $x\to y\to z$로
연장됨을 뜻한다.
\item $\text{Comp}(\mathcal{A})$의 사상 $y\to z$가
{\it 허용 전사상}이라는 것은 허용 짧은 완전열 $x\to y\to z$로
연장됨을 뜻한다.
\end{enumerate}
다음 보조정리는 공리 (A)와 (B)가 성립할 때 허용 짧은 완전열이
삼각형을 준다는 것을 말한다.

\begin{lemma}
\label{dga-lemma-get-triangle}
$\mathcal{A}$를 공리 (A)와 (B)를 만족하는 미분 등급 범주라 하자.
허용 짧은 완전열 $x \to y \to z$가 주어지면 (증명을 보라) 삼각형
$$
x \to y \to z \to x[1]
$$
을 $\text{Comp}(\mathcal{A})$에서 얻으며,
$z[-1] \to x \to y \to z \to x[1]$에서 연속한 두 사상의 합성은
모두 $K(\mathcal{A})$에서 영이다.
\end{lemma}

\begin{proof}
허용 짧은 완전열의 정의에서와 같이 등급 대상들의 동형
$y \cong x \oplus z$를 주는 도식
$$
\xymatrix{
x \ar[rr]_1 \ar[rd]_a & & x \\
& y \ar[ru]_\pi \ar[rd]^b & \\
z \ar[rr]^1 \ar[ru]^s & & z
}
$$
을 택하자. 이 상황에서는 다음 등식들이 성립한다.
\begin{enumerate}
\item $1 = \pi a$이고, 따라서 $\text{d}(\pi) a = 0$이다.
\item $1 = b s$이고, 따라서 $b \text{d}(s) = 0$이다.
\item $1 = a \pi + s b$이고, 따라서
$a \text{d}(\pi) + \text{d}(s) b = 0$이다.
\item $\pi s = 0$이고, 따라서
$\text{d}(\pi)s + \pi \text{d}(s) = 0$이다.
\item $\text{d}(s) = (a \pi + s b)\text{d}(s)$ 및
$b\text{d}(s) = 0$이므로 $\text{d}(s) = a \pi \text{d}(s)$이다.
\item $\text{d}(\pi) = \text{d}(\pi)(a \pi + s b)$ 및
$\text{d}(\pi)a = 0$이므로
$\text{d}(\pi) = \text{d}(\pi) s b$이다.
\item $\text{d}(\pi \text{d}(s)) = 0$이다. 실제로 이를 단사상 $a$와
뒤에서 합성하면
$\text{d}(a\pi \text{d}(s)) = \text{d}(\text{d}(s)) = 0$을 얻는다.
\item $\text{d}(\text{d}(\pi)s) = 0$이다. 실제로 (4)에 의해 이는
$\text{d}(\pi\text{d}(s))$의 음수이고, 후자는 (7)에 의해 $0$이다.
\end{enumerate}
$\text{d}(a) = 0$, $\text{d}(b) = 0$, $\text{d}(1) = 0$을 반복해서
사용하였다. (7)에 의해
$$
\delta = \pi \text{d}(s) = - \text{d}(\pi) s : z \to x[1]
$$
가 $\text{Comp}(\mathcal{A})$의 사상임을 알 수 있다. (5)에 의해 합성
$a \delta = a \pi \text{d}(s) = \text{d}(s)$는 영과 호모토픽이다.
(6)에 의해 합성 $\delta b = - \text{d}(\pi)sb = \text{d}(-\pi)$도
영과 호모토픽이다.
\end{proof}

\noindent
공리 (A)와 (B) 외에도 원뿔의 존재에 관한 공리가 필요하다. 모든 자료를
다음과 같이 형식화한다.

\begin{situation}
\label{dga-situation-ABC}
여기서 $R$은 환이고 $\mathcal{A}$는 공리 (A), (B), 그리고 다음 공리를
만족하는 $R$ 위의 미분 등급 범주이다.
\begin{enumerate}
\item[(C)] $\text{d}(f) = 0$인 차수 $0$의 화살표 $f : x \to y$가
주어지면, Lemma \ref{dga-lemma-get-triangle}의 사상 $x[1] \to y[1]$이
$f[1]$과 같아지도록 하는 $\text{Comp}(\mathcal{A})$의 허용 짧은 완전열
$y \to c(f) \to x[1]$이 존재한다.
\end{enumerate}
\end{situation}

\noindent
$c(f)$를 사상 $f$의 {\it 원뿔}이라 부르겠다. (A), (B), (C)가
성립하면 원뿔은 약한 의미에서 함자적이다.

\begin{lemma}
\label{dga-lemma-cone}
\begin{slogan}
호모토피 범주는 삼각범주이다. 이 보조정리는 삼각범주 공리들의 일부를
증명한다.
\end{slogan}
Situation \ref{dga-situation-ABC}에서 도식
$$
\xymatrix{
x_1 \ar[r]_{f_1} \ar[d]_a & y_1 \ar[d]^b \\
x_2 \ar[r]^{f_2} & y_2
}
$$
이 $\text{Comp}(\mathcal{A})$에서 호모토피까지 가환한다고 하자. 그러면
삼각형들의 사상
% P03-STACKS-E552: frozen source repeats the source triangle as the target.
$$
(a, b, c) : (x_1, y_1, c(f_1)) \to (x_1, y_1, c(f_1))
$$
을 $K(\mathcal{A})$에서 주는 사상 $c : c(f_1) \to c(f_2)$가 존재한다.
\end{lemma}

\begin{proof}
가정은 $\text{d}(h) = b f_1 - f_2 a$를 만족하는 차수 $-1$의 사상
$h : x_1 \to y_2$가 존재한다는 뜻이다. 사상
$y_i \to c(f_i) \to x_i[1]$과 양립하는 등급 대상들의 동형
$c(f_i) = y_i \oplus x_i[1]$을 택하자. 주어진 사상들을
$a_i : y_i \to c(f_i)$, $b_i : c(f_i) \to x_i[1]$,
$s_i : x_i[1] \to c(f_i)$, $\pi_i : c(f_i) \to y_i$로 나타내자.
$x_i[1] \to y_i[1]$은 $\pi_i \text{d}(s_i)$로 주어진다. 공리 (C)에
의해 이는
$$
f_i = \pi_i \text{d}(s_i) = - \text{d}(\pi_i) s_i
$$
라는 뜻이다(시프트 함자 $[1]$을 사용하여 $\Hom(x_i, y_i)$와
$\Hom(x_i[1], y_i[1])$을 동일시한다).
$c = a_2 b \pi_1 + s_2 a b_1 + a_2hb$로 두자. Lemma
\ref{dga-lemma-get-triangle}의 증명에서 얻은 등식들을 사용하면
\begin{align*}
\text{d}(c)
& =
a_2 b \text{d}(\pi_1) + \text{d}(s_2) a b_1 + a_2 \text{d}(h) b_1 \\
& =
- a_2 b f_1 b_1 + a_2 f_2 a b_1 + a_2 (b f_1 - f_2 a) b_1 \\
& = 0
\end{align*}
을 얻는다(특히
$\text{d}(\pi_1) = \text{d}(\pi_1) s_1 b_1 = f_1 b_1$ 및
$\text{d}(s_2) = a_2 \pi_2 \text{d}(s_2) = a_2 f_2$를 사용하였다).
따라서 $c$는 주어진 사상들 $y_i \to c(f_i) \to x_i[1]$과 양립하는
$\mathcal{A}$의 차수 $0$ 사상 $c : c(f_1) \to c(f_2)$이다.
\end{proof}

\noindent
Situation \ref{dga-situation-ABC}에서 $K(\mathcal{A})$의 삼각형
$(x, y, z, f, g, h)$가 {\it 특수 삼각형}이라는 것은, 허용 짧은 완전열
$x' \to y' \to z'$가 존재하여 $(x, y, z, f, g, h)$가
$K(\mathcal{A})$의 삼각형으로서 Lemma \ref{dga-lemma-get-triangle}에서 구성한
삼각형 $(x', y', z', x' \to y', y' \to z', \delta)$와 동형임을 뜻한다.
아래에서
$$
\boxed{
K(\mathcal{A})\text{ is a triangulated category}
}
$$
임을 보이겠다. 이 결과는 생각만큼 일반적이지는 않지만 Stacks
project에서 지금까지 다룬 경우들의 여러 자연스러운 일반화에 적용된다.
몇 가지 예를 들자.
\begin{enumerate}
\item $(X, \mathcal{O}_X)$를 환 달린 공간이라 하고 $(A, d)$를 미분 등급
$\mathcal{O}_X$-대수들의 층이라 하자. $\mathcal{A}$를 미분 등급
$A$-가군들의 미분 등급 범주라 하면 $K(\mathcal{A})$는 삼각범주이다.
\item $(\mathcal{C}, \mathcal{O})$를 환 달린 사이트라 하고 $(A, d)$를
미분 등급 $\mathcal{O}$-대수들의 층이라 하자. $\mathcal{A}$를 미분
등급 $A$-가군들의 미분 등급 범주라 하면 $K(\mathcal{A})$는
삼각범주이다. Differential Graded Sheaves, Proposition
\ref{sdga-proposition-homotopy-category-triangulated}을 보라.
\item 성격이 다른 두 예는 Examples, Section
\ref{examples-section-nongraded-differential-graded}에서 찾을 수 있다.
\end{enumerate}

\noindent
다음의 간단한 보조정리가 구성의 핵심이다.

\begin{lemma}
\label{dga-lemma-id-cone-null}
Situation \ref{dga-situation-ABC}에서 $\mathcal{A}$의 임의의 대상 $x$와
항등 사상 $1_x : x \to x$의 원뿔 $C(1_x)$가 주어지면,
$C(1_x)$ 위의 항등 사상은 영과 호모토픽이다.
\end{lemma}

\begin{proof}
공리 (C)가 주는 허용 짧은 완전열을 생각하자.
$$
\xymatrix{
x \ar@<0.5ex>[r]^a  &
C(1_x) \ar@<0.5ex>[l]^{\pi} \ar@<0.5ex>[r]^b &
x[1]\ar@<0.5ex>[l]^s
}
$$
그러면 시프트 아래에서 Hom 집합들을 동일시할 때 Lemma
\ref{dga-lemma-get-triangle}에 의해 $1_x=\pi d(s)=-d(\pi)s$이다. 여기서
$s$는 $\Hom_{\mathcal{A}}^{-1}(x,C(1_x))$의 사상으로 본다. 따라서
Lemma \ref{dga-lemma-get-triangle}의 공식 (5)를 사용하면
$a=a\pi d(s)=d(s)$이고, Lemma \ref{dga-lemma-get-triangle}의 공식 (6)에 의해
$b=-d(\pi)sb=-d(\pi)$이다. 그러므로
$$
1_{C(1_x)} = a\pi + sb = d(s)\pi - sd(\pi) = d(s\pi)
$$
이다. 여기서 $s$의 차수는 $-1$이다.
\end{proof}

\noindent
위 보조정리의 더 일반적인 판은 Lemma \ref{dga-lemma-cone-homotopy}에
나온다. 다음 보조정리는 Lemma \ref{dga-lemma-make-commute-map}의 유사물이다.

\begin{lemma}
\label{dga-lemma-homo-change}
Situation \ref{dga-situation-ABC}에서 도식
$$
\xymatrix{x\ar[r]^f\ar[d]_a & y\ar[d]^b\\
z\ar[r]^g & w}
$$
이 $\text{Comp}(\mathcal{A})$에서 호모토피까지 가환한다고 하자. 그러면
\begin{enumerate}
\item $f$가 허용 단사상이면 $b$는 이 도식을 가환하게 만드는 사상
$b'$와 호모토픽이다.
\item $g$가 허용 전사상이면 $a$는 이 도식을 가환하게 만드는 사상
$a'$와 호모토픽이다.
\end{enumerate}
\end{lemma}

\begin{proof}
(1)을 증명하자. 가정에 의해 차수 $-1$인 어떤
$h\in\Hom_{\mathcal{A}}(x,w)$가 존재하여 $bf-ga=d(h)$이다. $f$가 허용
단사상이므로 범주 $\mathcal{A}$에 차수 $0$의 사상
$\pi : y \to x$가 존재한다. $b' = b - d(h\pi)$로 두자. 그러면
\begin{align*}
b'f = bf - d(h\pi)f
= &
bf - d(h\pi f) \quad (\text{since }d(f) = 0) \\
= &
bf-d(h) \\
= &
ga
\end{align*}
이므로 원하는 결론을 얻는다. (2)의 증명은 생략한다.
\end{proof}

\noindent
다음 보조정리는 Lemma \ref{dga-lemma-make-injective}의 유사물이다.

\begin{lemma}
\label{dga-lemma-factor}
Situation \ref{dga-situation-ABC}에서 $\alpha : x \to y$를
$\text{Comp}(\mathcal{A})$의 사상이라 하자. 그러면
$\text{Comp}(\mathcal{A})$에 분해
$$
\xymatrix{
x \ar[r]^{\tilde{\alpha}}  &
\tilde{y} \ar@<0.5ex>[r]^{\pi} &
y\ar@<0.5ex>[l]^s
}
$$
가 존재하여 다음을 만족한다.
\begin{enumerate}
\item $\tilde{\alpha}$는 허용 단사상이고
$\pi\tilde{\alpha}=\alpha$이다.
\item $\pi s=1_y$이고 $s\pi$가 $1_{\tilde{y}}$와 호모토픽인
$\text{Comp}(\mathcal{A})$의 사상 $s:y\to\tilde{y}$가 존재한다.
\end{enumerate}
\end{lemma}

\begin{proof}
공리 (A)에 의해 $\tilde{y}$를 $y$와 $C(1_x)$의 미분 등급 직합으로
택할 수 있다. 즉 도식
$$
\xymatrix@C=3pc{
y \ar@<0.5ex>[r]^s  &
y\oplus C(1_x) \ar@<0.5ex>[l]^{\pi} \ar@<0.5ex>[r]^{p} &
C(1_x)\ar@<0.5ex>[l]^t
}
$$
이 존재하며 모든 사상은 차수 영이고 $\text{Comp}(\mathcal{A})$에
속한다. $\tilde{y} = y \oplus C(1_x)$로 두자. 그러면
$1_{\tilde{y}} = s\pi + tp$이다. 이제 도식
$$
\xymatrix{
x \ar[r]^{\tilde{\alpha}}  &
\tilde{y} \ar@<0.5ex>[r]^{\pi} &
y\ar@<0.5ex>[l]^s
}
$$
을 생각하자. 여기서 $\tilde{\alpha}$는 사상
$x\xrightarrow{\alpha}y$와 허용 짧은 완전열
$$
\xymatrix{
x \ar@<0.5ex>[r]  &
C(1_x) \ar@<0.5ex>[l] \ar@<0.5ex>[r] &
x[1]\ar@<0.5ex>[l]
}
$$
에 들어가는 자연 사상 $x\to C(1_x)$가 유도한다. 따라서 이 도식의
차수 0 사상 $C(1_x)\to x$와 영사상 $y\to x$는 차수 0 사상
$\beta : \tilde{y} \to x$를 유도한다. $\tilde{\alpha}$는 허용 짧은
완전열
$$
\xymatrix{
x\ar[r]^{\tilde{\alpha}} &
\tilde{y} \ar[r] &
x[1]
}
$$
에 들어가므로 허용 단사상이다. 또한 $\tilde{\alpha}$의 구성에 의해
$\pi\tilde{\alpha} = \alpha$이고, 첫째 도식에 의해 $\pi s = 1_y$이다.
$s\pi$가 $1_{\tilde{y}}$와 호모토픽임을 보이는 일만 남았다.
% P03-STACKS-E553: frozen source writes 1_x although the calculation acts on C(1_x).
어떤 차수 $-1$ 사상에 대하여 $1_x$를 $d(h)$로 쓰자. 그러면 마지막
주장은
\begin{align*}
1_{\tilde{y}} - s\pi
= &
tp \\
= &
t(dh)p\quad\text{(by Lemma \ref{dga-lemma-id-cone-null})} \\
= &
d(thp)
\end{align*}
에서 따른다. 실제로 $dt = dp = 0$이고 $t$의 차수는 영이다.
\end{proof}

\noindent
다음 보조정리는 Lemma \ref{dga-lemma-sequence-maps-split}의 유사물이다.

\begin{lemma}
\label{dga-lemma-analogue-sequence-maps-split}
Situation \ref{dga-situation-ABC}에서
$x_1 \to x_2 \to \ldots \to x_n$을 $\text{Comp}(\mathcal{A})$의
합성 가능한 사상들의 열이라 하자. 그러면 $\text{Comp}(\mathcal{A})$에
가환 도식
$$
\xymatrix{x_1\ar[r] & x_2\ar[r] & \ldots\ar[r] & x_n\\
y_1\ar[r]\ar[u] & y_2\ar[r]\ar[u] & \ldots\ar[r] & y_n\ar[u]}
$$
이 존재하여, 각 $y_i\to y_{i+1}$은 허용 단사상이고 각
$y_i\to x_i$는 호모토피 동치이다.
\end{lemma}

\begin{proof}
$n=1$인 경우는 자명하다. 단순히 $y_1 = x_1$로 두면 $x_1$의 항등
사상은 특히 호모토피 동치이다. $n = 2$인 경우는 Lemma
\ref{dga-lemma-factor}가 준다. $x_{n - 1}$까지 도식을 구성했다고 하자.
합성 $y_{n - 1} \to x_{n-1} \to x_n$에 Lemma \ref{dga-lemma-factor}를
적용하여 $y_n$을 얻는다. 그러면 $y_{n - 1} \to y_n$은 허용 단사상이고
$y_n \to x_n$은 호모토피 동치이다.
\end{proof}

\noindent
다음 보조정리는 Lemma \ref{dga-lemma-nilpotent}의 유사물이다.

\begin{lemma}
\label{dga-lemma-triseq}
Situation \ref{dga-situation-ABC}에서 $x_i \to y_i \to z_i$를
$\mathcal{A}$의 사상들이라 하자($i=1,2,3$). 이때
$x_2 \to y_2\to z_2$는 허용 짧은 완전열이라고 하자.
$b : y_1 \to y_2$와 $b' : y_2\to y_3$를 다음 두 도식이 호모토피까지
가환하도록 하는 $\text{Comp}(\mathcal{A})$의 사상들이라 하자.
$$
\vcenter{
\xymatrix{
x_1 \ar[d]_0 \ar[r] &
y_1 \ar[r] \ar[d]_b &
z_1 \ar[d]_0 \\
x_2 \ar[r] & y_2 \ar[r] & z_2
}
}
\quad\text{and}\quad
\vcenter{
\xymatrix{
x_2 \ar[d]^0 \ar[r] &
y_2 \ar[r] \ar[d]^{b'} &
z_2 \ar[d]^0 \\
x_3 \ar[r] & y_3 \ar[r] & z_3
}
}
$$
그러면 $b'\circ b$는 $0$과 호모토픽이다.
\end{lemma}

\begin{proof}
Lemma \ref{dga-lemma-homo-change}에 의해 $b$와 $b'$을 각각 호모토픽인
사상 $\tilde{b}$와 $\tilde{b}'$로 바꾸어, 왼쪽 도식의 오른쪽 사각형과
오른쪽 도식의 왼쪽 사각형이 가환하게 할 수 있다. $\mathcal{A}$의 차수
$-1$ 사상 $h$와 $h'$에 대하여
$b = \tilde{b} + d(h)$ 및 $b'=\tilde{b}'+d(h')$라고 쓰자. 그러면
$$
b'b = \tilde{b}'\tilde{b} + d(\tilde{b}'h + h'\tilde{b} + h'd(h))
$$
이다. 실제로 $d(\tilde{b})=d(\tilde{b}')=0$이다. 즉 $b'b$는
$\tilde{b}'\tilde{b}$와 호모토픽이다. 이제
$\tilde{b}'\tilde{b}=0$임을 보이자. $x_2\xrightarrow{f} y_2\xrightarrow{g} z_2$
가 허용 짧은 완전열이므로 $\text{id}_{y_2} = f\pi + sg$를 만족하는
차수 $0$ 사상 $\pi : y_2 \to x_2$와 $s : z_2 \to y_2$가 존재한다.
따라서
$$
\tilde{b}'\tilde{b} = \tilde{b}'(f\pi + sg)\tilde{b} = 0
$$
이다. 두 가환 사각형으로부터 $g\tilde{b} = 0$ 및
$\tilde{b}'f = 0$이 따르기 때문이다.
\end{proof}

\noindent
다음 보조정리는 Lemma \ref{dga-lemma-triangle-independent-splittings}의
유사물이다.

\begin{lemma}
\label{dga-lemma-analogue-triangle-independent-splittings}
Situation \ref{dga-situation-ABC}에서
$0 \to x \to y \to z \to 0$을 $\text{Comp}(\mathcal{A})$의 허용 짧은
완전열이라 하자. Lemma \ref{dga-lemma-get-triangle}에서 정의한
$\delta : z \to x[1]$를 갖는 삼각형
$$
\xymatrix{x\ar[r] & y\ar[r] & z\ar[r]^{\delta} & x[1]}
$$
은 $K(\mathcal{A})$에서 표준 동형까지 Lemma
\ref{dga-lemma-get-triangle}에서 한 선택들과 무관하다.
\end{lemma}

\begin{proof}
$\delta$가 분할
$$
\xymatrix{
x \ar@<0.5ex>[r]^{a} &
y \ar@<0.5ex>[r]^b\ar@<0.5ex>[l]^{\pi} &
z \ar@<0.5ex>[l]^s
}
$$
로 정의되고 $\delta'$가 $\pi,s$ 대신 $\pi',s'$을 사용한 분할로
정의된다고 하자. 그러면
$$
s'-s = (a\pi + sb)(s'-s) = a\pi s'
$$
이다. 실제로 $bs' = bs = 1_z$이고 $\pi s = 0$이다. 마찬가지로
$$
\pi' - \pi = (\pi' - \pi)(a\pi + sb) = \pi'sb
$$
이다. Lemma \ref{dga-lemma-get-triangle}에서 구성한 바와 같이
$\delta = \pi d(s)$이고 $\delta' = \pi'd(s')$이므로 다음과 같이
계산할 수 있다.
$$
\delta' = \pi'd(s') = (\pi + \pi'sb)d(s + a\pi s') = \delta + d(\pi s')
$$
여기서 $\pi a = 1_x$, $ba = 0$, 그리고 Lemma
\ref{dga-lemma-get-triangle}의 공식 (5)에 의한
$\pi'sbd(s') = \pi'sba\pi d(s') = 0$을 사용하였다.
\end{proof}

\noindent
다음 보조정리는 Lemma \ref{dga-lemma-rotate-cone}의 유사물이다.

\begin{lemma}
\label{dga-lemma-restate-axiom-c}
Situation \ref{dga-situation-ABC}에서 $f: x \to y$를
$\text{Comp}(\mathcal{A})$의 사상이라 하자. 삼각형
$(y, c(f), x[1], i, p, f[1])$은 허용 짧은 완전열
$$
\xymatrix{y\ar[r] & c(f) \ar[r] & x[1]}
$$
에 결부된 삼각형이며, 원뿔 $c(f)$는 Lemma
\ref{dga-lemma-get-triangle}에서와 같이 정의한다.
\end{lemma}

\begin{proof}
이는 공리 (C)에서 따른다.
\end{proof}

\noindent
다음 보조정리는 Lemma \ref{dga-lemma-rotate-triangle}의 유사물이다.

\begin{lemma}
\label{dga-lemma-cone-rotate-isom}
Situation \ref{dga-situation-ABC}에서 $\alpha : x \to y$와
$\beta : y \to z$가 $\text{Comp}(\mathcal{A})$의 허용 짧은 완전열
$$
\xymatrix{
x \ar[r] &
y\ar[r] &
z
}
$$
을 정의한다고 하자. $(x, y, z, \alpha, \beta, \delta)$를
$K(\mathcal{A})$에서 이에 결부된 삼각형이라 하자. 그러면 두 삼각형
$$
(z[-1], x, y, \delta[-1], \alpha, \beta)
\quad\text{and}\quad
(z[-1], x, c(\delta[-1]), \delta[-1], i, p)
$$
은 동형이다.
\end{lemma}

\begin{proof}
다음 꼴의 도식이 있다.
$$
\xymatrix{
z[-1]\ar[r]^{\delta[-1]}\ar[d]^1 &
x\ar@<0.5ex>[r]^{\alpha}\ar[d]^1 &
y\ar@<0.5ex>[r]^{\beta}\ar@{.>}[d]\ar@<0.5ex>[l]^{\tilde{\alpha}} &
z\ar[d]^1\ar@<0.5ex>[l]^{\tilde\beta} \\
z[-1] \ar[r]^{\delta[-1]} &
x\ar@<0.5ex>[r]^i &
c(\delta[-1]) \ar@<0.5ex>[r]^p\ar@<0.5ex>[l]^{\tilde i} &
z\ar@<0.5ex>[l]^{\tilde p}
}
$$
여기서 $\alpha, \beta, i$와 $p$의 분할은 각각
$\tilde{\alpha}, \tilde{\beta}, \tilde{i},$와 $\tilde{p}$로 주어진다.
$y \to c(\delta[-1])$를 $i\tilde{\alpha} + \tilde{p}\beta$로,
$c(\delta[-1]) \to y$를 $\alpha \tilde{i} + \tilde{\beta} p$로
정의하자. 먼저 이것들이 $\text{Comp}(\mathcal{A})$의 사상을 정의함을
확인하자. Lemma \ref{dga-lemma-get-triangle}의 등식들에 의해 관계식
$\delta[-1] = \tilde{\alpha}d(\tilde{\beta}) = -d(\tilde{\alpha})\tilde{\beta}$
및 $\delta[-1] = \tilde{i}d(\tilde{p})$가 성립한다. 그러면
\begin{align*}
d(\tilde{\alpha})
& =
d(\tilde{\alpha})\tilde{\beta}\beta \\
& =
-\delta[-1]\beta
\end{align*}
이다. 첫째 등식에는 Lemma \ref{dga-lemma-get-triangle}의 식 (6)을,
둘째 등식에는 앞의 관찰을 사용하였다. 마찬가지로
$d(\tilde{p}) = i\delta[-1]$을 얻는다. 따라서
\begin{align*}
d(i\tilde{\alpha} + \tilde{p}\beta)
& =
d(i)\tilde{\alpha} + id(\tilde{\alpha}) +
d(\tilde{p})\beta + \tilde{p}d(\beta) \\
& =
id(\tilde{\alpha}) + d(\tilde{p})\beta \\
& =
-i\delta[-1]\beta + i\delta[-1]\beta \\
& =
0
\end{align*}
이므로 $i\tilde{\alpha} + \tilde{p}\beta$는 실제로
$\text{Comp}(\mathcal{A})$의 사상이다. 비슷한 계산에 의해
$\alpha \tilde{i} + \tilde{\beta} p$도 $\text{Comp}(\mathcal{A})$의
사상이다. 이 사상들이 가환 도식에 들어감은 자명하다. 다음을 계산한다.
\begin{align*}
(i\tilde{\alpha} + \tilde{p}\beta)(\alpha \tilde{i} + \tilde{\beta} p)
& =
i\tilde{\alpha}\alpha\tilde{i} + i\tilde{\alpha}\tilde{\beta}p
+ \tilde{p}\beta\alpha\tilde{i} + \tilde{p}\beta\tilde{\beta}p \\
& =
i\tilde{i} + \tilde{p}p \\
& =
1_{c(\delta[-1])}
\end{align*}
여기서 Lemma \ref{dga-lemma-get-triangle}의 등식들을 자유롭게 사용하였다.
마찬가지로
$(\alpha \tilde{i} + \tilde{\beta} p)(i\tilde{\alpha} + \tilde{p}\beta) = 1_y$,
를 계산하므로 $y \cong c(\delta[-1])$이다. 따라서 문제의 두 삼각형은
동형이다.
\end{proof}

\noindent
다음 보조정리는 Lemma \ref{dga-lemma-third-isomorphism}의 유사물이다.

\begin{lemma}
\label{dga-lemma-analogue-third-isomorphism}
Situation \ref{dga-situation-ABC}에서 $f_1 : x_1 \to y_1$과
$f_2 : x_2 \to y_2$를 $\text{Comp}(\mathcal{A})$의 사상들이라 하자.
% P03-STACKS-E924: frozen target triangle retains i_1 and p_1.
$$
(a,b,c): (x_1,y_1,c(f_1), f_1, i_1, p_1) \to (x_2,y_2, c(f_2), f_2, i_1, p_1)
$$
를 $K(\mathcal{A})$에서 삼각형들의 임의의 사상이라 하자. $a$와 $b$가
호모토피 동치이면 $c$도 호모토피 동치이다.
\end{lemma}

\begin{proof}
$a$와 $b$는 호모토피 동치이므로 $K(\mathcal{A})$에서 가역이다.
$a^{-1}$와 $b^{-1}$를 $K(\mathcal{A})$에서의 역이라 하면 가환 도식
$$
\xymatrix{
x_2\ar[d]^{a^{-1}}\ar[r]^{f_2} &
y_2\ar[d]^{b^{-1}}\ar[r]^{i_2} &
c(f_2)\ar[d]^{c'} \\
x_1\ar[r]^{f_1} &
y_1 \ar[r]^{i_1} &
c(f_1)
}
$$
을 얻는다. 여기서 사상 $c'$은 위 도식의 왼쪽 가환 사각형에 Lemma
\ref{dga-lemma-cone}을 적용하여 정의한다. 도식은 $K(\mathcal{A})$에서
가환하므로 Lemma \ref{dga-lemma-triseq}에 의해 다음을 보이면 충분하다.
$K(\mathcal{A})$에서 삼각형들의 사상
$(1,1,c): (x,y,c(f),f,i,p)\to (x,y,c(f),f,i,p)$
가 주어지면 사상 $c$는 $K(\mathcal{A})$의 동형이다. $K(\mathcal{A})$에
가환 도식들이 있다.
$$
\vcenter{
\xymatrix{
y\ar[d]^{1}\ar[r] &
c(f)\ar[d]^{c}\ar[r] &
x[1]\ar[d]^{1} \\
y\ar[r] &
c(f) \ar[r] &
x[1]
}
}
\quad\Rightarrow\quad
\vcenter{
\xymatrix{
y\ar[d]^{0}\ar[r] &
c(f)\ar[d]^{c-1}\ar[r] &
x[1]\ar[d]^{0} \\
y\ar[r] &
c(f) \ar[r] &
x[1]
}
}
$$
각 행은 허용 짧은 완전열이므로 Lemma \ref{dga-lemma-triseq}에 의해
등식 $(c-1)^2 = 0$을 얻는다. 이로부터 $2-c$가 $K(\mathcal{A})$에서
$c$의 역임을 알 수 있으므로 $c$는 동형이다.
\end{proof}

\noindent
다음 보조정리는 Lemma \ref{dga-lemma-the-same-up-to-isomorphisms}의
유사물이다.

\begin{lemma}
\label{dga-lemma-cone-homotopy}
Situation \ref{dga-situation-ABC}에서 다음이 성립한다.
\begin{enumerate}
\item 허용 짧은 완전열
$x\xrightarrow{\alpha} y\xrightarrow{\beta} z$가 주어졌다고 하자.
그러면 도식
\begin{equation}
\label{dga-equation-cone-isom-triangle}
\vcenter{
\xymatrix{
x\ar[r]^{\alpha}\ar[d] &
y\ar[r]^{b}\ar[d] &
C(\alpha)\ar[r]^{-c}\ar@{.>}[d]^{e} &
x[1]\ar[d] \\
x\ar[r]^{\alpha} &
y\ar[r]^{\beta} &
z\ar[r]^{\delta} & x[1]
}
}
\end{equation}
이 $K(\mathcal{A})$에서 삼각형들의 동형을 정의하도록 하는 호모토피
동치 $e:C(\alpha)\to z$가 존재한다. 여기서
$y\xrightarrow{b}C(\alpha)\xrightarrow{c}x[1]$은 공리 (C)에서와 같이
주어진 허용 짧은 완전열이다.
\item $\alpha : x \to y$를 $\text{Comp}(\mathcal{A})$의 사상이라 하고,
$x \xrightarrow{\tilde{\alpha}} \tilde{y} \to y$를 Lemma
\ref{dga-lemma-factor}에서와 같이 주어진 분해라 하자. 여기서 허용 단사상
% P03-STACKS-E925: frozen source gives the codomain y rather than tilde y.
$x \xrightarrow{\tilde{\alpha}} y$는 허용 짧은 완전열
$$
\xymatrix{
x \ar[r]^{\tilde{\alpha}} &
\tilde{y} \ar[r] & z
}
$$
로 연장된다. 그러면 삼각형들의 동형
$$
\xymatrix{
x \ar[r]^{\tilde{\alpha}} \ar[d] &
\tilde{y} \ar[r] \ar[d] &
z \ar[r]^{\delta} \ar@{.>}[d]^{e} &
x[1] \ar[d] \\
x \ar[r]^{\alpha} &
y \ar[r] &
C(\alpha) \ar[r]^{-c} &
x[1]
}
$$
이 존재한다. 여기서 위쪽 삼각형은 열
$x \xrightarrow{\tilde{\alpha}} \tilde{y} \to z$에 결부된 삼각형이다.
\end{enumerate}
\end{lemma}

\begin{proof}
(1)에 대하여 $c$의 부호를 바꾸지 \emph{않은} 다음의 더 완전한 도식을
생각하자.
$$
\xymatrix{
x\ar@<0.5ex>[r]^{\alpha} \ar[d] &
y\ar@<0.5ex>[l]^{\pi} \ar@<0.5ex>[r]^{b}\ar[d] &
C(\alpha)\ar@<0.5ex>[l]^{p} \ar@<0.5ex>[r]^{c}\ar@{.>}@<0.5ex>[d]^{e} &
x[1]\ar@<0.5ex>[l]^{\sigma} \ar[d]\ar@<0.5ex>[r]^{\alpha} &
y[1]\ar@<0.5ex>[l]^{\pi} \\
x\ar@<0.5ex>[r]^{\alpha} &
y\ar@<0.5ex>[r]^{\beta} \ar@<0.5ex>[l]^{\pi} &
z\ar[r]^{\delta}\ar@<0.5ex>[l]^{s} \ar@{.>}@<0.5ex>[u]^{f} &
x[1]
}
$$
여기서 허용 짧은 완전열 $x\xrightarrow{\alpha} y\xrightarrow{\beta} z$에는
분할 $\pi$, $s$가 주어져 있고, 허용 짧은 완전열
$y\xrightarrow{b}C(\alpha)\xrightarrow{c}x[1]$에는 분할 $p$, $\sigma$가
주어져 있다. 시프트 아래에서 Hom 집합들을 동일시하면 Lemma
\ref{dga-lemma-get-triangle}의 구성에 의해
$$
\alpha = pd(\sigma) = -d(p)\sigma,\quad
\delta = \pi d(s) = -d(\pi)s
$$
임에 유의하라.

\medskip\noindent
$e=\beta p$ 및 $f=bs-\sigma\delta$로 정의한다. 먼저 이들이
$\text{Comp}(\mathcal{A})$의 사상임을 확인하자. $d(e)=\beta d(p)$가
영임을 보이려면 $\beta d(p)b$와 $\beta d(p)\sigma$가 모두 영임을
보이면 충분하다. 그런데
$$
\beta d(p)b = \beta d(pb) = \beta d(1_y) = 0,\quad
\beta d(p)\sigma = -\beta\alpha = 0
$$
이다. 마찬가지로 $d(f)=bd(s)-d(\sigma)\delta$가 영임을 확인하려면
$p$ 및 $c$와의 뒤쪽 합성이 모두 영임을 확인하면 충분하다. 그런데
\begin{align*}
pbd(s) - pd(\sigma)\delta
= &
d(s)-\alpha\delta = d(s)-\alpha\pi d(s) = 0 \\
cbd(s)-cd(\sigma)\delta
= &
-cd(\sigma)\delta = -d(c\sigma)\delta = 0
\end{align*}
이다. 도식 \ref{dga-equation-cone-isom-triangle}의 왼쪽 두 사각형의 가환성은
정의에서 곧바로 따른다. 오른쪽 사각형이 호모토피까지 가환함을
증명하기 전에 먼저 $e$가 호모토피 동치임을 확인하자. 명백히
$$
ef=\beta p (bs-\sigma\delta)=\beta s=1_z
$$
이다. $fe$가 $1_{C(\alpha)}$와 호모토픽임을 확인하기 위해 먼저
다음을 관찰한다.
% P03-STACKS-E926: frozen first identity retains d(alpha).
$$
b\alpha = bpd(\alpha) = d(\sigma),\quad
\alpha c = -d(p)\sigma c = -d(p),\quad
d(\pi)p = d(\pi)s\beta p = -\delta\beta p
$$
이 등식들을 사용하면
\begin{align*}
1_{C(\alpha)} = &
bp + \sigma c
\quad (\text{from }y \xrightarrow{b} C(\alpha) \xrightarrow{c} x[1]) \\
= &
b(\alpha\pi + s\beta)p + \sigma(\pi\alpha)c
\quad (\text{from }x \xrightarrow{\alpha} y \xrightarrow{\beta} z) \\
= &
d(\sigma)\pi p + bs\beta p - \sigma\pi d(p)
\quad (\text{by the first two identities above}) \\
= &
d(\sigma)\pi p + bs\beta p - \sigma\delta\beta p
+ \sigma\delta\beta p - \sigma\pi d(p) \\
= &
(bs - \sigma\delta)\beta p + d(\sigma)\pi p
- \sigma d(\pi)p - \sigma\pi d(p)\quad
(\text{by the third identity above}) \\
= &
fe + d(\sigma \pi p)
\end{align*}
을 계산한다. 여기서 $\sigma \in \Hom^{-1}(x, C(\alpha))$이다
(Lemma \ref{dga-lemma-id-cone-null}의 증명을 비교하라). 따라서 $e$와 $f$는
서로 호모토피 역이다. 마지막으로 도식
\ref{dga-equation-cone-isom-triangle}의 오른쪽 사각형이 호모토피까지
가환함을 확인하려면 $-cf=\delta$임을 확인하면 충분하다. 이는
$$
-cf = -c(bs-\sigma\delta) = c\sigma\delta = \delta
$$
에서 따른다. 실제로 $cb=0$이다.

\medskip\noindent
(2)에 대하여 Lemma \ref{dga-lemma-factor}에서와 같이 주어진 분해
$x\xrightarrow{\tilde{\alpha}}\tilde{y}\to y$를 생각하자. 따라서 둘째
사상은 호모토피 동치이다. Lemmas \ref{dga-lemma-cone} 및
\ref{dga-lemma-analogue-third-isomorphism}에 의해 두 삼각형
$$
x \xrightarrow{\alpha} y \to C(\alpha) \to x[1]
\quad\text{and}\quad
x \xrightarrow{\tilde{\alpha}} \tilde{y} \to C(\tilde{\alpha}) \to x[1]
$$
사이에 동형이 존재한다. 삼각형들의 동형을 합성할 수 있으므로
$\alpha$를 $\tilde{\alpha}$로, $y$를 $\tilde{y}$로, $C(\alpha)$를
$C(\tilde{\alpha})$로 바꾸어 $\alpha$가 허용 단사상이라고 가정해도
된다. 이 경우 결과는 (1)에서 따른다.
\end{proof}

\noindent
다음 보조정리는 Lemma \ref{dga-lemma-homotopy-category-pre-triangulated}의
유사물이다.

\begin{lemma}
\label{dga-lemma-analogue-homotopy-category-pre-triangulated}
Situation \ref{dga-situation-ABC}에서 호모토피 범주 $K(\mathcal{A})$는 그
자연 이동 함자들과 특수 삼각형들을 갖춘 전삼각범주이다.
\end{lemma}

\begin{proof}
TR1, TR2, TR3을 각각 확인하겠다.

\medskip\noindent
TR1의 증명. 정의에 의해 특수 삼각형과 동형인 모든 삼각형은 특수
삼각형이다.
$$
\xymatrix{x\ar[r]^{1_x} & x\ar[r] & 0}
$$
가 허용 짧은 완전열이므로 $(x, x, 0, 1_x, 0, 0)$은 특수 삼각형이다.
또한 $\text{Comp}(\mathcal{A})$의 사상 $\alpha : x \to y$가 주어지면
$(x, y, c(\alpha), \alpha, i, -p)$가 주는 삼각형은 Lemma
\ref{dga-lemma-cone-homotopy}에 의해 특수 삼각형이다.

\medskip\noindent
TR2의 증명. $(x,y,z,\alpha,\beta,\gamma)$를 삼각형이라 하고
$(y,z,x[1],\beta,\gamma,-\alpha[1])$이 특수 삼각형이라고 하자. 그러면
허용 짧은 완전열 $0 \to x' \to y' \to z' \to 0$이 존재하여 이에
결부된 삼각형 $(x',y',z',\alpha',\beta',\gamma')$이
$(y,z,x[1],\beta,\gamma,-\alpha[1])$과 동형이다. 회전하면
$(x,y,z,\alpha,\beta,\gamma)$가
$(z'[-1],x',y', \gamma'[-1], \alpha',\beta')$와 동형임을 알 수 있다.
Lemma \ref{dga-lemma-cone-rotate-isom}에 의해
$(z'[-1],x',y', \gamma'[-1], \alpha',\beta')$는
$(z'[-1],x',c(\gamma'[-1]), \gamma'[-1], i, p)$와 동형이다. 다음
도식에 나타낸 부호 변경과 함께 두 동형을 합성한다.
% P03-STACKS-E557: frozen diagram retains unprimed x in both lower rows.
$$
\xymatrix@C=3pc{
x\ar[r]^{\alpha}\ar[d] &
y\ar[r]^{\beta}\ar[d] &
z\ar[r]^{\gamma}\ar[d] &
x[1]\ar[d] \\
z'[-1]\ar[r]^{-\gamma'[-1]}\ar[d]_{-1_{z'[-1]}} &
x \ar[r]^{\alpha'}\ar@{=}[d] &
y' \ar[r]^{\beta'} \ar[d] &
z'\ar[d]^{-1_{z'}} \\
z'[-1]\ar[r]^{\gamma'[-1]} &
x \ar[r]^{\alpha'} &
c(\gamma'[-1]) \ar[r]^{-p} &
z'
}
$$
Lemma \ref{dga-lemma-cone-homotopy} (2)에 의해
$(x,y,z,\alpha,\beta,\gamma)$가 특수 삼각형임을 얻는다. 역으로
$(x,y,z,\alpha,\beta,\gamma)$가 특수 삼각형이라고 하자. 그러면 Lemma
\ref{dga-lemma-cone-homotopy} (1)에 의해 이는 어떤
$\text{Comp}(\mathcal{A})$의 사상 $\alpha': x' \to y'$에 대하여
$(x',y', c(\alpha'), \alpha', i, -p)$ 꼴의 삼각형과 동형이다.
% P03-STACKS-E927: frozen model triangles retain alpha[1], not alpha'[1].
회전된 삼각형 $(y,z,x[1],\beta,\gamma, -\alpha[1])$은 삼각형
$(y',c(\alpha'), x'[1], i, -p, -\alpha[1])$과 동형이고, 이는 다시
$(y',c(\alpha'), x'[1], i, p, \alpha[1])$과 동형이다. Lemma
\ref{dga-lemma-restate-axiom-c}에 의해 이 삼각형은 특수 삼각형이므로
$(y,z,x[1], \beta,\gamma, -\alpha[1])$도 특수 삼각형이다.

\medskip\noindent
TR3의 증명. $(x,y,z, \alpha,\beta,\gamma)$와
$(x',y',z',\alpha',\beta',\gamma')$가
% P03-STACKS-E928: frozen source places distinguished triangles in Comp(A).
$\text{Comp}(\mathcal{A})$의 특수 삼각형들이라고 하자. 또한
$f: x \to x'$와 $g: y \to y'$가
$\alpha' \circ f = g \circ \alpha$를 만족하는 사상들이라 하자. Lemma
\ref{dga-lemma-cone-homotopy}에 의해
$(x,y,z,\alpha,\beta,\gamma)= (x,y,c(\alpha),\alpha, i, -p)$
및 $(x',y',z', \alpha',\beta',\gamma')= (x',y',c(\alpha'), \alpha',i',-p')$
이라고 가정해도 된다. 이제 Lemma \ref{dga-lemma-cone}을 적용하면 증명이
끝난다.
\end{proof}

\noindent
다음 보조정리는 Lemma \ref{dga-lemma-two-split-injections}의 유사물이다.

\begin{lemma}
\label{dga-lemma-dgc-analogue-tr4}
Situation \ref{dga-situation-ABC}에서 $\mathcal{A}$의 허용 단사상
% P03-STACKS-E933: frozen source uses A although admissibility is defined in Comp(A).
$x \xrightarrow{\alpha} y$, $y \xrightarrow{\beta} z$가 주어졌다고 하자.
그러면 TR4가 성립하는 특수 삼각형
$(x,y,q_1,\alpha,p_1,\delta_1)$,
$(x,z,q_2,\beta\alpha,p_2,\delta_2)$ 및
$(y,z,q_3,\beta,p_3,\delta_3)$가 존재한다.
\end{lemma}

\begin{proof}
허용 단사상 $x\xrightarrow{\alpha} y$와 $y\xrightarrow{\beta}z$가
주어지면, 이들을 허용 짧은 완전열로 연장하여 다음 특수 삼각형들을
찾을 수 있다.
$$
\xymatrix{
x\ar@<0.5ex>[r]^{\alpha} &
y\ar@<0.5ex>[l]^{\pi_1} \ar@<0.5ex>[r]^{p_1} &
q_1 \ar[r]^{\delta_1} \ar@<0.5ex>[l]^{s_1} &
x[1]
}
$$
$$
\xymatrix{
x\ar@<0.5ex>[r]^{\beta\alpha} &
z\ar@<0.5ex>[l]^{\pi_1\pi_3} \ar@<0.5ex>[r]^{p_2} &
q_2 \ar[r]^{\delta_2} \ar@<0.5ex>[l]^{s_2} &
x[1]
}
$$
% P03-STACKS-E929: frozen third triangle ends at x[1], not y[1].
$$
\xymatrix{
y\ar@<0.5ex>[r]^{\beta} &
z\ar@<0.5ex>[l]^{\pi_3} \ar@<0.5ex>[r]^{p_3} &
q_3 \ar[r]^{\delta_3} \ar@<0.5ex>[l]^{s_3} &
x[1]
}
$$
이 도식들에서 사상 $\delta_i$는 Lemma \ref{dga-lemma-get-triangle}에서
정의한 사상들과 마찬가지로
% P03-STACKS-E930: frozen uniform formula invokes an unintroduced pi_2.
$\delta_i = \pi_i d(s_i)$로 정의한다. 이들은 다음 실선 가환 도식에
들어간다.
$$
\xymatrix@C=5pc@R=3pc{
x\ar@<0.5ex>[r]^{\alpha} \ar@<0.5ex>[dr]^{\beta\alpha} &
y\ar@<0.5ex>[d]^{\beta} \ar@<0.5ex>[l]^{\pi_1} \ar@<0.5ex>[r]^{p_1} &
q_1 \ar[r]^{\delta_1} \ar@<0.5ex>[l]^{s_1} \ar@{.>}[dd]^{p_2\beta s_1} &
x[1] \\
 &
z \ar@<0.5ex>[u]^{\pi_3}\ar@<0.5ex>[d]^{p_3}
\ar@<0.5ex>[dr]^{p_2} \ar@<0.5ex>[ul]^{\pi_1\pi_3} & & \\
 &
q_3\ar@<0.5ex>[u]^{s_3} \ar[d]^{\delta_3} &
q_2 \ar@{.>}[l]^{p_3s_2} \ar@<0.5ex>[ul]^{s_2} \ar[dr]^{\delta_2} \\
 &
y[1] & & x[1]}
$$
여기서 점선 화살표들은 표시된 대로 정의하였다. 분명히 이들의 합성은
$p_3s_2p_2\beta s_1 = 0$이다.
% P03-STACKS-E932: frozen source justifies this using the false identity s_2p_2=0.
실제로 $s_2p_2 = 0$이다. 두 점선 화살표가 모두
$\text{Comp}(\mathcal{A})$의 사상임을 주장한다. Lemma
\ref{dga-lemma-get-triangle}의 등식들을 사용하여 이를 확인할 수 있다.
$$
d(p_2\beta s_1) = p_2\beta d(s_1) = p_2\beta\alpha\pi_1 d(s_1) = 0
$$
이는 $p_2\beta\alpha = 0$이기 때문이고,
$$
d(p_3s_2) = p_3d(s_2) = p_3\beta\alpha\pi_1\pi_3 d(s_2) = 0
$$
이는 $p_3\beta = 0$이기 때문이다. $q_1\to q_2\to q_3$가 허용 짧은
완전열임을 확인하려면, 바탕 등급 범주에서
$q_2 = q_1\oplus q_3$이고 위의 두 사상이 각각 여사영과 사영임을
보이는 일만 남았다. 이를 위해 바탕 등급 범주 $\mathcal{C}$에서
$$
y = x\oplus q_1,\quad
z = y\oplus q_3 = x\oplus q_1\oplus q_3
$$
가 성립함을 관찰하자. 여기서 $\pi_1\pi_3$는 첫째 인자로의 사영 사상
% P03-STACKS-E931: frozen prose retains the reversed projection arrow.
$x\oplus q_1\oplus q_3\to z$를 준다. $\mathcal{A}$에 대한 공리 (A)에
의해 $\mathcal{C}$는 가법 범주이므로 Homology, Lemma
\ref{homology-lemma-additive-cat-biproduct-kernel}을 적용하여
$$
\Ker(\pi_1\pi_3) = q_1\oplus q_3
$$
를 $\mathcal{C}$에서 얻는다. $z = x\oplus q_2$에 Homology, Lemma
\ref{homology-lemma-additive-cat-biproduct-kernel}을 다시 적용하면
$\Ker(\pi_1\pi_3) = q_2$를 얻는다. 따라서 $\mathcal{C}$에서
$q_2\cong q_1\oplus q_3$이다. 위에서 정의한 점선 사상들이 여사영과
사영을 준다는 것은 명백하다.

\medskip\noindent
마지막으로 허용 짧은 완전열 $q_1\to q_2\to q_3$가 유도하는 사상
$\delta : q_3 \to q_1[1]$이 $p_1\delta_3$와 일치함을 확인해야 한다.
Lemma \ref{dga-lemma-get-triangle}의 구성에 의해 사상 $\delta$는
\begin{align*}
p_1\pi_3s_2d(p_2s_3)
= &
p_1\pi_3s_2p_2d(s_3) \\
= &
p_1\pi_3(1-\beta\alpha\pi_1\pi_3)d(s_3) \\
= &
p_1\pi_3d(s_3)\quad (\text{since }\pi_3\beta = 0) \\
= &
p_1\delta_3
\end{align*}
로 주어진다. 이것이 원하는 등식이므로 증명이 끝났다.
\end{proof}

\noindent
모든 결과를 종합하면 마침내 Proposition
\ref{dga-proposition-homotopy-category-triangulated}의 유사물을 얻는다.

\begin{proposition}
\label{dga-proposition-ABC-homotopy-category-triangulated}
Situation \ref{dga-situation-ABC}에서 호모토피 범주 $K(\mathcal{A})$는 그
자연 이동 함자들과 특수 삼각형들을 갖춘 삼각범주이다.
\end{proposition}

\begin{proof}
Lemma \ref{dga-lemma-analogue-homotopy-category-pre-triangulated}에 의해
$K(\mathcal{A})$는 전삼각범주이다. Lemmas
\ref{dga-lemma-analogue-sequence-maps-split} 및
\ref{dga-lemma-dgc-analogue-tr4}를 Derived Categories, Lemma
\ref{derived-lemma-easier-axiom-four}와 결합하면 $K(\mathcal{A})$가
삼각범주라는 결론을 얻는다.
\end{proof}

\begin{lemma}
\label{dga-lemma-functor-between-ABC}
$R$을 환이라 하자. $F : \mathcal{A} \to \mathcal{B}$를 공리 (A), (B),
(C)를 만족하는 $R$ 위의 미분 등급 범주들 사이의 함자라 하고
$F(x[1]) = F(x)[1]$이라고 하자. 그러면 $F$는 삼각범주들의 완전 함자
$K(\mathcal{A}) \to K(\mathcal{B})$를 유도한다.
\end{lemma}

\begin{proof}
실제로 $x \to y \to z$가 $\text{Comp}(\mathcal{A})$의 허용 짧은
완전열이면 $F(x) \to F(y) \to F(z)$는
$\text{Comp}(\mathcal{B})$의 허용 짧은 완전열이다. 또한 Lemma
\ref{dga-lemma-get-triangle}에서 구성한 ``경계'' 사상
$\delta = \pi\text{d}(s) : z \to x[1]$은 사상
$F(\delta) : F(z) \to F(x[1]) = F(x)[1]$을 주며, 이는 허용 짧은
완전열 $F(x) \to F(y) \to F(z)$의 경계 사상
$F(\pi) \text{d}(F(s))$와 같다.
\end{proof}





\section{쌍가군}
\label{dga-section-bimodules}

\noindent
Section \ref{dga-section-tensor-product}에서 시작한 논의를 계속한다.

\begin{definition}
\label{dga-definition-bimodule}
쌍가군을 정의한다. $R$을 환이라 하자.
\begin{enumerate}
\item $A$와 $B$를 $R$-대수라 하자. {\it $(A, B)$-쌍가군}이란
다음 $R$-쌍선형 사상들을 갖춘 $R$-가군 $M$로서
% P03-STACKS-E934: frozen English misspells “equipped” as “equippend”.
$$
A \times M \to M, (a, x) \mapsto ax
\quad\text{및}\quad
M \times B \to M, (x, b) \mapsto xb
$$
다음을 만족하는 것을 말한다.
\begin{enumerate}
\item $a'(ax) = (a'a)x$ 및 $(xb)b' = x(bb')$,
\item $a(xb) = (ax)b$,
\item $1 x = x = x 1$.
\end{enumerate}
\item $A$와 $B$를 $\mathbf{Z}$-등급 $R$-대수라 하자.
{\it 등급 $(A, B)$-쌍가군}이란 등급
$M = \bigoplus M^n$을 가지며
$A^n M^m \subset M^{n + m}$ 및 $M^n B^m \subset M^{n + m}$을
만족하는 $(A, B)$-쌍가군 $M$을 말한다.
\item $A$와 $B$를 미분 등급 $R$-대수라 하자.
{\it 미분 등급 $(A, B)$-쌍가군}이란 차수 $1$인 동차 미분
$\text{d} : M \to M$을 갖추고, 동차 원소 $a \in A$, $x \in M$,
$b \in B$에 대하여
$\text{d}(ax) = \text{d}(a)x + (-1)^{\deg(a)}a\text{d}(x)$ 및
$\text{d}(xb) = \text{d}(x)b + (-1)^{\deg(x)}x\text{d}(b)$을
만족하는 등급 $(A, B)$-쌍가군을 말한다.
\end{enumerate}
\end{definition}

\noindent
미분 등급 $(A, B)$-쌍가군 $M$은 등급과 미분이 일치하고
$A$-가군 구조와 $B$-가군 구조가 서로 가환하는, 왼쪽 미분 등급
$A$-가군이기도 한 오른쪽 미분 등급 $B$-가군과 같은 것임을 주목하자.
정확한 명제는 다음과 같다.

\begin{lemma}
\label{dga-lemma-what-makes-a-bimodule-dg}
$R$을 환이라 하자. $(A, \text{d})$와 $(B, \text{d})$를
$R$ 위의 미분 등급 대수라 하자. $M$을 오른쪽 미분 등급 $B$-가군이라 하자.
주어진 미분 등급 $B$-가군 구조와 양립하는 $M$ 위의 $(A, B)$-쌍가군
구조와 미분 등급 $R$-대수의 준동형
$$
A
\longrightarrow
\Hom_{\text{Mod}^{dg}_{(B, \text{d})}}(M, M)
$$
사이에는 일대일($1$ 대 $1$) 대응이 있다.
\end{lemma}

\begin{proof}
$\mu : A \times M \to M$가 $M$의 바탕 $R$-가군 복합체
$M^\bullet$ 위에 왼쪽 미분 등급 $A$-가군 구조를 정의한다고 하자.
Lemma \ref{dga-lemma-left-module-structure}에 의하면 구조 $\mu$는 미분 등급
$R$-대수의 사상 $\gamma : A \to \Hom^\bullet(M^\bullet, M^\bullet)$에
대응한다. 이 보조정리의 주장은 $\mu$가 $B$의 작용과 가환하는 것과
$\gamma$의 상이
$$
\Hom_{\text{Mod}^{dg}_{(B, \text{d})}}(M, M) \subset
\Hom^\bullet(M^\bullet, M^\bullet)
$$
에 포함되는 것이 동치라는 말이다. 자세한 계산은 생략한다.
\end{proof}

\noindent
$M$을 미분 등급 $(A, B)$-쌍가군이라 하자. Section
\ref{dga-section-left-modules}에서 왼쪽 미분 등급 $A$-가군 구조는 오른쪽
미분 등급 $A^{opp}$-가군 구조에 대응함을 상기하자. $A$-가군 구조와
$B$-가군 구조가 서로 가환하므로, 이는 $M$에 미분 등급
$A^{opp} \otimes_R B$-가군 구조를 준다.
$$
x \cdot (a \otimes b) = (-1)^{\deg(a)\deg(x)} axb
$$
역으로 미분 등급 $A^{opp} \otimes_R B$-가군 $M$이 주어지면,
위 식을 사용하여 미분 등급 $(A, B)$-쌍가군을 얻을 수 있다.

\begin{lemma}
\label{dga-lemma-bimodule-over-tensor}
$R$을 환이라 하자. $(A, \text{d})$와 $(B, \text{d})$를
$R$ 위의 미분 등급 대수라 하자. 위 구성은 다음 범주의 동치를 정의한다.
$$
\begin{matrix}
\text{미분 등급}\\
(A, B)\text{-쌍가군}
\end{matrix}
\longleftrightarrow
\begin{matrix}
\text{오른쪽 미분 등급 }\\
A^{opp} \otimes_R B\text{-가군}
\end{matrix}
$$
\end{lemma}

\begin{proof}
위의 논의에서 곧바로 따른다.
% P03-STACKS-E935: frozen English says “from discussion the above”.
\end{proof}

\noindent
$R$을 환이라 하자. $(A, \text{d})$와 $(B, \text{d})$를 미분 등급
$R$-대수라 하자. $P$를 미분 등급 $(A, B)$-쌍가군이라 하자.
$P$가 {\it 성질 (P)를 갖는다}고 함은 미분 등급 $(A, B)$-쌍가군들에 의한
필트레이션
% P03-STACKS-E936: frozen English duplicates “if it there exists”.
$$
0 = F_{-1}P \subset F_0P \subset F_1P \subset \ldots \subset P
$$
이 존재하여 다음을 만족하는 것을 말한다.
\begin{enumerate}
\item $P = \bigcup F_pP$,
\item 포함사상 $F_iP \to F_{i + 1}P$는 등급 $(A, B)$-쌍가군 사상으로서
분할되고,
\item 몫 $F_{i + 1}P/F_iP$는 미분 등급 $(A, B)$-쌍가군으로서
$(A \otimes_R B)[k]$들의 직합과 동형이다.
\end{enumerate}

\begin{lemma}
\label{dga-lemma-bimodule-resolve}
$R$을 환이라 하자. $(A, \text{d})$와 $(B, \text{d})$를 미분 등급
$R$-대수라 하자. $M$을 미분 등급 $(A, B)$-쌍가군이라 하자.
$P$가 위에서 정의한 성질 (P)를 가지며, 유사동형인 미분 등급
$(A, B)$-쌍가군 준동형 $P \to M$이 존재한다.
\end{lemma}

\begin{proof}
Lemmas \ref{dga-lemma-bimodule-over-tensor}와
\ref{dga-lemma-resolve}에서 곧바로 따른다.
\end{proof}

\begin{lemma}
\label{dga-lemma-bimodule-property-P-sequence}
$R$을 환이라 하자. $(A, \text{d})$와 $(B, \text{d})$를 미분 등급
$R$-대수라 하자. $P$를 성질 (P)를 갖는 미분 등급 $(A, B)$-쌍가군이라 하고,
대응하는 필트레이션을 $F_\bullet$이라 하자. 그러면 미분 등급
$(A, B)$-쌍가군들의 짧은 완전열
% P03-STACKS-E937: frozen English uses a comma-spliced “Let ..., then”.
$$
0 \to
\bigoplus\nolimits F_iP \to
\bigoplus\nolimits F_iP \to P \to 0
$$
을 얻으며, 이는 등급 $(A, B)$-쌍가군들의 열로서 분할된다.
\end{lemma}

\begin{proof}
Lemmas \ref{dga-lemma-bimodule-over-tensor}와
\ref{dga-lemma-property-P-sequence}에서 곧바로 따른다.
\end{proof}















\section{쌍가군과 텐서곱}
\label{dga-section-bimodules-tensor}

\noindent
$R$을 환이라 하자. $A$와 $B$를 $R$-대수라 하자. $M$을 오른쪽
$A$-가군이라 하자. $N$을 $(A, B)$-쌍가군이라 하자. 그러면
$M \otimes_A N$은 오른쪽 $B$-가군이다.

\medskip\noindent
앞 단락의 상황에서 $A$와 $B$가 $\mathbf{Z}$-등급 대수이고,
$M$이 등급 $A$-가군이며, $N$이 등급 $(A, B)$-쌍가군이면,
$M \otimes_A N$은 오른쪽 등급 $B$-가군이다. 이 구성은 $M$에 대하여
함자적이며 다음 함자를 정의한다.
$$
- \otimes_A N :
\text{Mod}^{gr}_A
\longrightarrow
\text{Mod}^{gr}_B
$$
이는 Example \ref{dga-example-gm-gr-cat}에서와 같은 등급 범주들의 함자이다.
실제로 $M$과 $M'$이 등급 $A$-가군이고 $f : M \to M'$가 차수 $n$인
동차 $A$-가군 준동형이면,
$f \otimes \text{id}_N : M \otimes_A N \to M' \otimes_A N$은
차수 $n$인 동차 $B$-가군 준동형이다.

\medskip\noindent
앞 단락의 상황에서 $(A, \text{d})$와 $(B, \text{d})$가 미분 등급 대수이고,
$M$이 미분 등급 $A$-가군이며, $N$이 미분 등급 $(A, B)$-쌍가군이면,
$M \otimes_A N$은 오른쪽 미분 등급 $B$-가군이다.

\begin{lemma}
\label{dga-lemma-tensor}
$R$을 환이라 하자. $(A, \text{d})$와 $(B, \text{d})$를
$R$ 위의 미분 등급 대수라 하자. $N$을 미분 등급 $(A, B)$-쌍가군이라 하자.
그러면 $M \mapsto M \otimes_A N$은 다음 함자를 정의한다.
$$
- \otimes_A N :
\text{Mod}^{dg}_{(A, \text{d})}
\longrightarrow
\text{Mod}^{dg}_{(B, \text{d})}
$$
이는 미분 등급 범주들의 함자이다. Lemma \ref{dga-lemma-functorial}을 적용하면
이 함자는 다음 함자들을 유도한다.
$$
\text{Mod}_{(A, \text{d})} \to \text{Mod}_{(B, \text{d})}
\quad\text{및}\quad
K(\text{Mod}_{(A, \text{d})}) \to K(\text{Mod}_{(B, \text{d})})
$$
\end{lemma}

\begin{proof}
위에서 이 구성이 바탕 등급 범주들의 함자를 어떻게 정의하는지 보았다.
따라서 이 구성이 미분과 양립함을 보이면 충분하다. $M$과 $M'$을 미분 등급
$A$-가군이라 하고, $f : M \to M'$를 차수 $n$인 동차 $A$-가군
준동형이라 하자. 그러면
$$
\text{d}(f) = \text{d}_{M'} \circ f - (-1)^n f \circ \text{d}_M
$$
이다. 한편,
$$
\text{d}(f \otimes \text{id}_N) =
\text{d}_{M' \otimes_A N} \circ
(f \otimes \text{id}_N)
- (-1)^n
(f \otimes \text{id}_N) \circ \text{d}_{M \otimes_A N}
$$
이다. 이를 $x \in M$과 $y \in N$이 동차인 원소 $x \otimes y$에
적용하면 다음을 얻는다.
\begin{align*}
\text{d}(f \otimes \text{id}_N)(x \otimes y)
= &
\text{d}_{M'}(f(x)) \otimes y + (-1)^{n + \deg(x)}f(x) \otimes \text{d}_N(y) \\
& - (-1)^n f(\text{d}_M(x)) \otimes y
- (-1)^{n + \deg(x)}f(x) \otimes \text{d}_N(y) \\
= &
\text{d}(f) (x \otimes y)
\end{align*}
% P03-STACKS-E939: frozen final line is ill-typed; retained literally.
따라서 $\text{d}(f) \otimes \text{id}_N =
\text{d}(f \otimes \text{id}_N)$임을 알 수 있고, 증명이 끝난다.
\end{proof}

\begin{remark}
\label{dga-remark-shift-tensor-no-sign}
$R$을 환이라 하자. $(A, \text{d})$와 $(B, \text{d})$를
$R$ 위의 미분 등급 대수라 하자. $N$을 미분 등급 $(A, B)$-쌍가군이라 하고,
$M$을 오른쪽 미분 등급 $A$-가군이라 하자. 그러면 모든
$k \in \mathbf{Z}$에 대하여 동형
$$
(M \otimes_A N)[k] \longrightarrow
M[k] \otimes_A N
$$
이 있으며, 이는 부호를 도입하지 않고 정의되는 오른쪽 미분 등급
$B$-가군의 동형이다. More on Algebra, Section
\ref{more-algebra-section-sign-rules}을 보라.
\end{remark}

\noindent
환 $R$, $R$-대수 $A$, $B$, $C$, 오른쪽 $A$-가군 $M$,
$(A, B)$-쌍가군 $N$, 그리고 $(B, C)$-쌍가군 $N'$이 주어지면,
$N \otimes_B N'$은 $(A, C)$-쌍가군이고 다음이 성립한다.
$$
(M \otimes_A N) \otimes_B N' = M \otimes_A (N \otimes_B N')
$$
이 등식은 등급의 경우와 미분 등급의 경우에도 계속 성립한다.
% P03-STACKS-E938: frozen English says “continuous to hold”.
부호 규칙은 More on Algebra, Section \ref{more-algebra-section-sign-rules}을 보라.






\section{쌍가군과 내부 Hom}
\label{dga-section-bimodules-hom}

\noindent
$R$을 환이라 하자. $A$가 $R$-대수이고(Section
\ref{dga-section-conventions}의 관례를 보라) $M$, $M'$이 오른쪽
$A$-가군이면, 통상적으로 다음과 같이 정의한다.
$$
\Hom_A(M, M') = \{f : M \to M' \mid f \text{는 }A\text{-선형이다}\}
$$

\medskip\noindent
$R$을 환이라 하자. $A$와 $B$를 $R$-대수라 하자. $N$을
$(A, B)$-쌍가군이라 하고 $N'$을 오른쪽 $B$-가군이라 하자.
% P03-STACKS-E940: frozen English twice says “R-be”.
이 상황에서 전합성을 이용하여
$$
\Hom_B(N, N')
$$
을 오른쪽 $A$-가군으로 생각한다.

\medskip\noindent
$R$을 환이라 하자. $A$와 $B$를 $\mathbf{Z}$-등급 $R$-대수라 하자.
$N$을 등급 $(A, B)$-쌍가군이라 하고, $N'$을 오른쪽 등급 $B$-가군이라 하자.
이 상황에서 Example \ref{dga-example-gm-gr-cat}에 정의된 등급 $R$-가군
$$
\Hom_{\text{Mod}^{gr}_B}(N, N')
$$
을 전합성을 이용하여 오른쪽 등급 $A$-가군으로 생각한다. 이 구성은
$N'$에 대하여 함자적이며 다음 함자를 정의한다.
$$
\Hom_{\text{Mod}^{gr}_B}(N, -) :
\text{Mod}^{gr}_B
\longrightarrow
\text{Mod}^{gr}_A
$$
이는 Example \ref{dga-example-gm-gr-cat}에서와 같은 등급 범주들의 함자이다.
실제로 $N_1$과 $N_2$가 등급 $B$-가군이고 $f : N_1 \to N_2$가 차수
$n$인 동차 $B$-가군 준동형이면, 유도된 사상
$\Hom_{\text{Mod}^{gr}_B}(N, N_1) \to \Hom_{\text{Mod}^{gr}_B}(N, N_2)$
은 차수 $n$인 동차 $A$-가군 준동형이다.

\medskip\noindent
$R$을 환이라 하자. $A$와 $B$를 미분 $\mathbf{Z}$-등급 $R$-대수라 하자.
$N$을 미분 등급 $(A, B)$-쌍가군이라 하고, $N'$을 오른쪽 미분 등급
$B$-가군이라 하자. 이 상황에서 Example \ref{dga-example-dgm-dg-cat}에
정의된 미분 등급 $R$-가군
$$
\Hom_{\text{Mod}^{dg}_{(B, \text{d})}}(N, N')
$$
을 등급의 경우와 같이 전합성을 이용하여 오른쪽 미분 등급 $A$-가군으로
생각한다. 곱셈은 다음 합성이므로 이는 미분과 양립한다.
% P03-STACKS-E941: frozen display omits (B,d) in four Mod-category subscripts.
$$
\begin{aligned}
&
\Hom_{\text{Mod}^{dg}_B}(N, N') \otimes_R A \to
\Hom_{\text{Mod}^{dg}_B}(N, N') \otimes_R \\
&
\Hom_{\text{Mod}^{dg}_B}(N, N) \to
\Hom_{\text{Mod}^{dg}_B}(N, N')
\end{aligned}
$$
첫 번째 화살표는 Lemma \ref{dga-lemma-what-makes-a-bimodule-dg}의 사상을
사용하고, 두 번째 화살표는 미분 등급 범주
$\text{Mod}^{dg}_{(B, \text{d})}$에서의 합성이다.

\begin{lemma}
\label{dga-lemma-hom}
$R$을 환이라 하자. $(A, \text{d})$와 $(B, \text{d})$를
$R$ 위의 미분 등급 대수라 하자. $N$을 미분 등급 $(A, B)$-쌍가군이라 하자.
위 구성은 다음 함자를 정의한다.
$$
\Hom_{\text{Mod}^{dg}_{(B, \text{d})}}(N, -) :
\text{Mod}^{dg}_{(B, \text{d})}
\longrightarrow
\text{Mod}^{dg}_{(A, \text{d})}
$$
이는 미분 등급 범주들의 함자이다. Lemma \ref{dga-lemma-functorial}을 적용하면
이 함자는 다음 함자들을 유도한다.
$$
\text{Mod}_{(B, \text{d})} \to \text{Mod}_{(A, \text{d})}
\quad\text{및}\quad
K(\text{Mod}_{(B, \text{d})}) \to K(\text{Mod}_{(A, \text{d})})
$$
\end{lemma}

\begin{proof}
위에서 이 구성이 바탕 등급 범주들의 함자를 어떻게 정의하는지 보았다.
따라서 이 구성이 미분과 양립함을 보이면 충분하다. $N_1$과 $N_2$를
미분 등급 $B$-가군이라 하자. 다음과 같이 쓰자.
$$
\begin{aligned}
&
H_{12} = \Hom_{\text{Mod}^{dg}_{(B, \text{d})}}(N_1, N_2),\quad
H_1 = \Hom_{\text{Mod}^{dg}_{(B, \text{d})}}(N, N_1),\quad \\
&
H_2 = \Hom_{\text{Mod}^{dg}_{(B, \text{d})}}(N, N_2)
\end{aligned}
$$
미분 등급 범주 $\text{Mod}^{dg}_{(B, \text{d})}$에서의 합성
$$
c : H_{12} \otimes_R H_1 \longrightarrow H_2
$$
을 생각하자. $f : N_1 \to N_2$를 차수 $n$인 동차 $B$-가군 준동형이라
하자. 다시 말해 $f \in H_{12}^n$이다. 보조정리의 함자는 $f$를
$c_f : H_1 \to H_2$, $g \mapsto c(f, g)$로 보낸다. $\text{d}(f)$에
대해서도 마찬가지이다. 한편,
$$
\Hom_{\text{Mod}^{dg}_{(A, \text{d})}}(H_1, H_2)
$$
위의 미분은 $c_f$를
$\text{d}_{H_2} \circ c_f - (-1)^n c_f \circ \text{d}_{H_1}$로 보낸다.
$c$는 $R$-가군 복합체의 사상이므로
$\text{d} c(f, g) = c(\text{d}f, g) + (-1)^n c(f, \text{d}g)$이다.
따라서
\begin{align*}
(\text{d}c_f)(g)
& =
\text{d}c(f,g) - (-1)^n c(f, \text{d}g) \\
& =
c(\text{d}f, g) + (-1)^n c(f, \text{d}g)  - (-1)^n c(f, \text{d}g) \\
& =
c(\text{d}f, g) = c_{\text{d}f}(g)
\end{align*}
이고, 증명이 끝난다.
\end{proof}

\begin{remark}
\label{dga-remark-shift-hom-no-sign}
$R$을 환이라 하자. $(A, \text{d})$와 $(B, \text{d})$를
$R$ 위의 미분 등급 대수라 하자. $N$을 미분 등급 $(A, B)$-쌍가군이라 하고,
$N'$을 오른쪽 미분 등급 $B$-가군이라 하자. 그러면 모든
$k \in \mathbf{Z}$에 대하여 동형
$$
\Hom_{\text{Mod}^{gr}_B}(N, N')[k]
\longrightarrow
\Hom_{\text{Mod}^{gr}_B}(N, N'[k])
$$
이 있으며, 이는 부호를 도입하지 않고 정의되는 오른쪽 미분 등급
$A$-가군의 동형이다. More on Algebra, Section
\ref{more-algebra-section-sign-rules}을 보라.
\end{remark}

\begin{lemma}
\label{dga-lemma-tensor-hom-adjunction}
$R$을 환이라 하자. $A$와 $B$를 $R$-대수라 하자.
$M$을 오른쪽 $A$-가군, $N$을 $(A, B)$-쌍가군, $N'$을 오른쪽
$B$-가군이라 하자. 그러면 표준 동형
$$
\Hom_B(M \otimes_A N, N') = \Hom_A(M, \Hom_B(N, N'))
$$
인 $R$-가군의 동형이 있다. $A$, $B$, $M$, $N$, $N'$이 서로 양립하게
등급지어졌다면 표준 동형
$$
\Hom_{\text{Mod}_B^{gr}}(M \otimes_A N, N') =
\Hom_{\text{Mod}_A^{gr}}(M, \Hom_{\text{Mod}_B^{gr}}(N, N'))
$$
인 등급 $R$-가군의 동형이 있다.
% P03-STACKS-E942: frozen English omits terminal punctuation here.
$A$, $B$, $M$, $N$, $N'$이 서로 양립하게 미분 등급지어졌다면 표준 동형
$$
\Hom_{\text{Mod}^{dg}_{(B, \text{d})}}(M \otimes_A N, N') =
\Hom_{\text{Mod}^{dg}_{(A, \text{d})}}(M,
\Hom_{\text{Mod}^{dg}_{(B, \text{d})}}(N, N'))
$$
인 $R$-가군 복합체의 동형이 있다.
\end{lemma}

\begin{proof}
생략한다. 힌트: 등급 없는 경우에는 양변을 오른쪽에서 $B$-선형인
$A$-쌍선형 사상 $\psi : M \times N \to N'$로 해석하라.
% P03-STACKS-E943: frozen source says A-bilinear rather than A-balanced.
(미분) 등급의 경우에는 More on Algebra, Lemma
\ref{more-algebra-lemma-compose}의 동형을 사용하고, 그것이 가군 구조와
양립함을 확인하라. 또는 Lemma \ref{dga-lemma-characterize-hom}의 동형을
사용하여 그것이 $B$-가군 구조와 양립함을 보여라.
\end{proof}







\section{유도 Hom}
\label{dga-section-restriction}

\noindent
이 절은 More on Algebra, Section \ref{more-algebra-section-RHom}과 유사하다.

\medskip\noindent
$R$을 환이라 하자. $(A, \text{d})$와 $(B, \text{d})$를
$R$ 위의 미분 등급 대수라 하자. $N$을 미분 등급 $(A, B)$-쌍가군이라 하자.
Section \ref{dga-section-bimodules-hom}의 함자
\begin{equation}
\label{dga-equation-restriction}
\Hom_{\text{Mod}^{dg}_{(B, \text{d})}}(N, -) :
\text{Mod}_{(B, \text{d})}
\longrightarrow
\text{Mod}_{(A, \text{d})}
\end{equation}
를 생각하자.

\begin{lemma}
\label{dga-lemma-restriction-homotopy}
함자 (\ref{dga-equation-restriction})은 삼각범주들의 완전 함자
$K(\text{Mod}_{(B, \text{d})}) \to K(\text{Mod}_{(A, \text{d})})$
를 정의한다.
\end{lemma}

\begin{proof}
Lemma \ref{dga-lemma-hom}과 Remark \ref{dga-remark-shift-hom-no-sign}을 통하여,
이는 Lemma \ref{dga-lemma-functor-between-ABC}의 일반 원리에서 따른다.
\end{proof}

\noindent
Lemma \ref{dga-lemma-restriction-homotopy}에 의하면 삼각범주들의 완전 함자
$$
\Hom_{\text{Mod}^{dg}_{(B, \text{d})}}(N, -) :
K(\text{Mod}_{(B, \text{d})}) \to K(\text{Mod}_{(A, \text{d})})
$$
가 있음을 상기하자. 다음 도식을 생각하자.
$$
\xymatrix{
K(\text{Mod}_{(B, \text{d})}) \ar[d] \ar[rr]_{\text{위 참조}} \ar[rrd]_F & &
K(\text{Mod}_{(A, \text{d})}) \ar[d] \\
D(B, \text{d}) \ar@{..>}[rr] & &
D(A, \text{d})
}
$$
합성 $F$의 {\it 오른쪽 유도함자}로서 점선 화살표를 구성하고자 한다.
({\it 주의}: 흥미로운 대부분의 경우 이 도식은 가환하지 않는다.)
즉 Derived Categories, Section \ref{derived-section-derived-functors}의
일반적 설정에서, $K(\text{Mod}_{(A, \text{d})})$ 안의 유사동형들로 이루어진
곱셈적 계에 관하여 $F$의 오른쪽 유도함자를 계산하고자 한다.
% P03-STACKS-E567: frozen source names K(Mod_(A,d)) although F starts at K(Mod_(B,d)).

\begin{lemma}
\label{dga-lemma-derived-restriction}
위 상황에서 $F$의 오른쪽 유도함자가 존재한다. 이를
$R\Hom(N, -) : D(B, \text{d}) \to D(A, \text{d})$로 나타낸다.
\end{lemma}

\begin{proof}
이를 증명하기 위해 Derived Categories, Lemma
\ref{derived-lemma-find-existence-computes}를 사용한다. 객체들의 모음
$\mathcal{I}$로 성질 (I)를 갖는 객체들을 사용한다. 성질 (1)은 Lemma
\ref{dga-lemma-right-resolution}에서 보였다. 성질 (2)도 성립한다. 실제로
$s : I \to I'$가 성질 (I)를 갖는 가군들의 유사동형이면, Lemma
\ref{dga-lemma-hom-derived}에 의해 $s$는 호모토피 동치이다.
\end{proof}

\begin{lemma}
\label{dga-lemma-functoriality-derived-restriction}
$R$을 환이라 하자. $(A, \text{d})$와 $(B, \text{d})$를 미분 등급
$R$-대수라 하자. $f : N \to N'$를 미분 등급 $(A, B)$-쌍가군의
준동형이라 하자. 그러면 $f$는 함자들의 사상
$$
- \circ f : R\Hom(N', -) \longrightarrow R\Hom(N, -)
$$
을 유도한다. $f$가 유사동형이면 $f \circ -$는 함자들의 동형이다.
% P03-STACKS-E568: frozen source reverses the composition in the preceding sentence and proof.
\end{lemma}

\begin{proof}
미분 등급 $B$-가군들의 미분 등급 범주를
$\mathcal{B} = \text{Mod}^{dg}_{(B, \text{d})}$로 쓰자. Example
\ref{dga-example-dgm-dg-cat}를 보라.
% P03-STACKS-E944: frozen English omits the linking preposition in this notation clause.
$I$를 성질 (I)를 갖는 미분 등급 $B$-가군이라 하자. 그러면
$f \circ - : \Hom_\mathcal{B}(N', I) \to \Hom_\mathcal{B}(N, I)$는
미분 등급 $A$-가군의 사상이다. 더욱이 이는 $I$에 대하여 함자적이다.
함자 $ R\Hom(N', -)$와 $R\Hom(N, -)$는 성질 (I)를 갖는 객체에
$\Hom_\mathcal{B}$를 적용하여 계산되므로(Lemma
\ref{dga-lemma-derived-restriction}), 표시한 함자 변환을 얻는다.

\medskip\noindent
$f$가 유사동형이라고 가정하자. $F_\bullet$을 $I$ 위의 주어진 필트레이션이라
하자. $I = \lim I/F_pI$이므로
$\Hom_\mathcal{B}(N', I) = \lim \Hom_\mathcal{B}(N', I/F_pI)$ 및
$\Hom_\mathcal{B}(N, I) = \lim \Hom_\mathcal{B}(N, I/F_pI)$임을 알 수 있다.
계 $I/F_pI$의 전이 사상들은 등급 가군으로서 분할되므로, 계
$\Hom_\mathcal{B}(N', I/F_pI)$와 $\Hom_\mathcal{B}(N, I/F_pI)$의 전이
사상들은 전사이다. 따라서 $\Hom_\mathcal{B}(N', I)$ 및 각각
$\Hom_\mathcal{B}(N, I)$를 아벨군의 복합체로 보면 복합체 계
$\Hom_\mathcal{B}(N', I/F_pI)$ 및 각각 $\Hom_\mathcal{B}(N, I/F_pI)$의
$R\lim$을 계산한다. More on Algebra, Lemma
\ref{more-algebra-lemma-compute-Rlim}를 보라. 그러므로 각 사상
$$
\Hom_\mathcal{B}(N', I/F_pI) \to \Hom_\mathcal{B}(N, I/F_pI)
$$
이 유사동형임을 보이면 충분하다. 전사
$I/F_{p + 1}I \to I/F_pI$는 등급 $B$-가군의 사상으로서 분할되므로
$$
0 \to \Hom_\mathcal{B}(N', F_pI/F_{p + 1}I) \to
\Hom_\mathcal{B}(N', I/F_{p + 1}I) \to
\Hom_\mathcal{B}(N', I/F_pI) \to 0
$$
은 미분 등급 $A$-가군들의 짧은 완전열이다. $N$에 대해서도 유사한 열이
있고, $f$는 짧은 완전열들의 사상을 유도한다. 따라서 $p$에 대한 귀납법을
사용하면($I/F_0I = 0$인 $p = 0$에서 시작한다), 사상
$\Hom_\mathcal{B}(N', F_pI/F_{p + 1}I) \to \Hom_\mathcal{B}(N, F_pI/F_{p + 1}I)$
이 유사동형임을 보이면 충분하다는 결론을 얻는다. $F_pI/F_{p + 1}I$는
$A^\vee$의 이동들의 곱이므로
% P03-STACKS-E569/P03-STACKS-E570: frozen source says A^vee and “it suffice”.
$\Hom_\mathcal{B}(N', B^\vee[k]) \to \Hom_\mathcal{B}(N, B^\vee[k])$
이 유사동형임을 보이면 충분하다. Lemma
\ref{dga-lemma-hom-into-shift-dual-free}에 의하면
$(N')^\vee \to N^\vee$가 유사동형임을 보이면 충분하다. 이는 $f$가
유사동형이고 $(\ )^\vee$가 완전 함자이므로 참이다.
\end{proof}

\begin{lemma}
\label{dga-lemma-derived-restriction-exts}
$(A, \text{d})$와 $(B, \text{d})$를 환 $R$ 위의 미분 등급 대수라 하자.
$N$을 미분 등급 $(A, B)$-쌍가군이라 하자. 그러면 모든
$n \in \mathbf{Z}$에 대하여 동형
$$
H^n(R\Hom(N, M)) = \Ext^n_{D(B, \text{d})}(N, M)
$$
인 $R$-가군의 동형이 있으며, 이는 $M$에 대하여 함자적이다. 또한 Lemma
\ref{dga-lemma-functoriality-derived-restriction}에 기술된 연산에 관하여
$N$에 대해서도 함자적이다.
\end{lemma}

\begin{proof}
Lemma \ref{dga-lemma-derived-restriction}의 증명에서
$$
R\Hom(N, M) = 
\Hom_{\text{Mod}^{dg}_{(B, \text{d})}}(N, I)
$$
임을 미분 등급 $A$-가군으로서 보았다. 여기서 $M \to I$는 $M$에서 성질
(I)를 갖는 미분 등급 $B$-가군으로 가는 유사동형이다. 따라서 Lemma
\ref{dga-lemma-hom-derived}에 의해 이 복합체는 올바른 코호몰로지 가군들을
갖는다. 이러한 동일시의 함자적 성질에 대한 논의는 생략한다.
\end{proof}

\begin{lemma}
\label{dga-lemma-compute-derived-restriction}
$R$을 환이라 하자. $(A, \text{d})$와 $(B, \text{d})$를 미분 등급
$R$-대수라 하자. $N$을 미분 등급 $(A, B)$-쌍가군이라 하자. 모든
$N' \in K(B, \text{d})$에 대하여
$\Hom_{D(B, \text{d})}(N, N') = \Hom_{K(\text{Mod}_{(B, \text{d})})}(N, N')$
이면, 예를 들어 $N$이 미분 등급 $B$-가군으로서 성질 (P)를 가지면,
% P03-STACKS-E571: frozen source quantifies N' in K(B,d), not K(Mod_(B,d)).
$$
R\Hom(N, M) = \Hom_{\text{Mod}^{dg}_{(B, \text{d})}}(N, M)
$$
가 $D(B, \text{d})$ 안의 $M$에 대하여 함자적으로 성립한다.
\end{lemma}

\begin{proof}
구성에 따라(Lemma \ref{dga-lemma-derived-restriction}), $R\Hom(N, M)$을
구하기 위해 성질 (I)를 갖는 미분 등급 $B$-가군 $I$로 가는 유사동형
$M \to I$를 택하고
$R\Hom(N, M) = \Hom_{\text{Mod}^{dg}_{(B, \text{d})}}(N, I)$로 둔다.
가정에 의하면 $M \to I$가 유도하는 사상
$$
\Hom_{\text{Mod}^{dg}_{(B, \text{d})}}(N, M) \longrightarrow
\Hom_{\text{Mod}^{dg}_{(B, \text{d})}}(N, I)
$$
은 유사동형이다. Example \ref{dga-example-dgm-dg-cat}의 논의를 보라.
이로써 보조정리가 증명된다. $N$이 $B$-가군으로서 성질 (P)를 가지면,
Lemma \ref{dga-lemma-hom-derived}에 의해 가정이 만족됨을 알 수 있다.
\end{proof}




\section{유도 Hom의 변형}
\label{dga-section-variant}

\noindent
$\mathcal{A}$를 아벨 범주라 하자. $\mathcal{A}$의 복합체들로 이루어진
미분 등급 범주 $\text{Comp}^{dg}(\mathcal{A})$를 생각하자. Example
\ref{dga-example-category-complexes}를 보라. $K^\bullet$을 $\mathcal{A}$의
복합체라 하자. 다음과 같이 두자.
$$
(E, \text{d}) = \Hom_{\text{Comp}^{dg}(\mathcal{A})}(K^\bullet, K^\bullet)
$$
그리고 Lemma \ref{dga-lemma-construction}의 미분 등급 범주들의 함자
$$
\text{Comp}^{dg}(\mathcal{A}) \longrightarrow \text{Mod}^{dg}_{(E, \text{d})},
\quad
X^\bullet
\longmapsto
\Hom_{\text{Comp}^{dg}(\mathcal{A})}(K^\bullet, X^\bullet)
$$
를 생각하자.

\begin{lemma}
\label{dga-lemma-existence-of-derived}
위의 상황을 가정하자. $\Hom(K^\bullet, -) : K(\mathcal{A}) \to D(\textit{Ab})$의
오른쪽 유도함자 $R\Hom(K^\bullet, -)$가 $D(\mathcal{A})$ 전체에서
정의되면, 삼각범주들의 표준 완전 함자
% P03-STACKS-E945: frozen English opens with a sentence fragment.
$$
R\Hom(K^\bullet, -) : D(\mathcal{A}) \longrightarrow D(E, \text{d})
$$
를 얻으며, 연관된 아벨군 복합체를 취하면 통상적인 함자로 환원된다.
% P03-STACKS-E946: frozen English contains a doubled space in “one  on”.
\end{lemma}

\begin{proof}
Lemma \ref{dga-lemma-construction}에 의해 연관된 함자
$K(\mathcal{A}) \to K(\text{Mod}_{(E, \text{d})})$가 있음에 주목하자.
이 함자가 삼각범주들의 완전 함자임을 주장한다. 실제로
$f : A^\bullet \to B^\bullet$를 $\mathcal{A}$의 복합체들의 사상이라 하자.
그러면 계산으로 다음을 알 수 있다.
$$
\begin{aligned}
&
\Hom_{\text{Comp}^{dg}(\mathcal{A})}(K^\bullet, C(f)^\bullet)
\\
&=
C\left(
\Hom_{\text{Comp}^{dg}(\mathcal{A})}(K^\bullet, A^\bullet) \to
\Hom_{\text{Comp}^{dg}(\mathcal{A})}(K^\bullet, B^\bullet)
\right)
\end{aligned}
$$
여기서 우변은 이 장의 앞부분에서 정의한 $\text{Mod}_{(E, \text{d})}$의
원뿔이다. 이는 우리의 함자가 원뿔과 양립하고, 따라서 특수 삼각형과
양립함을 보인다. $X^\bullet$을 $K(\mathcal{A})$의 객체라 하자.
유사동형 $s : X^\bullet \to Y^\bullet$들의 범주를 생각하자. 함자
$(s : X^\bullet \to Y^\bullet) \mapsto \Hom_\mathcal{A}(K^\bullet, Y^\bullet)$
가 $D(\textit{Ab})$에서 보았을 때 본질적으로 상수임이 주어져 있다.
% P03-STACKS-E947: frozen source switches here to undeclared Hom_A; retained literally.
그런데 망각 함자 $D(E, \text{d}) \to D(\textit{Ab})$는 코호몰로지를
취하는 것과 양립하므로, $D(E, \text{d})$에서도 마찬가지이다.
이로써 보조정리가 증명된다.
\end{proof}

\noindent
{\bf 주의:} 이 보조정리는 서술된 그대로 성립하며 그 형태로 유용할 수
있지만, 모든 $n \in \mathbf{Z}$에 대하여
$H^n(E) = \Ext^n_{D(\mathcal{A})}(K^\bullet, K^\bullet)$
가 성립하지 않는 한 미분 대수 $E$는 ``올바른'' 것이 아니다.





\section{유도 텐서곱}
\label{dga-section-base-change}

\noindent
이 절은 More on Algebra, Section
\ref{more-algebra-section-derived-base-change}와 유사하다.

\medskip\noindent
$R$을 환이라 하자. $(A, \text{d})$와 $(B, \text{d})$를
$R$ 위의 미분 등급 대수라 하자. $N$을 미분 등급 $(A, B)$-쌍가군이라 하자.
Section \ref{dga-section-bimodules-tensor}에서 정의한 함자
\begin{equation}
\label{dga-equation-bc}
\text{Mod}_{(A, \text{d})}
\longrightarrow
\text{Mod}_{(B, \text{d})},\quad
M \longmapsto M \otimes_A N
\end{equation}
를 생각하자.

\begin{lemma}
\label{dga-lemma-bc-homotopy}
함자 (\ref{dga-equation-bc})은 삼각범주들의 완전 함자
$K(\text{Mod}_{(A, \text{d})}) \to K(\text{Mod}_{(B, \text{d})})$.
를 정의한다.
\end{lemma}

\begin{proof}
Lemma \ref{dga-lemma-tensor}와 Remark \ref{dga-remark-shift-tensor-no-sign}을
통하여, 이는 Lemma \ref{dga-lemma-functor-between-ABC}의 일반 원리에서 따른다.
\end{proof}

\noindent
이제 다음 도식을 생각할 수 있다.
$$
\xymatrix{
K(\text{Mod}_{(A, \text{d})}) \ar[d] \ar[rr]_{- \otimes_A N} \ar[rrd]_F & &
K(\text{Mod}_{(B, \text{d})}) \ar[d] \\
D(A, \text{d}) \ar@{..>}[rr] & &
D(B, \text{d})
}
$$
아래에서 구성할 점선 화살표는 합성 $F$의 {\it 왼쪽 유도함자}가 된다.
({\it 주의}: 이 도식은 가환하지 않는다.) 즉 Derived Categories, Section
\ref{derived-section-derived-functors}의 일반적 설정에서,
$K(\text{Mod}_{(A, \text{d})})$ 안의 유사동형들로 이루어진 곱셈적 계에
관하여 $F$의 왼쪽 유도함자를 계산하고자 한다.

\begin{lemma}
\label{dga-lemma-derived-bc}
위 상황에서 $F$의 왼쪽 유도함자가 존재한다. 이를
$- \otimes_A^\mathbf{L} N : D(A, \text{d}) \to D(B, \text{d})$.
로 나타낸다.
\end{lemma}

\begin{proof}
이를 증명하기 위해 Derived Categories, Lemma
\ref{derived-lemma-find-existence-computes}를 사용한다. 객체들의 모음
$\mathcal{P}$로 성질 (P)를 갖는 객체들을 사용한다. 성질 (1)은 Lemma
\ref{dga-lemma-resolve}에서 보였다. 성질 (2)도 성립한다. 실제로
$s : P \to P'$가 성질 (P)를 갖는 가군들의 유사동형이면, Lemma
\ref{dga-lemma-hom-derived}에 의해 $s$는 호모토피 동치이다.
\end{proof}

\begin{lemma}
\label{dga-lemma-functoriality-bc}
$R$을 환이라 하자. $(A, \text{d})$와 $(B, \text{d})$를 미분 등급
$R$-대수라 하자. $f : N \to N'$를 미분 등급 $(A, B)$-쌍가군의
준동형이라 하자. 그러면 $f$는 함자들의 사상
$$
1\otimes f :
- \otimes_A^\mathbf{L} N
\longrightarrow
- \otimes_A^\mathbf{L} N'
$$
을 유도한다. $f$가 유사동형이면 $1 \otimes f$는 함자들의 동형이다.
\end{lemma}

\begin{proof}
$M$을 성질 (P)를 갖는 미분 등급 $A$-가군이라 하자. 그러면
$1 \otimes f : M \otimes_A N \to M \otimes_A N'$는 미분 등급
$B$-가군의 사상이다. 더욱이 이는 $M$에 대하여 함자적이다. 함자
$- \otimes_A^\mathbf{L} N$과 $- \otimes_A^\mathbf{L} N'$는 성질 (P)를
갖는 객체에서 텐서곱을 취하여 계산되므로(Lemma
\ref{dga-lemma-derived-bc}), 표시한 함자 변환을 얻는다.

\medskip\noindent
$f$가 유사동형이라고 가정하자. $F_\bullet$을 $M$ 위의 주어진 필트레이션이라
하자. 다음을 주목하자.
$M \otimes_A N = \colim F_i(M) \otimes_A N$ 및
$M \otimes_A N' = \colim F_i(M) \otimes_A N'$.
따라서
$F_n(M) \otimes_A N \to F_n(M) \otimes_A N'$
이 유사동형임을 보이면 충분하다(필터 여극한은 완전하다. Algebra, Lemma
\ref{algebra-lemma-directed-colimit-exact}를 보라). 포함사상
$F_n(M) \to F_{n + 1}(M)$은 등급 $A$-가군의 사상으로서 분할되므로
$$
0 \to F_n(M) \otimes_A N \to F_{n + 1}(M) \otimes_A N \to
F_{n + 1}(M)/F_n(M) \otimes_A N \to 0
$$
은 미분 등급 $B$-가군들의 짧은 완전열이다. $N'$에 대해서도 유사한 열이
있고, $f$는 짧은 완전열들의 사상을 유도한다. 따라서 $n$에 대한
귀납법을 사용하면($F_{-1}(M) = 0$인 $n = -1$에서 시작한다), 사상
$F_{n + 1}(M)/F_n(M) \otimes_A N \to F_{n + 1}(M)/F_n(M) \otimes_A N'$
이 유사동형임을 보이면 충분하다는 결론을 얻는다.
$F_{n + 1}(M)/F_n(M)$는 $A$의 이동들의 직합이고
$f : N \to N'$가 유사동형이므로 $A[k]$에 대해서는 그 결과가 참이다.
\end{proof}

\begin{lemma}
\label{dga-lemma-compute-bc}
$R$을 환이라 하자. $(A, \text{d})$와 $(B, \text{d})$를 미분 등급
$R$-대수라 하자. $N$을 왼쪽 미분 등급 $A$-가군으로서 성질 (P)를 갖는
미분 등급 $(A, B)$-쌍가군이라 하자. 그러면 모든 미분 등급
$A$-가군 $M$에 대하여 $M \otimes_A^\mathbf{L} N$은
$M \otimes_A N$으로 계산된다.
\end{lemma}

\begin{proof}
$f : M \to M'$를 유사동형인 미분 등급 $A$-가군의 준동형이라 하자.
$f \otimes \text{id} : M \otimes_A N \to M' \otimes_A N$이
유사동형임을 주장한다. 이것이 참이면 Lemma \ref{dga-lemma-derived-bc}의
증명에서 유도 텐서곱을 구성한 방식에 의해 원하는 결과를 얻는다.
사상 $f \otimes \text{id}$의 구성은 $N$ 위의 왼쪽 미분 등급 $A$-가군
구조에만 의존한다. 더욱이 $N$이 미분 등급 $A^{opp}$-가군으로서 성질
(P)를 가지므로 $M \otimes_A N = N \otimes_{A^{opp}} M =
N \otimes_{A^{opp}}^\mathbf{L} M$이다. 따라서 주장은 Lemma
\ref{dga-lemma-functoriality-bc}에서 따른다.
\end{proof}

\begin{lemma}
\label{dga-lemma-tensor-hom-adjoint}
$R$을 환이라 하자. $(A, \text{d})$와 $(B, \text{d})$를 미분 등급
$R$-대수라 하자. $N$을 미분 등급 $(A, B)$-쌍가군이라 하자. 그러면
Lemma \ref{dga-lemma-derived-bc}의 함자
$$
- \otimes_A^\mathbf{L} N : D(A, \text{d}) \longrightarrow D(B, \text{d})
$$
는 Lemma \ref{dga-lemma-derived-restriction}의 함자
$$
R\Hom(N, -) : D(B, \text{d}) \longrightarrow D(A, \text{d})
$$
의 왼쪽 수반이다.
\end{lemma}

\begin{proof}
이는 Derived Categories, Lemma
\ref{derived-lemma-pre-derived-adjoint-functors-general}와,
Lemma \ref{dga-lemma-tensor-hom-adjunction}에 의해 $- \otimes_A N$과
$\Hom_{\text{Mod}^{dg}_{(B, \text{d})}}(N, -)$가 수반이라는 사실에서 따른다.
\end{proof}

\begin{example}
\label{dga-example-map-hom-tensor}
$R$을 환이라 하자. $(A, \text{d}) \to (B, \text{d})$를 미분 등급
$R$-대수의 준동형이라 하자. 그러면 $B$를 미분 등급 $(A, B)$-쌍가군으로
볼 수 있고, 다음 함자를 얻는다.
% P03-STACKS-E948: frozen display omits the derived L; retained literally.
$$
- \otimes_A B : D(A, \text{d}) \longrightarrow D(B, \text{d})
$$
Lemma \ref{dga-lemma-tensor-hom-adjoint}에 의해 이것의 왼쪽 수반은 함자
$R\Hom(B, -)$이다.
% P03-STACKS-E573: frozen source calls RHom the left rather than right adjoint.
미분 등급 $B$-가군 $N$에 대하여, $A \to B$를 통한 제한으로 얻는 미분
등급 $A$-가군을 $N_A$로 나타내자.
% P03-STACKS-E949: frozen English denote-construction omits “by”.
그러면 분명히 $B$-가군 $N$에 대하여 함자적인 표준 동형
$$
\Hom_{\text{Mod}^{dg}_{(B, \text{d})}}(B, N) \longrightarrow N_A,\quad
f \longmapsto f(1)
$$
이 있다. 따라서 $R\Hom(B, -)$가 제한 함자임을 알 수 있고,
$$
\Hom_{D(A, \text{d})}(M, N_A) =
\Hom_{D(B, \text{d})}(M \otimes^\mathbf{L}_A B, N)
$$
을 가환환의 경우와 꼭 마찬가지로 $M$과 $N$에 대하여 이중함자적으로
얻는다. 마지막으로 ${}_BB_A$가 $B$를 미분 등급 $(B, A)$-쌍가군으로 본
것일 때 $N_A = N \otimes_B {}_BB_A = N \otimes_B^\mathbf{L} {}_BB_A$이므로,
제한 역시 텐서 함자임을 주목하자.
\end{example}

\begin{lemma}
\label{dga-lemma-tensor-with-compact-fully-faithful}
Lemma \ref{dga-lemma-tensor-hom-adjoint}의 표기와 가정을 사용하고 다음을
가정하자.
% P03-STACKS-E950: frozen English opens with a prepositional sentence fragment.
\begin{enumerate}
\item $N$은 $D(B, \text{d})$의 콤팩트 객체를 정의하고,
\item 사상 $H^k(A) \to \Hom_{D(B, \text{d})}(N, N[k])$는 모든
$k \in \mathbf{Z}$에 대하여 동형이다.
\end{enumerate}
그러면 함자 $-\otimes_A^\mathbf{L} N$은 충실충만하다.
\end{lemma}

\begin{proof}
Lemma \ref{dga-lemma-tensor-hom-adjoint}에 의해 우리의 함자는
$R\Hom(N, -)$로 주어지는 왼쪽 수반을 갖는다.
% P03-STACKS-E574: frozen source calls RHom the left rather than right adjoint.
Categories, Lemma \ref{categories-lemma-adjoint-fully-faithful}에 의해,
미분 등급 $A$-가군 $M$에 대하여 사상
$$
M \longrightarrow R\Hom(N, M \otimes_A^\mathbf{L} N)
$$
이 $D(A, \text{d})$에서 동형임을 보이면 충분하다. 이를 위해
$$
H^n(M) \longrightarrow
\text{Ext}^n_{D(B, \text{d})}(N, M \otimes_A^\mathbf{L} N)
$$
이 동형임을 보이면 충분하다. Lemma
\ref{dga-lemma-derived-restriction-exts}를 보라. $N$은 콤팩트 객체이므로
우변은 직합과 가환한다. 따라서 Remark \ref{dga-remark-P-resolution}에 의해
$M = A[k]$일 때 이 사상이 동형임을 증명하면 충분하다. Remark
\ref{dga-remark-shift-tensor-no-sign}에 의해
$(A[k] \otimes_A^\mathbf{L} N) = N[k]$이므로, $N$에 대한 가정 (2)는 바로
이들에 대해 결과가 성립한다는 것이다.
\end{proof}

\begin{lemma}
\label{dga-lemma-base-change-K-flat}
$R \to R'$을 환 사상이라 하자. $(A, \text{d})$를 미분 등급
$R$-대수라 하고, $(A', \text{d})$를 기저변환, 즉
$A' = A \otimes_R R'$이라 하자. $A$가 $R$-가군 복합체로서 K-평탄이면
다음이 성립한다.
\begin{enumerate}
\item $- \otimes_A^\mathbf{L} A' : D(A, \text{d}) \to D(A', \text{d})$
은 다음 함자의 오른쪽 유도함자와 같다.
% P03-STACKS-E575: frozen source says right rather than left derived functor.
$$
K(A, \text{d}) \longrightarrow K(A', \text{d}),\quad
M \longmapsto M \otimes_R R'
$$
\item 도식
$$
\xymatrix{
D(A, \text{d}) \ar[r]_{- \otimes_A^\mathbf{L} A'} \ar[d]_{restriction} &
D(A', \text{d}) \ar[d]^{restriction} \\
D(R) \ar[r]^{- \otimes_R^\mathbf{L} R'} & D(R')
}
$$
은 가환하고,
\item $M$이 $R$-가군 복합체로서 K-평탄이면, 미분 등급 $A'$-가군
$M \otimes_R R'$은 $M \otimes_A^\mathbf{L} A'$을 나타낸다.
\end{enumerate}
\end{lemma}

\begin{proof}
임의의 미분 등급 $A$-가군 $M$에 대하여 표준 사상
$$
c_M : M \otimes_R R' \longrightarrow M \otimes_A A'
$$
이 있다. $P$를 성질 (P)를 갖는 미분 등급 $A$-가군이라 하자.
$c_P$가 동형이고 $P$가 $R$-가군 복합체로서 K-평탄임을 주장한다.
Lemma \ref{dga-lemma-derived-bc}의 유도 텐서곱 정의와 More on Algebra,
Section \ref{more-algebra-section-derived-tensor-product}를 사용하는
형식적 논증으로부터, 이는 보조정리에 서술된 모든 결과를 증명한다.

\medskip\noindent
$F_\bullet$을 $P$ 위의 필트레이션으로서, $P$가 성질 (P)를 가짐을
보이는 것이라 하자.
가정에 의해 $c_A$는 동형이고 $A$는 $R$-가군 복합체로서 K-평탄이다.
따라서 $A$의 이동들의 직합에 대해서도 마찬가지이다(원한다면 직합을
다루기 위해 More on Algebra, Lemma
\ref{more-algebra-lemma-colimit-K-flat}을 사용할 수 있다). 그러므로 복합체
$F_{p + 1}P/F_pP$에 대해서도 성립한다. 짧은 완전열
$$
0 \to F_pP \to F_{p + 1}P \to F_{p + 1}P/F_pP \to 0
$$
은 등급 가군의 열로서 분할 완전하므로, 귀납법으로 모든 $p$에 대하여
$c_{F_pP}$가 동형이고 $F_pP$가 $R$-가군 복합체로서 K-평탄임을 보일 수
있다(More on Algebra, Lemma
\ref{more-algebra-lemma-K-flat-two-out-of-three}을 사용하라). 마지막으로
$P = \colim F_pP$임을 사용하면, $c_P$가 동형이고 $P$가 $R$-가군
복합체로서 K-평탄이라는 결론을 얻는다(More on Algebra, Lemma
\ref{more-algebra-lemma-colimit-K-flat}을 사용하라).
\end{proof}

\begin{lemma}
\label{dga-lemma-RHom-is-tensor}
$R$을 환이라 하자. $(A, \text{d})$와 $(B, \text{d})$를 미분 등급
$R$-대수라 하자. $T$를 미분 등급 $(A, B)$-쌍가군이라 하자. 다음을
가정하자.
\begin{enumerate}
\item $T$는 $D(B, \text{d})$의 콤팩트 객체를 정의하고,
\item $S = \Hom_{\text{Mod}^{dg}_{(B, \text{d})}}(T, B)$
는 $D(A, \text{d})$에서 $R\Hom(T, B)$를 나타낸다.
\end{enumerate}
그러면 $S$는 미분 등급 $(B, A)$-쌍가군 구조를 가지며, 동형
$$
N \otimes_B^\mathbf{L} S \longrightarrow R\Hom(T, N)
$$
이 있고 이는 $D(B, \text{d})$ 안의 $N$에 대하여 함자적이다.
\end{lemma}

\begin{proof}
다음과 같이 쓰자: $\mathcal{B} = \text{Mod}^{dg}_{(B, \text{d})}$.
$S$ 위의 오른쪽 $A$-가군 구조는 사상
$A \to \Hom_\mathcal{B}(T, T)$과 Example
\ref{dga-example-dgm-dg-cat}에서 정의한 합성
$\Hom_\mathcal{B}(T, B) \otimes \Hom_\mathcal{B}(T, T)
\to \Hom_\mathcal{B}(T, B)$에서 온다. 이 곱셈을 한 번 더 사용하면 사상
$$
c_N :
N \otimes_B S = \Hom_\mathcal{B}(B, N) \otimes_B \Hom_\mathcal{B}(T, B)
\longrightarrow
\Hom_\mathcal{B}(T, N)
$$
이 있고 이는 $N$에 대하여 함자적이다. $N$이 주어지면 유사동형
$P \to N \to I$를 택할 수 있는데, $P$ 및 각각 $I$는 성질 (P) 및 각각
(I)를 갖는 미분 등급 $B$-가군이다. 그러면 $c_N$을 사용하여 사상
$P \otimes_B S \to \Hom_\mathcal{B}(T, I)$을 얻는데, 이는
$S \otimes_B^\mathbf{L} N$과 $R\Hom(T, N)$을 나타내는
객체들 사이의 사상이다.
% P03-STACKS-E576: frozen source reverses the factors in this represented tensor product.
이는 분명히 보조정리에서와 같은 함자 변환 $c$를 정의한다.

\medskip\noindent
$c$가 함자들의 동형임을 증명하기 위해, $N$이 성질 (P)를 갖는 미분 등급
$B$-가군이라고 가정해도 된다. $T$는 $D(B, \text{d})$의 콤팩트 객체를
정의하고 화살표의 양변은 삼각범주들의 완전 함자를 정의하므로, Lemma
\ref{dga-lemma-property-P-sequence}을 사용하여 $N$이 유한 필트레이션을 갖고
그 등급 조각들이 $B[k]$들의 직합인 경우로 환원한다. 화살표의 양변이
삼각범주들의 완전 함자라는 것과 $T$의 콤팩트성을 다시 사용하면
$N = B[k]$인 경우로 환원된다. 가정 (2)는 바로 이 경우 $c$가
동형이라는 가정이다.
\end{proof}






\section{유도 텐서곱의 합성}
\label{dga-section-compose-tensor-functors}

\noindent
독자에게 이 절을 건너뛰기를 권한다.

\medskip\noindent
$R$을 환이라 하자. $(A, \text{d})$, $(B, \text{d})$ 및
$(C, \text{d})$를 미분 등급 $R$-대수라 하자.
$N$을 미분 등급 $(A, B)$-쌍가군이라 하자.
% P03-STACKS-E951: frozen source twice calls N' a (B,C)-module rather
% than the bimodule required by both tensor actions and used later.
$N'$을 미분 등급 $(B, C)$-가군이라 하자.
% P03-STACKS-E952: frozen English omits “by” in both temporary denote clauses.
$A$-구조를 잊고 미분 등급 $B$-가군으로 본 쌍가군 $N$을
$N_B$로 쓰자. 그러면 표준 사상
\begin{equation}
\label{dga-equation-plain-versus-derived}
N_B \otimes_B^\mathbf{L} N'
\longrightarrow
(N \otimes_B N')_C
\end{equation}
이 $D(C, \text{d})$에 있다. 여기서 $(N \otimes_B N')_C$는
$(A, C)$-쌍가군 $N \otimes_B N'$을 미분 등급 $C$-가군으로
본 것을 뜻한다. 즉, 이 사상은 유도 텐서곱에서 보통 텐서곱으로
가는 사상이 항상 존재한다는 사실에서 온다(유도 텐서곱은 왼쪽
유도함자이기 때문이다).

\begin{lemma}
\label{dga-lemma-compose-tensor-functors-general}
$R$을 환이라 하자. $(A, \text{d})$, $(B, \text{d})$ 및
$(C, \text{d})$를 미분 등급 $R$-대수라 하자.
$N$을 미분 등급 $(A, B)$-쌍가군이라 하자.
% P03-STACKS-E951: the same frozen missing “bi” recurs in this lemma.
$N'$을 미분 등급 $(B, C)$-가군이라 하자.
(\ref{dga-equation-plain-versus-derived})가 동형이라고 가정하자.
그러면 합성
$$
\xymatrix{
D(A, \text{d}) \ar[rr]^{- \otimes_A^\mathbf{L} N} & &
D(B, \text{d}) \ar[rr]^{- \otimes_B^\mathbf{L} N'} & &
D(C, \text{d})
}
$$
은 $- \otimes_A^\mathbf{L} N''$과 동형이다. 여기서
$N'' = N \otimes_B N'$은 $(A, C)$-쌍가군으로 본다.
\end{lemma}

\begin{proof}
함자 변환
$$
(- \otimes_A^\mathbf{L} N) \otimes_B^\mathbf{L} N'
\longrightarrow
- \otimes_A^\mathbf{L} N''
$$
을 정의하자. 이를 위해 $M$을 성질 (P)를 갖는 미분 등급
$A$-가군이라 하자. Lemma \ref{dga-lemma-derived-bc}의 증명에서
함자 $- \otimes_A^\mathbf{L} N''$을 구성한 방식에 따르면, 보통
텐서곱 $M \otimes_A N''$은 $D(C, \text{d})$에서
$M \otimes_A^\mathbf{L} N''$을 표현한다. 또한
$$
M \otimes_A N'' =
M \otimes_A (N \otimes_B N') =
(M \otimes_A N) \otimes_B N'
$$
으로 쓸 수 있다. 가군 $M \otimes_A N$은 $D(B, \text{d})$에서
$M \otimes_A^\mathbf{L} N$을 표현한다. $Q$가 성질 (P)를 갖는
미분 등급 $B$-가군이 되도록 유사동형 $Q \to M \otimes_A N$을
택하자. 그러면 $Q \otimes_B N'$은 $D(C, \text{d})$에서
$(M \otimes_A^\mathbf{L} N) \otimes_B^\mathbf{L} N'$을
표현한다. 따라서 사상을
$$
(M \otimes_A^\mathbf{L} N) \otimes_B^\mathbf{L} N' =
Q \otimes_B N' \to
M \otimes_A N \otimes_B N' =
M \otimes_A^\mathbf{L} N''
$$
으로 정의할 수 있다. 이 사상의 구성은 $M$에 대해 함자적이고
특수 삼각형 및 직합과 양립한다. 세부사항은 생략한다.
이 사상이 동형이라는 것을 나타내는 $D(A, \text{d})$의 대상
$M$의 성질을 $T$라 하자. 그러면
\begin{enumerate}
\item 각 $M_i$에 대해 $T$가 성립하면 $\bigoplus M_i$에 대해서도
$T$가 성립하고,
% P03-STACKS-E953: frozen English conflates the two-out-of-three
% hypothesis with the closure rule; the canonical claim is retained.
\item 특수 삼각형에서 셋 중 $2$에 대해 $T$가 성립하면 $3$번째에도
성립하며,
\item $A[k]$에 대해 $T$가 성립한다. 실제로 이 경우에는 동형이라고
가정한 사상 (\ref{dga-equation-plain-versus-derived})의 시프트를 얻는다.
\end{enumerate}
따라서 Remark \ref{dga-remark-P-resolution}에 의해 성질 $T$는 항상
성립하며 증명이 끝난다.
\end{proof}

\noindent
$R$을 환이라 하자. $(A, \text{d})$, $(B, \text{d})$ 및
$(C, \text{d})$를 미분 등급 $R$-대수라 하자.
% P03-STACKS-E952: the second frozen temporary-denote clause also omits “by”.
미분 등급 대수 $A \otimes_R B$를 (오른쪽) 미분 등급 $B$-가군으로
본 것을 잠시 $(A \otimes_R B)_B$로 쓰고, 미분 등급 대수
$B \otimes_R C$를 미분 등급 $(B, C)$-쌍가군으로 본 것을
${}_B(B \otimes_R C)_C$로 쓰자. 그러면 표준 사상
\begin{equation}
\label{dga-equation-plain-versus-derived-algebras}
(A \otimes_R B)_B \otimes_B^\mathbf{L} {}_B(B \otimes_R C)_C
\longrightarrow
(A \otimes_R B \otimes_R C)_C
\end{equation}
이 $D(C, \text{d})$에 있다. 여기서 $(A \otimes_R B \otimes_R C)_C$는
미분 등급 $R$-대수 $A \otimes_R B \otimes_R C$를
(오른쪽) 미분 등급 $C$-가군으로 본 것을 뜻한다. 즉, 이 사상은
동일시
$$
(A \otimes_R B)_B \otimes_B {}_B(B \otimes_R C)_C =
(A \otimes_R B \otimes_R C)_C
$$
와 유도 텐서곱에서 보통 텐서곱으로 가는 사상이 항상 존재한다는
사실에서 온다(유도 텐서곱은 왼쪽 유도함자이기 때문이다).

\begin{lemma}
\label{dga-lemma-compose-tensor-functors-general-algebra}
$R$을 환이라 하자. $(A, \text{d})$, $(B, \text{d})$ 및
$(C, \text{d})$를 미분 등급 $R$-대수라 하자.
(\ref{dga-equation-plain-versus-derived-algebras})가 동형이라고 가정하자.
$N$을 미분 등급 $(A, B)$-쌍가군, $N'$을 미분 등급
$(B, C)$-쌍가군이라 하자. 그러면 합성
$$
\xymatrix{
D(A, \text{d}) \ar[rr]^{- \otimes_A^\mathbf{L} N} & &
D(B, \text{d}) \ar[rr]^{- \otimes_B^\mathbf{L} N'} & &
D(C, \text{d})
}
$$
은 증명에서 설명할 미분 등급 $(A, C)$-쌍가군 $N''$에 대한
$- \otimes_A^\mathbf{L} N''$과 동형이다.
\end{lemma}

\begin{proof}
Lemma \ref{dga-lemma-functoriality-bc}에 의해 $N$과 $N'$을 유사동형인
쌍가군으로 바꾸어도 된다. 따라서 $N$ 및 각각 $N'$이 미분 등급
$(A, B)$-쌍가군 및 각각 $(B, C)$-쌍가군으로서 성질 (P)를
갖는다고 가정해도 된다. Lemma \ref{dga-lemma-bimodule-resolve}를
보라. $(A, C)$-쌍가군 $N'' = N \otimes_B N'$에 대해 보조정리가
성립한다고 주장한다. 이를 증명하려면
$$
N_B \otimes_B^\mathbf{L} N' \longrightarrow (N \otimes_B N')_C
$$
이 $D(C, \text{d})$에서 동형임을 보이면 충분하다. Lemma
\ref{dga-lemma-compose-tensor-functors-general}을 보라.

\medskip\noindent
$F_\bullet$을 쌍가군에 대한 성질 (P)에서와 같은 $N$ 위의
필트레이션이라 하자. Lemma \ref{dga-lemma-bimodule-property-P-sequence}에
의해 미분 등급 $(A, B)$-쌍가군의 짧은 완전열
$$
0 \to
\bigoplus\nolimits F_iN \to
\bigoplus\nolimits F_iN \to N \to 0
$$
이 있고, 이는 등급 $(A, B)$-쌍가군의 열로서 분할된다. 특히 이는
미분 등급 $B$-가군의 허용 짧은 완전열이며, $D(B, \text{d})$에서
특수 삼각형
$$
\bigoplus\nolimits F_iN_B \to
\bigoplus\nolimits F_iN_B \to N_B \to
\bigoplus\nolimits F_iN_B[1]
$$
을 준다. $- \otimes_B^\mathbf{L} N'$이 삼각범주들의 완전 함자이고
직합과 가환한다는 것, 그리고 $- \otimes_B N'$이 허용 완전열을
허용 완전열로 보내며 직합과 가환한다는 것을 사용하면, 모든 $p$에
대해
$$
(F_pN)_B \otimes_B^\mathbf{L} N' \longrightarrow (F_pN)_B \otimes_B N'
$$
이 유사동형임을 보이는 것으로 환원된다. 등급 $(A, B)$-쌍가군으로서
분할되는 $(A, B)$-쌍가군의 짧은 완전열
$$
0 \to F_pN \to F_{p + 1}N \to F_{p + 1}N/F_pN \to 0
$$
을 사용해 논증을 반복하면 $F_{p + 1}N/F_pN$에 대한 같은 명제를
보이는 것으로 환원된다. 이 가군들은 $(A \otimes_R B)_B$의
시프트들의 직합이므로,
$$
(A \otimes_R B)_B \otimes_B^\mathbf{L} N'
\longrightarrow
(A \otimes_R B)_B \otimes_B N'
$$
이 유사동형임을 보이는 것으로 환원된다.

\medskip\noindent
$F_\bullet$을 쌍가군에 대한 성질 (P)에서와 같은 $N'$ 위의
필트레이션으로 택하자. $P$가 성질 (P)를 갖도록 미분 등급
$B$-가군의 유사동형 $P \to (A \otimes_R B)_B$를 택하자.
Lemma \ref{dga-lemma-derived-bc}의 구성에 의해 $P \otimes_B N'$은
$D(C, \text{d})$에서
$(A \otimes_R B)_B \otimes_B^\mathbf{L} N'$을 표현하므로,
$P \otimes_B N' \to (A \otimes_R B)_B \otimes_B N'$이
유사동형임을 보여야 한다. $N' = \colim F_pN'$이므로,
$P \otimes_B F_pN' \to (A \otimes_R B)_B \otimes_B F_pN'$이
유사동형임을 보이면 충분하다. 등급 $(B, C)$-쌍가군으로서 분할되는
짧은 완전열
$0 \to F_pN' \to F_{p + 1}N' \to F_{p + 1}N'/F_pN' \to 0$을
사용하면
$P \otimes_B F_{p + 1}N'/F_pN' \to
(A \otimes_R B)_B \otimes_B F_{p + 1}N'/F_pN'$
이 모든 $p$에 대해 유사동형임을 보이는 것으로 환원된다.
마지막으로 $F_{p + 1}N'/F_pN'$이
${}_B(B \otimes_R C)_C$의 시프트들의 직합이라는 것을 사용하면,
$$
P \otimes_B {}_B(B \otimes_R C)_C \to
(A \otimes_R B)_B \otimes_B {}_B(B \otimes_R C)_C
$$
이 유사동형임을 보이면 충분하다는 결론을 얻는다.
$P \to (A \otimes_R B)_B$는 성질 (P)를 만족하는 가군에 의한
분해이므로, 이 미분 등급 $C$-가군의 사상은 $D(C, \text{d})$에서
사상 (\ref{dga-equation-plain-versus-derived-algebras})을 표현한다.
이로써 증명이 끝난다.
\end{proof}

\begin{lemma}
\label{dga-lemma-compose-tensor-functors}
$R$을 환이라 하자. $(A, \text{d})$, $(B, \text{d})$ 및
$(C, \text{d})$를 미분 등급 $R$-대수라 하자.
$C$가 $R$-가군 복합체로서 K-평탄이면,
(\ref{dga-equation-plain-versus-derived-algebras})
는 동형이고 Lemma
\ref{dga-lemma-compose-tensor-functors-general-algebra}의 결론이 성립한다.
\end{lemma}

\begin{proof}
$P$가 성질 (P)를 갖도록 미분 등급 $B$-가군의 유사동형
$P \to (A \otimes_R B)_B$를 택하자. 그러면
$$
P \otimes_B (B \otimes_R C) \longrightarrow
(A \otimes_R B) \otimes_B (B \otimes_R C)
$$
이 유사동형임을 보여야 한다. 동치로 사상
$$
P \otimes_R C \longrightarrow
A \otimes_R B \otimes_R C
$$
을 살펴보는 것이다. More on Algebra, Lemma
\ref{more-algebra-lemma-K-flat-quasi-isomorphism}에 의해 $C$가
$R$-가군 복합체로서 K-평탄이면 이는 유사동형이다.
\end{proof}





\section{유도 텐서곱의 변형}
\label{dga-section-variant-base-change}

\noindent
$(\mathcal{C}, \mathcal{O})$를 환 달린 사이트라 하자. 그러면 함자들
$$
\text{Comp}(\mathcal{O}) \to K(\mathcal{O}) \to D(\mathcal{O})
$$
이 있고, 위에서 본 바와 같이 미분 등급 향상
$\text{Comp}^{dg}(\mathcal{O})$가 있다. 즉, 이는 아벨 범주
$\textit{Mod}(\mathcal{O})$에 결부된 Example
\ref{dga-example-category-complexes}의 미분 등급 범주이다.
$K^\bullet$을 $\mathcal{O}$-가군의 복합체, 다시 말해
$\text{Comp}^{dg}(\mathcal{O})$의 대상이라 하자. 다음과 같이 두자.
$$
(E, \text{d}) =
\Hom_{\text{Comp}^{dg}(\mathcal{O})}(K^\bullet, K^\bullet)
$$
이는 미분 등급 $\mathbf{Z}$-대수이다. 이 상황에도 유도 기저변환의
유사체가 있다고 주장한다.

\begin{lemma}
\label{dga-lemma-tensor-with-complex}
위 상황에는 미분 등급 범주들의 함자
$$
- \otimes_E K^\bullet :
\text{Mod}^{dg}_{(E, \text{d})}
\longrightarrow
\text{Comp}^{dg}(\mathcal{O})
$$
가 있다. 이 함자는 $E$를 $K^\bullet$로 보내고 직합과 가환한다.
\end{lemma}

\begin{proof}
$M$을 미분 등급 $E$-가군이라 하자. $\mathcal{C}$의 모든 대상
$U$에 대해 복합체 $K^\bullet(U)$는 왼쪽 미분 등급 $E$-가군인
동시에 오른쪽 $\mathcal{O}(U)$-가군이다. 두 작용은 가환하므로
쌍가군을 얻는다. 따라서 Sections \ref{dga-section-tensor-product}와
\ref{dga-section-bimodules}의 구성에 의해 텐서곱
$$
M \otimes_E K^\bullet(U)
$$
을 만들 수 있다. 이는 미분 등급 $\mathcal{O}(U)$-가군, 즉
$\mathcal{O}(U)$-가군의 복합체이다. 이 구성은 $U$에 대해
함자적이므로 층화하여 $\mathcal{O}$-가군의 복합체를 얻을 수 있고,
이를
$$
M \otimes_E K^\bullet
$$
로 쓴다. 더욱이 각 $U$에 대해 이 구성은 Lemma
\ref{dga-lemma-tensor}에 의해 미분 등급 범주들의 함자
$\text{Mod}^{dg}_{(E, \text{d})} \to \text{Comp}^{dg}(\mathcal{O}(U))$를
정한다. 따라서 보조정리에 서술된 함자를 얻음이 명백하다.
\end{proof}

\begin{lemma}
\label{dga-lemma-tensor-with-complex-homotopy}
Lemma \ref{dga-lemma-tensor-with-complex}의 함자는 삼각범주들의 완전 함자
$K(\text{Mod}_{(E, \text{d})}) \to K(\mathcal{O})$를 정한다.
\end{lemma}

\begin{proof}
Lemma \ref{dga-lemma-functorial}에 의해 이 함자는 호모토피 범주들 사이의
함자를 유도한다. $- \otimes_E K^\bullet$이 특수 삼각형을 특수
삼각형으로 보냄을 보여야 한다.
$0 \to K \to L \to M \to 0$이 미분 등급 $E$-가군의 허용 짧은
완전열이라고 가정하자.
% P03-STACKS-E577: frozen source calls the section M -> L a left inverse
% of L -> M twice, although its typed composition makes it a right inverse.
$s : M \to L$을 $L \to M$의 왼쪽 역인 등급 $E$-가군 준동형이라 하자.
그러면 $s$는 등급 $\mathcal{O}$-가군의 사상
$M \otimes_E K^\bullet \to L \otimes_E K^\bullet$을 정한다
(즉, $\mathcal{O}$-가군 구조와 등급은 보존하지만 미분은 보존하지
않을 수 있다). 이는
$L \otimes_E K^\bullet \to M \otimes_E K^\bullet$의 왼쪽 역이다.
따라서
$$
0 \to K \otimes_E K^\bullet \to L \otimes_E K^\bullet \to
M \otimes_E K^\bullet \to 0
$$
% P03-STACKS-E578: frozen English has number disagreement and a duplicated article.
은 항별로 분할되는 복합체의 짧은 완전열, 즉 $K(\mathcal{O})$에서
특수 삼각형을 정한다.
\end{proof}

\begin{lemma}
\label{dga-lemma-tensor-with-complex-derived}
Lemma \ref{dga-lemma-tensor-with-complex-homotopy}의 함자
$K(\text{Mod}_{(E, \text{d})}) \to K(\mathcal{O})$에는
$D(E, \text{d})$ 전체에 정의된 왼쪽 유도형이 있다. 이를
$- \otimes_E^\mathbf{L} K^\bullet : D(E, \text{d}) \to D(\mathcal{O})$로 쓴다.
\end{lemma}

\begin{proof}
이를 증명하기 위해 Derived Categories, Lemma
\ref{derived-lemma-find-existence-computes}를 사용한다. 대상들의 모음
$\mathcal{P}$로는 성질 (P)를 갖는 대상들을 사용한다.
성질 (1)은 Lemma \ref{dga-lemma-resolve}에서 증명되었다.
성질 (2)는 $s : P \to P'$가 성질 (P)를 갖는 가군들의 유사동형이면
Lemma \ref{dga-lemma-hom-derived}에 의해 $s$가 호모토피 동치이므로 성립한다.
\end{proof}

\begin{lemma}
\label{dga-lemma-upgrade-tensor-with-complex-derived}
$R$을 환이라 하자. $\mathcal{C}$를 사이트라 하고, $\mathcal{O}$를
가환 $R$-대수의 층이라 하자. $K^\bullet$을 $\mathcal{O}$-가군의
복합체라 하자. Lemma \ref{dga-lemma-tensor-with-complex-derived}의 함자는
다음 성질을 갖는다. $D(E, \text{d})$의 모든 $M$, $N$에 대해
표준 사상
$$
R\Hom(M, N)
\longrightarrow
R\Hom_\mathcal{O}(M \otimes_E^\mathbf{L} K^\bullet,
N \otimes_E^\mathbf{L} K^\bullet)
$$
이 $D(R)$에 있으며, 이는 코호몰로지 가군 위에서 함자
$- \otimes_E^\mathbf{L} K^\bullet$가 유도하는 사상
$$
\Ext^n_{D(E, \text{d})}(M, N) \to
\Ext^n_{D(\mathcal{O})}
(M \otimes_E^\mathbf{L} K^\bullet, N \otimes_E^\mathbf{L} K^\bullet)
$$
을 준다.
\end{lemma}

\begin{proof}
화살표의 오른쪽은 Cohomology on Sites, Section
\ref{sites-cohomology-section-global-RHom}에서 도입한 대역 유도 Hom이며,
이는 올바른 코호몰로지 가군들을 갖는다.
% P03-STACKS-E579: frozen source calls M an (R,A)-bimodule although
% this lemma places it in the E-module category.
왼쪽에 대해서는 $M$을 $(R, A)$-쌍가군으로 보고, Section
\ref{dga-section-restriction}에서 도입한 유도 $\Hom$을 사용한다. 이 역시
올바른 코호몰로지 가군들을 갖는다. 보조정리를 증명하기 위해 $M$과
$N$이 성질 (P)를 갖는 미분 등급 $E$-가군이라고 가정해도 된다.
Lemma \ref{dga-lemma-functoriality-derived-restriction}에 의해 이는 화살표의
왼쪽을 바꾸지 않는다. Lemma \ref{dga-lemma-compute-derived-restriction}에
의해 이때 화살표의 왼쪽은
% P03-STACKS-E580: both frozen Hom complexes below use B although the
% objects and the constructed DG functor are E-linear.
$\Hom_{\text{Mod}^{dg}_{(B, \text{d})}}(M, N)$
이 된다. Lemmas \ref{dga-lemma-tensor-with-complex},
\ref{dga-lemma-tensor-with-complex-homotopy} 및
\ref{dga-lemma-tensor-with-complex-derived}에서 미분 등급 범주들의 함자
$$
- \otimes_E K^\bullet :
\text{Mod}^{dg}_{(E, \text{d})}
\longrightarrow
\text{Comp}^{dg}(\mathcal{O})
$$
를 구성했고, $- \otimes_E^\mathbf{L} K^\bullet$이 이 함자를
성질 (P)를 갖는 미분 등급 $E$-가군에서 평가함으로써 계산됨을 보였다.
따라서 $R$-가군 복합체의 사상
$$
\Hom_{\text{Mod}^{dg}_{(B, \text{d})}}(M, N)
\longrightarrow
\Hom_{\text{Comp}^{dg}(\mathcal{O})}
(M \otimes_E K^\bullet, N \otimes_E K^\bullet)
$$
을 얻는다. $\mathcal{O}$-가군의 임의의 복합체
$\mathcal{F}^\bullet$, $\mathcal{G}^\bullet$에 대해 표준 사상
$$
\Hom_{\text{Comp}^{dg}(\mathcal{O})}
(\mathcal{F}^\bullet, \mathcal{G}^\bullet) =
\Gamma(\mathcal{C},
\SheafHom^\bullet(\mathcal{F}^\bullet, \mathcal{G}^\bullet))
\longrightarrow
R\Hom_\mathcal{O}(\mathcal{F}^\bullet, \mathcal{G}^\bullet).
$$
이 있다. 이 사상들을 합성하여 보조정리의 원하는 사상을 얻는다.
\end{proof}

\begin{lemma}
\label{dga-lemma-tensor-with-complex-hom-adjoint}
$(\mathcal{C}, \mathcal{O})$를 환 달린 사이트라 하자.
$K^\bullet$을 $\mathcal{O}$-가군의 복합체라 하자. 그러면
Lemma \ref{dga-lemma-tensor-with-complex-derived}의 함자
$$
- \otimes_E^\mathbf{L} K^\bullet :
D(E, \text{d})
\longrightarrow
D(\mathcal{O})
$$
는 Lemma \ref{dga-lemma-existence-of-derived}의 함자
$$
R\Hom(K^\bullet, -) : D(\mathcal{O}) \longrightarrow D(E, \text{d})
$$
의 왼쪽 수반이다.
\end{lemma}

\begin{proof}
이 명제는
$$
\Hom_{D(E, \text{d})}(M, R\Hom(K^\bullet, L^\bullet)) =
\Hom_{D(\mathcal{O})}(M \otimes^\mathbf{L}_E K^\bullet, L^\bullet)
$$
가 $M$과 $L^\bullet$에 대해 쌍함자적으로 성립한다는 뜻이다. 이를
보기 위해 $M$을 성질 (P)를 갖는 미분 등급 $E$-가군 $P$로 바꾸어도
된다. 또한 $L^\bullet$을 $\mathcal{O}$-가군의 K-단사 복합체
$I^\bullet$로 바꾸어도 된다. 명제에서 인용한 보조정리들이 주는
유도함자의 계산을 Lemma \ref{dga-lemma-hom-derived}와 결합하면 위 등식은
$$
\Hom_{K(\text{Mod}_{(E, \text{d})})}
(P, \Hom_\mathcal{B}(K^\bullet, I^\bullet)) =
\Hom_{K(\mathcal{O})}(P \otimes_E K^\bullet, I^\bullet)
$$
이 된다. 여기서 $\mathcal{B} = \text{Comp}^{dg}(\mathcal{O})$이다.
오른쪽에서 왼쪽으로 가며 $P$와 $I^\bullet$에 대해 함자적인 평가
사상이 있다(세부사항은 생략한다). $F_\bullet$을 성질 (P)의 정의에서와
같은 $P$ 위의 필트레이션으로 택하자. Lemma
\ref{dga-lemma-property-P-sequence}와 등식 양변이
$K(\text{Mod}_{(E, \text{d})})$ 위에서 $P$에 대한 호몰로지 함자라는
사실에 의해, $P$를 미분 등급 $E$-가군 $\bigoplus F_iP$로 바꾼
경우로 환원된다. 양변은 변수 $P$의 직합을 직곱으로 바꾸므로,
$P$가 미분 등급 $E$-가군들 $F_iP$ 중 하나인 경우로 환원된다.
각 $F_iP$에는 허용 단사상들로 주어지고 그 등급 조각들이 등급 사영
$E$-가군인 유한 필트레이션이 있으므로, $P$가 등급 사영 $E$-가군인
경우로 환원된다. 이 경우에는 명백히
$$
\Hom_{\text{Mod}^{dg}_{(E, \text{d})}}
(P, \Hom_\mathcal{B}(K^\bullet, I^\bullet)) =
\Hom_{\text{Comp}^{dg}(\mathcal{O})}(P \otimes_E K^\bullet, I^\bullet)
$$
가 등급 $\mathbf{Z}$-가군으로서 성립한다(이 명제는 명백한 경우
$P = E[k]$로 환원되기 때문이다). 이 동형은 미분과 양립하므로
결론을 얻는다.
\end{proof}

\begin{lemma}
\label{dga-lemma-fully-faithful-in-compact-case}
$(\mathcal{C}, \mathcal{O})$를 환 달린 사이트라 하자.
$K^\bullet$을 $\mathcal{O}$-가군의 복합체라 하자. 다음을 가정하자.
\begin{enumerate}
\item $K^\bullet$은 $D(\mathcal{O})$의 콤팩트 객체를 표현하고,
\item $E = \Hom_{\text{Comp}^{dg}(\mathcal{O})}(K^\bullet, K^\bullet)$
는 $D(\mathcal{O})$에서 $K^\bullet$의 Ext 군들을 계산한다.
\end{enumerate}
그러면 Lemma \ref{dga-lemma-tensor-with-complex-derived}의 함자
$$
- \otimes_E^\mathbf{L} K^\bullet :
D(E, \text{d})
\longrightarrow
D(\mathcal{O})
$$
는 충실충만하다.
\end{lemma}

\begin{proof}
% P03-STACKS-E581: frozen proof calls RHom the left adjoint although the
% cited lemma makes tensor left adjoint and RHom right adjoint.
Lemma \ref{dga-lemma-tensor-with-complex-hom-adjoint}에 의해 우리 함자는
$R\Hom(K^\bullet, -)$로 주어지는 왼쪽 수반을 가지므로, 미분 등급
$E$-가군 $M$에 대해 사상
$$
H^0(M) \longrightarrow
\Hom_{D(\mathcal{O})}(K^\bullet, M \otimes_E^\mathbf{L} K^\bullet)
$$
이 동형임을 보이면 충분하다. $M = P$가 성질 (P)를 갖는 미분 등급
$E$-가군이라고 가정해도 된다. $K^\bullet$이 콤팩트 객체를 정하므로,
Lemma \ref{dga-lemma-property-P-sequence}를 사용하여 $P$가 유한 필트레이션을
갖고 그 등급 조각들이 $E[k]$들의 직합인 경우로 환원된다. 다시
콤팩트성을 사용하면 $P = E[k]$인 경우로 환원된다. $K^\bullet$에 대한
가정은 바로 이 경우들에 결론이 성립한다는 것이다.
\end{proof}











\section{콤팩트 객체의 특징짓기}
\label{dga-section-compact}

\noindent
가법 범주의 콤팩트 객체는 Derived Categories, Definition
\ref{derived-definition-compact-object}에서 정의된다.
이 절에서는 미분 등급 대수의 유도범주에서 콤팩트 객체들을 특징짓는다.

\begin{remark}
\label{dga-remark-source-graded-projective}
$(A, \text{d})$를 미분 등급 대수라 하자. 모든 미분 등급
$A$-가군 $M$에 대해
$$
\Hom_{K(A, \text{d})}(P, M) =  \Hom_{D(A, \text{d})}(P, M)
$$
가 성립하는 미분 등급 $A$-가군 $P$들을 특징지을 수 있는가?
위 조건을 만족하는 대상 $P$들을 대상으로 하는 충만 부분범주를
$\mathcal{D} \subset K(A, \text{d})$라 하자. 그러면 $\mathcal{D}$는
$K(A, \text{d})$의 엄밀 충만 포화 삼각부분범주이다.
% P03-STACKS-E582: frozen English uses “where” for the membership question.
$P$가 등급 $A$-가군으로서 사영이면, $P$가 $\mathcal{D}$의 대상인지
확인하기 위해서는 $M$이 비순환일 때마다
$\Hom_{K(A, \text{d})}(P, M) = 0$임을 확인하면 충분하다.
그러나 일반적으로 $P$가 등급 $A$-가군으로서 사영이라고 가정하는
것만으로는 충분하지 않다. 예를 들어 $k$를 체라 하고
$A = R = k[\epsilon]$로 두자. 여기서
$k[\epsilon] = k[x]/(x^2)$는 쌍대수의 환이다.
$P$를 모든 $n \in \mathbf{Z}$에 대해 $P^n = R$이고 미분이
$\epsilon$에 의한 곱셈으로 주어지는 대상으로 하자. 그러면
$\text{id}_P \in \Hom_{K(A, \text{d})}(P, P)$는 영이 아닌 원소이지만
$P$는 비순환이다.
\end{remark}

\begin{remark}
\label{dga-remark-graded-projective-is-compact}
$(A, \text{d})$를 미분 등급 대수라 하자. 미분 등급 $A$-가군
$M$이 오른쪽 $A$-가군으로서 유한 개의 원소로 생성되면 $M$을
{\it 유한}이라 하자. $P$가 유한 등급 사영인 미분 등급 $A$-가군이면
$P$가 $D(A, \text{d})$의 콤팩트 객체를 주는가를 물을 수 있다.
아마도 일반적으로 참이 아니겠지만 우리는 반례를 알지 못한다.
그러나 $P$가 Remark \ref{dga-remark-source-graded-projective}의 범주
$\mathcal{D}$의 대상이기도 하면 이는 참이다(이는
$D(A, \text{d})$에서 직합이 가군들의 직합으로 주어진다는 사실에서
따른다. 세부사항은 생략한다).
\end{remark}

\begin{lemma}
\label{dga-lemma-factor-through-nicer}
$(A, \text{d})$를 미분 등급 대수라 하자. $E$를
$D(A, \text{d})$의 콤팩트 객체라 하자. $P$를 미분 등급
부분가군들에 의한 유한 필트레이션
$$
0 = F_{-1}P \subset F_0P \subset F_1P \subset \ldots \subset F_nP = P
$$
을 갖는 미분 등급 $A$-가군이라 하자. 어떤 집합들 $J_i$와 정수들
$k_{i, j}$에 대해 미분 등급 $A$-가군으로서
$$
F_{i + 1}P/F_iP \cong \bigoplus\nolimits_{j \in J_i} A[k_{i, j}]
$$
라고 하자. $E \to P$를 $D(A, \text{d})$의 사상이라 하자.
그러면 어떤 유한 부분집합들 $J'_i \subset J_i$에 대해
$F_{i + 1}P \cap P'/(F_iP \cap P')$이
$\bigoplus_{j \in J'_i} A[k_{i, j}]$와 같고, $E \to P$가
$P'$을 통해 분해되도록 하는 미분 등급 부분가군 $P' \subset P$가
존재한다.
\end{lemma}

\begin{proof}
$-1 \leq m \leq n$에 대한 귀납법으로 다음을 만족하는 미분 등급
부분가군 $P' \subset P$가 존재함을 증명한다.
\begin{enumerate}
\item $F_mP \subset P'$이다.
\item $i \geq m$에 대해 몫 $F_{i + 1}P \cap P'/(F_iP \cap P')$은
어떤 유한 부분집합 $J'_i \subset J_i$에 대해
$\bigoplus_{j \in J'_i} A[k_{i, j}]$와 동형이다.
\item $E \to P$는 $P'$을 통해 분해된다.
\end{enumerate}
기초 단계 $m = n$에서는 $P' = P$로 택할 수 있다.

\medskip\noindent
귀납 단계. $P'$이 $m$에 대해 조건을 만족한다고 가정하자.
$i \geq m$ 및 $j \in J'_i$에 대해
$x_{i, j} \in F_{i + 1}P \cap P'$를, $k_{i, j}$차의 동차 원소이면서
$F_{i + 1}P \cap P'/(F_iP \cap P')$에서의 상이
$j \in J_i$에 대응하는 직합인자의 생성원이 되도록 택하자.
$x_{i, j}$들은 $P'/F_mP$를 $A$-가군으로서 생성한다. 다음과 같이 쓰자.
$$
\text{d}(x_{i, j}) = \sum x_{i', j'} a_{i, j}^{i', j'} + y_{i, j}
$$
여기서 $y_{i, j} \in F_mP$이고 $a_{i, j}^{i', j'} \in A$이다.
각 $y_{i, j}$가 $F_mP/F_{m - 1}P$의 부분가군
$\bigoplus_{j \in J'_{m - 1}} A[k_{m - 1, j}]$의 원소로 사상되도록
하는 유한 부분집합 $J'_{m - 1} \subset J_{m - 1}$이 존재한다.
$P'' \subset F_mP$를 사상 $F_mP \to F_mP/F_{m - 1}P$에 의한
$\bigoplus_{j \in J'_{m - 1}} A[k_{m - 1, j}]$의 역상이라 하자.
그러면 $A$-부분가군
$$
P'' + \sum x_{i, j}A
$$
은 우리가 찾는 유형의 미분 등급 부분가군이다. 더욱이
$$
P'/(P'' + \sum x_{i, j}A) =
\bigoplus\nolimits_{j \in J_{m - 1} \setminus J'_{m - 1}} A[k_{m - 1, j}]
$$
이다. $E$가 콤팩트이므로 주어진 사상 $E \to P'$과 몫사상의 합성은
위에 표시한 가군의 유한 부분직합을 통해 분해된다. 따라서
$J'_{m - 1}$을 확대하면, 원하는 대로 $E \to P'$이
$P'' + \sum x_{i, j}A$를 통해 분해된다고 가정할 수 있다.
\end{proof}

\noindent
$D(A, \text{d})$의 모든 콤팩트 객체가 유한 등급 사영 미분 등급
$A$-가군에서 오는 것은 아니다. Examples, Section
\ref{examples-section-interesting-compact}를 보라.

\begin{proposition}
\label{dga-proposition-compact}
$(A, \text{d})$를 미분 등급 대수라 하고 $E$를 $D(A, \text{d})$의
대상이라 하자. 다음은 서로 동치이다.
\begin{enumerate}
\item $E$는 콤팩트 객체이다.
\item $E$는 $D(A, \text{d})$의 어떤 대상의 직합인자이고, 그 대상은
미분 등급 부분가군들에 의한 유한 필트레이션 $F_\bullet$을 가지며
각 $F_iP/F_{i - 1}P$가 $A$의 시프트들의 유한 직합인 미분 등급
가군 $P$로 표현된다.
\end{enumerate}
\end{proposition}

\begin{proof}
$E$가 콤팩트 객체라고 가정하자. Lemma \ref{dga-lemma-resolve}에 의해
$E$가 성질 (P)를 갖는 미분 등급 $A$-가군 $P$로 표현된다고
가정해도 된다. Lemma \ref{dga-lemma-property-P-sequence}의 허용 짧은
완전열에서 오는 특수 삼각형
$$
\bigoplus F_iP \to \bigoplus F_iP \to P
\xrightarrow{\delta} \bigoplus F_iP[1]
$$
을 생각하자. $E$가 콤팩트이므로 어떤
$\delta_i : P \to F_iP[1]$들에 대해
$\delta = \sum_{i = 1, \ldots, n} \delta_i$이다. $\delta$와 사상
$\bigoplus F_iP[1] \to \bigoplus F_iP[1]$의 합성은 영이므로
(Derived Categories, Lemma \ref{derived-lemma-composition-zero}),
$\delta = 0$이다(이는 $\bigoplus F_iP \to \bigoplus F_iP$가
직합인자 $F_iP$를 $\text{id}$와 $F_{i - 1}P$로 가는 포함사상의
차를 통해 보내는 데서 따른다). 따라서 $E$ 위의 항등사상은
$D(A, \text{d})$에서 $\bigoplus F_iP$를 통해 분해된다
(Derived Categories, Lemma \ref{derived-lemma-split}).
다음으로 $P$의 콤팩트성을 다시 사용하면, 사상
$E \to \bigoplus F_iP$는 어떤 $n$에 대해
$\bigoplus_{i = 1, \ldots, n} F_iP$를 통해 분해된다. 다시 말해
$E$ 위의 항등사상은 $\bigoplus_{i = 1, \ldots, n} F_iP$를 통해
분해된다. Lemma \ref{dga-lemma-factor-through-nicer}에 의해 $E$의
항등사상은 $E \to P \to E$로 분해되며, 여기서 $P$는 보조정리
명제의 (2)에 나오는 것과 같다. 따라서 (1)이 (2)를 함의함을 증명했다.

\medskip\noindent
(2)를 가정하자. Derived Categories, Lemma
\ref{derived-lemma-compact-objects-subcategory}에 의해 $P$가 콤팩트
객체를 줌을 보이면 충분하다. $P$는 성질 (P)를 가지므로
$$
\Hom_{D(A, \text{d})}(P, M) = \Hom_{K(A, \text{d})}(P, M)
$$
가 임의의 미분 등급 가군 $M$에 대해 성립한다(Lemma
\ref{dga-lemma-hom-derived}). $D(A, \text{d})$의 직합은 등급 가군들의
직합으로 주어지므로(Lemma \ref{dga-lemma-derived-products}),
$\Hom_{K(A, \text{d})}(P, M)$이 직합과 가환함을 보이는 것으로
환원된다. $K(A, \text{d})$가 삼각범주라는 것, $\Hom$이 첫째
변수에 대한 코호몰로지 함자라는 것, 그리고 $P$ 위의 필트레이션을
사용하면 $P$가 $A$의 시프트들의 유한 직합인 경우로 환원된다.
따라서 명백한 경우 $P = A[k]$로 환원된다.
\end{proof}

\begin{lemma}
\label{dga-lemma-compact-implies-bounded}
$(A, \text{d})$를 미분 등급 대수라 하자.
$D(A, \text{d})$의 모든 콤팩트 객체 $E$에 대해 다음 성질을 갖는
정수 $a \leq b$가 존재한다. $i \in [a, b]$에 대해
$H^i(M) = 0$이면 $\Hom_{D(A, \text{d})}(E, M) = 0$이다.
\end{lemma}

\begin{proof}
그러한 정수 쌍이 존재하는 $D(A, \text{d})$의 대상들의 모음은
$D(A, \text{d})$의 포화된 엄밀 충만 삼각부분범주임을 주목하자.
따라서 Proposition \ref{dga-proposition-compact}에 의해, $E$가 미분 등급
부분가군들에 의한 유한 필트레이션 $F_\bullet$을 가지며
$F_iP/F_{i - 1}P$가 $A$의 시프트들의 유한 직합인 미분 등급 가군
$P$로 표현되는 경우를 증명하면 충분하다. 삼각형과의 양립성을
사용하면 $P = A$인 경우를 증명하면 충분하다. 이 경우
$\Hom_{D(A, \text{d})}(A, M) = H^0(M)$이고 $a = b = 0$으로
결론이 성립한다.
\end{proof}

\noindent
$(A, \text{d})$가 미분이 영이고 $0$차에 놓인 대수이거나, 더
일반적으로 유한 개의 차수에만 놓여 있으면 콤팩트 객체들을 더
정밀하게 기술할 수 있다.

\begin{lemma}
\label{dga-lemma-compact}
$(A, \text{d})$를 미분 등급 대수라 하자. $|n| \gg 0$일 때
$A^n = 0$이라고 가정하자. $E$를 $D(A, \text{d})$의 대상이라 하자.
다음은 서로 동치이다.
\begin{enumerate}
\item $E$는 콤팩트 객체이다.
\item $E$는 등급 $A$-가군으로서 유한 사영이고, 모든 미분 등급
$A$-가군 $M$에 대해
$\Hom_{K(A, \text{d})}(P, M) = \Hom_{D(A, \text{d})}(P, M)$
을 만족하는 미분 등급 $A$-가군 $P$로 표현될 수 있다.
\end{enumerate}
\end{lemma}

\begin{proof}
$\mathcal{D} \subset K(A, \text{d})$를 Remark
\ref{dga-remark-source-graded-projective}에서 논한 삼각부분범주라 하자.
$P$를 $\mathcal{D}$의 대상이면서 등급 $A$-가군으로서 유한 사영인
대상이라 하자. 그러면 Remark \ref{dga-remark-graded-projective-is-compact}에
의해 $P$는 $D(A, \text{d})$의 콤팩트 객체를 표현한다.

\medskip\noindent
역을 증명하기 위해 $E$를 $D(A, \text{d})$의 콤팩트 객체라 하자.
Lemma \ref{dga-lemma-compact-implies-bounded}에서와 같은 $a \leq b$를
고정하자. 필요하다면 $a$를 줄이고 $b$를 늘려
$i \not \in [a, b]$에 대해 $H^i(E) = 0$이라고 가정해도 된다
(이는 Proposition \ref{dga-proposition-compact}와 $A$에 대한 가정에서
따른다). 더욱이 $|n| \geq c$이면 $A^n = 0$이 되도록 하는 정수
$c > 0$을 고정하자.

\medskip\noindent
Proposition \ref{dga-proposition-compact}에 의해 $E$는
$D(A, \text{d})$에서 다음 미분 등급 $A$-가군 $P$의 직합인자이다.
그 가군에는 미분 등급 부분가군들에 의한 유한 필트레이션
$F_\bullet$이 있고, $F_iP/F_{i - 1}P$들은 $A$의 시프트들의 유한
직합이다. 특히 $P$는 성질 (P)를 가지며, Lemma
\ref{dga-lemma-hom-derived}에 의해 임의의 미분 등급 가군 $M$에 대해
$\Hom_{D(A, \text{d})}(P, M) = \Hom_{K(A, \text{d})}(P, M)$이다.
다시 말해 $P$는 Remark \ref{dga-remark-source-graded-projective}에서
논한 삼각부분범주 $\mathcal{D} \subset K(A, \text{d})$의 대상이다.
$P$가 등급 $A$-가군으로서 유한 자유임을 주목하자.

\medskip\noindent
$b + 4c - n < a$가 되도록 $n > 0$을 택하자.
$E$ 위로의 사영자를 미분 등급 $A$-가군의 자기준동형
$\varphi : P \to P$로 표현하자. $K(A, \text{d})$에서 첫째 화살표의
원뿔을 $C$라 할 때의 특수 삼각형
$$
P \xrightarrow{1 - \varphi} P \to C \to P[1]
$$
을 생각하자. 그러면 $C$는 $\mathcal{D}$의 대상이고,
$D(A, \text{d})$에서 $C \cong E \oplus E[1]$이며, $C$는 유한
등급 자유 $A$-가군이다. 다음으로 $K(A, \text{d})$에서 특수 삼각형
$$
C[1] \to C \to C' \to C[2]
$$
을 생각하자. 여기서 $C'$은 $D(A, \text{d})$에서 합성
$$
C[1] \cong E[1] \oplus E[2] \to E[1] \to E \oplus E[1] \cong C
$$
을 표현하는 사상 $C[1] \to C$의 원뿔이다. 그러면 $C'$은
$E \oplus E[2]$를 표현한다. 이와 같이 계속하면 $\mathcal{D}$의
대상이고,
% P03-STACKS-E583: frozen English contains the spurious article in
% “is a finite free as a graded A-module”.
등급 $A$-가군으로서 유한 자유이며 $E \oplus E[n]$을 표현하는
미분 등급 $A$-가군 $P$를 찾을 수 있다.

\medskip\noindent
$P$의 $A$-가군으로서의 동차 원소 기저 $x_i$, $i \in I$를 택하고,
$d_i = \deg(x_i)$라 하자. $P_1$을 $d_i \leq a - c - 1$인
$x_i$와 $\text{d}(x_i)$로 생성되는 $P$의 $A$-부분가군이라 하자.
$P_2$를 $d_i \geq b - n + c$인 $x_i$와 $\text{d}(x_i)$로
생성되는 $P$의 $A$-부분가군이라 하자. 다음을 알 수 있다.
\begin{enumerate}
\item $P_1$과 $P_2$는 $P$의 미분 등급 부분가군이다.
\item $t \geq a$이면 $P_1^t = 0$이다.
\item $t \leq a - 2c$이면 $P_1^t = P^t$이다.
\item $t \leq b - n$이면 $P_2^t = 0$이다.
\item $t \geq b - n + 2c$이면 $P_2^t = P^t$이다.
\end{enumerate}
% P03-STACKS-E584: the frozen inequality below is reversed relative to
% b + 4c - n < a; it is preserved and the correction remains sidecar-only.
$n$의 선택에 의해 $b - n + 2c \geq a - 2c$이므로, 미분 등급
$A$-가군의 짧은 완전열
$$
0 \to P_1 \cap P_2 \to P_1 \oplus P_2 \xrightarrow{\pi} P \to 0
$$
을 얻는다. $P$는 등급 $A$-가군으로서 사영이므로 이는 허용 짧은
완전열이다(Lemma \ref{dga-lemma-target-graded-projective}). 따라서
$K(A, \text{d})$에서 경계사상
$\delta : P \to (P_1 \cap P_2)[1]$을 얻는다. Lemma
\ref{dga-lemma-admissible-ses}를 보라. $P = E \oplus E[n]$이고
$P_1 \cap P_2$가 차수 $(b - n, a)$에 놓이므로
$\Hom_{D(A, \text{d})}(E \oplus E[n], (P_1 \cap P_2)[1])$은
영이다. 따라서 $P$가 $\mathcal{D}$의 대상이므로 $\delta = 0$은
$K(A, \text{d})$의 사상으로서 영이다. Derived Categories, Lemma
\ref{derived-lemma-split}에 의해 어떤 $-1$차 사상 $h : P \to P$에
대해 $\pi \circ s = \text{id}_P + \text{d}h + h\text{d}$가
되도록 하는 사상 $s : P \to P_1 \oplus P_2$를 찾을 수 있다.
$P_1 \oplus P_2 \to P$가 전사이고 $P$가 등급 $A$-가군으로서
사영이므로 $h$의 동차 리프트 $\tilde h : P \to P_1 \oplus P_2$를
택할 수 있다. 이제 $s$를 $s + \text{d} \tilde h + \tilde h \text{d}$로
바꾸면 $\pi \circ s = \text{id}_P$를 얻는다. 즉, 직합분해
$P = s^{-1}(P_1) \oplus s^{-1}(P_2)$를 얻는다.
$s^{-1}(P_2)$는 차수 $\geq b - n + 2c$에서 $P$와 같으므로
$s^{-1}(P_2) \to P \to E$는 유사동형, 즉
$D(A, \text{d})$의 동형이다. 이로써 증명이 끝난다.
\end{proof}








\section{유도범주의 동치}
\label{dga-section-equivalence}

\noindent
$R$을 환이라 하자. $(A, \text{d})$와 $(B, \text{d})$를 미분 등급
$R$-대수라 하자. 자연스럽게 생기는 질문은 $D(A, \text{d})$와
$D(B, \text{d})$가 $R$-선형 삼각범주로서 동치라는 것이 무엇을
뜻하는가이다. 이는 상당히 미묘한 질문이며, 언제나 올바른 질문은
아님이 드러날 것이다.
%% P03-STACKS-E585
그럼에도 이 절에서는 이를 보장하는 조건을 몇 가지 모은다.

\medskip\noindent
독자에게 이 주제에 관한 획기적인 논문 \cite{Rickard}를 반드시
살펴보기를 강력히 권한다.

\begin{lemma}
\label{dga-lemma-qis-equivalence}
$R$을 환이라 하자. $(A, \text{d}) \to (B, \text{d})$를
코호몰로지 대수의 동형을 유도하는 미분 등급 $R$-대수의
준동형이라 하자. 그러면
$$
- \otimes_A^\mathbf{L} B : D(A, \text{d}) \to D(B, \text{d})
$$
는 삼각범주의 $R$-선형 동치를 주며, 그 준역함자는 제한 함자
$N \mapsto N_A$이다.
\end{lemma}

\begin{proof}
Lemma \ref{dga-lemma-tensor-with-compact-fully-faithful}에 의해
함자 $M \longmapsto M \otimes_A^\mathbf{L} B$는 충실충만이다.
Lemma \ref{dga-lemma-tensor-hom-adjoint}에 의해 함자
$N \longmapsto R\Hom(B, N) = N_A$는 오른쪽 수반이다.
Example \ref{dga-example-map-hom-tensor}를 보라.
$R\Hom(B, -)$의 핵이 영임은 자명하다.
따라서 결과는 Derived Categories, Lemma
\ref{derived-lemma-fully-faithful-adjoint-kernel-zero}에서 따른다.
\end{proof}

\noindent
위 증명을 분석하면 다음 일반화도 즉시 얻어진다.

\begin{lemma}
\label{dga-lemma-tilting-equivalence}
$R$을 환이라 하자. $(A, \text{d})$와 $(B, \text{d})$를
미분 등급 $R$-대수라 하자. $N$을 미분 등급
$(A, B)$-쌍가군이라 하자. 다음을 가정하자.
\begin{enumerate}
\item $N$은 $D(B, \text{d})$의 콤팩트 객체를 정의한다.
\item $N' \in D(B, \text{d})$이고 모든 $n \in \mathbf{Z}$에 대해
$\Hom_{D(B, \text{d})}(N, N'[n]) = 0$이면 $N' = 0$이다.
\item 모든 $k \in \mathbf{Z}$에 대해 사상
$H^k(A) \to \Hom_{D(B, \text{d})}(N, N[k])$는 동형이다.
\end{enumerate}
그러면
$$
- \otimes_A^\mathbf{L} N : D(A, \text{d}) \to D(B, \text{d})
$$
는 삼각범주의 $R$-선형 동치를 준다.
\end{lemma}

\begin{proof}
Lemma \ref{dga-lemma-tensor-with-compact-fully-faithful}에 의해
함자 $M \longmapsto M \otimes_A^\mathbf{L} N$은 충실충만이다.
Lemma \ref{dga-lemma-tensor-hom-adjoint}에 의해 함자
$N' \longmapsto R\Hom(N, N')$는 오른쪽 수반이다.
%% P03-STACKS-E586
가정 (3)에 의해 $R\Hom(N, -)$의 핵은 영이다.
따라서 결과는 Derived Categories, Lemma
\ref{derived-lemma-fully-faithful-adjoint-kernel-zero}에서 따른다.
\end{proof}

\begin{remark}
\label{dga-remark-tilting-equivalence}
Lemma \ref{dga-lemma-tilting-equivalence}에서 조건 (2)를 $N$이
$D_{compact}(B, d)$의 고전적 생성자라는 조건으로 바꿀 수 있다.
Derived Categories, Proposition
\ref{derived-proposition-generator-versus-classical-generator}을 보라.
또한 $R\Hom(N, B)$가 $D(A, \text{d})$의 콤팩트 객체임을
알고 있다면, $N$이 $D_{compact}(B, \text{d})$의 약한 생성자임을
확인하는 것으로 충분하다. 증명은 생략한다. Stacks project에서
이 결과가 필요하게 되면 여기에 증명을 추가할 것이다.
\end{remark}

\noindent
아래 보조정리의 $B$-가군 $P$를 때때로
“$(A, B)$-틸팅 복합체”라고 부른다.

\begin{lemma}
\label{dga-lemma-rickard}
$R$을 환이라 하자. $(A, \text{d})$와 $(B, \text{d})$를
미분 등급 $R$-대수라 하자. $A = H^0(A)$라고 가정하자.
다음은 서로 동치이다.
\begin{enumerate}
\item $D(A, \text{d})$와 $D(B, \text{d})$는 $R$-선형
삼각범주로서 동치이다.
\item 다음을 만족하는 $D(B, \text{d})$의 객체 $P$가 존재한다.
\begin{enumerate}
\item $P$는 $D(B, \text{d})$의 콤팩트 객체이다.
\item $N \in D(B, \text{d})$이고 모든 $i \in \mathbf{Z}$에 대해
$\Hom_{D(B, \text{d})}(P, N[i]) = 0$이면 $N = 0$이다.
\item $i \not = 0$이면 $\Hom_{D(B, \text{d})}(P, P[i]) = 0$이고,
$i = 0$이면 $A$와 같다.
\end{enumerate}
\end{enumerate}
(2)에서 구성한 동치 $D(A, \text{d}) \to D(B, \text{d})$는
$A$를 $P$로 보낸다.
\end{lemma}

\begin{proof}
$F : D(A, \text{d}) \to D(B, \text{d})$를 동치라 하자.
그러면 $F$는 콤팩트 객체를 콤팩트 객체로 보낸다. 따라서
$P = F(A)$는 콤팩트하므로 (2)(a)가 성립한다.
조건 (2)(b)와 (2)(c)는 $F$가 동치라는 사실에서 즉시 따른다.

\medskip\noindent
$P$를 (2)의 객체라 하자. $P$를 성질 (P)를 갖는 미분 등급
가군으로 나타내자. 다음과 같이 놓는다.
$$
(E, \text{d}) = \Hom_{\text{Mod}^{dg}_{(B, \text{d})}}(P, P)
$$
그러면 Lemma \ref{dga-lemma-hom-derived}와 가정 (2)(c)에 의해
$H^0(E) = A$이고 $k \not = 0$이면 $H^k(E) = 0$이다.
$P$를 $(E, B)$-쌍가군으로 보고 Lemma
\ref{dga-lemma-tilting-equivalence}와 가정 (2)(b)를 사용하면,
$E$를 $P$로 보내는 동치
$$
D(E, \text{d}) \to D(B, \text{d})
$$
를 얻는다. $E' \subset E$를 다음을 만족하는 미분 등급
$R$-부분대수라 하자.
$$
(E')^i = \left\{
\begin{matrix}
E^i & \text{if }i < 0 \\
\Ker(E^0 \to E^1) & \text{if }i = 0 \\
0 & \text{if }i > 0
\end{matrix}
\right.
$$
그러면 미분 등급 대수의 유사동형
$(A, \text{d}) \leftarrow (E', \text{d}) \rightarrow (E, \text{d})$가
존재한다. 따라서 Lemma \ref{dga-lemma-qis-equivalence}에 의해 동치들
$$
D(A, \text{d}) \leftarrow D(E', \text{d}) \rightarrow D(E, \text{d})
\rightarrow D(B, \text{d})
$$
을 얻는다.
\end{proof}

\begin{remark}
\label{dga-remark-lift-equivalence-to-dga}
$R$을 환이라 하자. $(A, \text{d})$와 $(B, \text{d})$를 미분 등급
$R$-대수라 하자. 삼각범주의 $R$-선형 동치
$$
F : D(A, \text{d}) \longrightarrow D(B, \text{d})
$$
가 주어졌다고 하자. $N = F(A)$로 놓자. 그러면 $N$은 미분 등급
$B$-가군이다. $F$가 동치이고 $A$가 $D(A, \text{d})$의 콤팩트
객체이므로, $N$은 $D(B, \text{d})$의 콤팩트 객체이다.
$A$가 $D(A, \text{d})$를 생성하고 $F$가 동치이므로,
$N$은 $D(B, \text{d})$를 생성한다.
%% P03-STACKS-E955
마지막으로 $H^k(A) = \Hom_{D(A, \text{d})}(A, A[k])$이고
$F$가 동치이므로, 모든 $k$에 대해 $F$는 동형
$H^k(A) = \Hom_{D(B, \text{d})}(N, N[k])$를 유도한다.
Lemma \ref{dga-lemma-tilting-equivalence}의 구성에서 생기는 동치
$D(A, \text{d}) \longrightarrow D(B, \text{d})$가 존재한다고
결론내리기 위해 필요한 것은 $N$ 위의 왼쪽 $A$-가군 구조가
미분과 양립하고 주어진 오른쪽 $B$-가군 구조와 가환한다는 것뿐이다.
사실 $N$을 유사동형인 미분 등급 $B$-가군으로 바꾼 뒤 이를
구성하는 것으로 충분하다. 이 가군 구조는 특정한 경우에 구성할
수 있다. 예를 들어 $F$가 결합된 미분 등급 범주들 사이의
미분 등급 함자
$$
F^{dg} :
\text{Mod}^{dg}_{(A, \text{d})}
\longrightarrow
\text{Mod}^{dg}_{(B, \text{d})}
$$
로 올라간다고 가정하면(표기는 Example \ref{dga-example-dgm-dg-cat} 참조),
이 조건이 성립한다. 다른 경우는 아래 명제에서 다룬다.
\end{remark}

\begin{proposition}
\label{dga-proposition-rickard}
$R$을 환이라 하자. $(A, \text{d})$와 $(B, \text{d})$를
미분 등급 $R$-대수라 하자.
$F : D(A, \text{d}) \to D(B, \text{d})$를 삼각범주의
$R$-선형 동치라 하자. 다음을 가정하자.
\begin{enumerate}
\item $A = H^0(A)$이다.
\item $B$는 $R$-가군의 복합체로서 K-평탄이다.
\end{enumerate}
그러면 Lemma \ref{dga-lemma-tilting-equivalence}에서와 같은
$(A, B)$-쌍가군 $N$이 존재한다.
\end{proposition}

\begin{proof}
위 Remark \ref{dga-remark-lift-equivalence-to-dga}에서와 같이
$D(B, \text{d})$에서 $N = F(A)$로 놓는다. $N$이 성질 (P)를 갖는
미분 등급 $B$-가군이라고 가정해도 된다. 다음과 같이 놓는다.
$$
(E, \text{d}) = \Hom_{\text{Mod}^{dg}_{(B, \text{d})}}(N, N)
$$
그러면 Lemma \ref{dga-lemma-hom-derived}에 의해 $H^0(E) = A$이고
$k \not = 0$이면 $H^k(E) = 0$이다. 또한 Remark
\ref{dga-remark-lift-equivalence-to-dga}의 논의와 Lemma
\ref{dga-lemma-tilting-equivalence}에 의해, $(E, B)$-쌍가군으로서의
$N$은 동치
$- \otimes_E^\mathbf{L} N : D(E, \text{d}) \to D(B, \text{d})$를
유도한다. $E' \subset E$를 다음을 만족하는 미분 등급
$R$-부분대수라 하자.
$$
(E')^i = \left\{
\begin{matrix}
E^i & \text{if }i < 0 \\
\Ker(E^0 \to E^1) & \text{if }i = 0 \\
0 & \text{if }i > 0
\end{matrix}
\right.
$$
그러면 미분 등급 대수의 유사동형
$(A, \text{d}) \leftarrow (E', \text{d}) \rightarrow (E, \text{d})$가
존재한다. 따라서 Lemma \ref{dga-lemma-qis-equivalence}에 의해 동치들
$$
D(A, \text{d}) \leftarrow D(E', \text{d}) \rightarrow D(E, \text{d})
\rightarrow D(B, \text{d})
$$
을 얻는다.
%% P03-STACKS-E956
왼쪽 수직 화살표의 준역함자 $D(A, \text{d}) \to D(E', \text{d})$는,
$A$를 $(A, E')$-쌍가군으로 보았을 때
$M \mapsto M \otimes_A^\mathbf{L} A$로 주어진다.
Example \ref{dga-example-map-hom-tensor}를 보라. 한편 함자
$D(E', \text{d}) \to D(B, \text{d})$는 위의 $N$을 사용하여
$M \mapsto M \otimes_{E'}^\mathbf{L} N$으로 주어진다.
Lemma \ref{dga-lemma-compose-tensor-functors}에 의해 결론이 따른다.
\end{proof}

\begin{remark}
\label{dga-remark-rickard}
$A, B, F, N$을 Proposition \ref{dga-proposition-rickard}에서와 같이
두자. $F$와 함자 $G(-) = - \otimes_A^\mathbf{L} N$이
동형인지는 분명하지 않다. 구성에 의해 $D(B, \text{d})$ 안의 동형
$N = G(A) \to F(A)$가 존재한다.
%% P03-STACKS-E957
$M$이 등급 사영인 미분 등급 $A$-가군인 경우(예를 들어 $A$의
시프트들의 합인 경우), 이를 함자적 동형 $G(M) \to F(M)$으로
확장하는 일은 어렵지 않다. 이어서 $M$이 그러한 가군들 사이 사상의
원뿔이면 $G(M) \cong F(M)$라고 결론낼 수 있다.
일반적으로 더 강한 사실이 성립하는지는 알지 못한다.
\end{remark}

\begin{lemma}
\label{dga-lemma-rickard-rings}
$R$을 환이라 하자. $A$와 $B$를 $R$-대수라 하자.
다음은 서로 동치이다.
\begin{enumerate}
\item 삼각범주의 $R$-선형 동치 $D(A) \to D(B)$가 존재한다.
\item 다음을 만족하는 $D(B)$의 객체 $P$가 존재한다.
\begin{enumerate}
\item $P$는 유한 사영 $B$-가군들의 유한 복합체로 나타낼 수 있다.
\item $K \in D(B)$이고 모든 $i \in \mathbf{Z}$에 대해
$\Ext^i_B(P, K) = 0$이면 $K = 0$이다.
\item $i \not = 0$이면 $\Ext^i_B(P, P) = 0$이고,
$i= 0$이면 $A$와 같다.
\end{enumerate}
\end{enumerate}
더 나아가 $B$가 $R$-가군으로서 평탄하면, 이는 다음 조건과도
동치이다.
\begin{enumerate}
\item[(3)] $- \otimes_A^\mathbf{L} N : D(A) \to D(B)$가
동치가 되게 하는 $(A, B)$-쌍가군 $N$이 존재한다.
\end{enumerate}
\end{lemma}

\begin{proof}
(1)과 (2)의 동치는 Lemma \ref{dga-lemma-rickard}의 특수한 경우와
$D(B)$의 콤팩트 객체를 특징짓는 Lemma \ref{dga-lemma-compact}의
결과를 결합하면 따른다(작은 세부사항은 생략한다).
$B$가 $R$-평탄인 경우 (3)과의 동치는 Proposition
\ref{dga-proposition-rickard}에서 따른다.
\end{proof}

\begin{remark}
\label{dga-remark-centers}
$R$을 환이라 하자. $A$와 $B$를 $R$-대수라 하자.
$D(A)$와 $D(B)$가 $R$-선형 삼각범주로서 동치이면,
$A$와 $B$의 중심들은 $R$-대수로서 동형이다.
특히 $A$와 $B$가 가환이면 $A \cong B$이다.
상당히 까다로운 증명은 \cite[Proposition 9.2]{Rickard} 또는
\cite[Proposition 6.3.2]{KZ}에서 찾을 수 있다.
다른 접근으로는 호흐실트 코호몰로지를 사용할 수도 있다
(아래 주석 참조).
\end{remark}

\begin{remark}
\label{dga-remark-hochschild-cohomology}
$R$을 환이라 하자. $(A, \text{d})$와 $(B, \text{d})$를
유도 동치인 미분 등급 $R$-대수라 하자. 즉, 삼각범주의
$R$-선형 동치 $D(A, \text{d}) \to D(B, \text{d})$가 존재한다고
하자. 우리는 $(A, \text{d})$와 $(B, \text{d})$의 특정 불변량들이
일치함을 보이고자 한다. 많은 상황에서는 이 사정을 더 잘 통제할
수 있다. 예를 들어 어떤 미분 등급 $(A, B)$-쌍가군 $\Omega$에 대해
다음 꼴의 동치가 존재할 수 있다.
%% P03-STACKS-E587
$$
- \otimes_A \Omega : D(A, \text{d}) \longrightarrow D(B, \text{d})
$$
(Proposition \ref{dga-proposition-rickard}의 상황에서 이런 일이 일어나며,
동치가 기하적 구성에서 나오는 경우에도 흔히 그러하다.)
%% P03-STACKS-E958
또한 우리 함자의 준역함자가 어떤 미분 등급
$(B, A)$-쌍가군 $\Omega'$에 대해
$$
- \otimes_B^\mathbf{L} \Omega' : D(B, \text{d}) \longrightarrow D(A, \text{d})
$$
로 주어진다고 하자(앞에서와 마찬가지로 실제 상황에서 그러한
가군 $\Omega'$이 존재하는 경우가 많다). 이 경우 쌍가군의
유도범주 위의 함자
$$
D(A^{opp} \otimes_R A, \text{d})
\longrightarrow
D(B^{opp} \otimes_R B, \text{d}),\quad
M \longmapsto \Omega' \otimes^\mathbf{L}_A M \otimes_A^\mathbf{L} \Omega
$$
를 생각할 수 있다(쌍가군을 오른쪽 가군으로 바꾸려면
Lemma \ref{dga-lemma-bimodule-over-tensor}를 사용하라).
이 함자는 $(A, A)$-쌍가군 $A$를 $(B, B)$-쌍가군 $B$로 보낸다.
적절한 조건(예를 들어 $R$ 위에서 $A$, $B$, $\Omega$의 평탄성 등)
아래에서는 이 함자도 동치가 된다. 이 경우 호흐실트 코호몰로지
군들의 동형
$$
\begin{aligned}
HH^i(A, \text{d}) =
&\Hom_{D(A^{opp} \otimes_R A, \text{d})}(A, A[i]) \\
&
\longrightarrow
\Hom_{D(B^{opp} \otimes_R B, \text{d})}(B, B[i]) \\
&=
HH^i(B, \text{d}).
\end{aligned}
$$
을 얻는다. 예를 들어 $A = H^0(A)$이면 $HH^0(A, \text{d})$는
$A$의 중심과 같고, 이는 Remark \ref{dga-remark-centers}에서 언급한
결과의 개념적 증명을 준다. 이 주석이 필요하게 되면 여기에 정확한
명제와 자세한 증명을 제시할 것이다.
\end{remark}


\section{미분 등급 대수의 분해}
\label{dga-section-resolution-dgas}

\noindent
$R$을 환이라 하자. 우리의 가정 아래 집합 $S$ 위의 자유
$R$-대수 $R\langle S \rangle$는 다음 식들을 $R$-기저로 갖는
대수이다.
$$
s_1 s_2 \ldots s_n
$$
%% P03-STACKS-E959
여기서 $n \geq 0$이고 $s_1, \ldots, s_n \in S$는 $S$의
원소들로 이루어진 열이다. 곱셈은 이어붙이기로 주어진다.
$$
(s_1 s_2 \ldots s_n) \cdot (s'_1 s'_2 \ldots s'_m) =
s_1 \ldots s_n s'_1 \ldots s'_m
$$
이 대수는 모든 $R$-대수 $A$에 대해 사상
$$
\Mor_{R\text{-alg}}(R\langle S \rangle, A) \to
\text{Map}(S, A),\quad
\varphi \longmapsto (s \mapsto \varphi(s))
$$
이 전단사라는 성질로 특징지어진다.

\medskip\noindent
등급 $R$-대수의 범주에서는 집합 $S$에 등급이 주어져야 하며,
이를 사상 $\deg : S \to \mathbf{Z}$로 생각한다. 그러면
$R\langle S\rangle$는 단항식의 차수가
$$
\deg(s_1 s_2 \ldots s_n) = \deg(s_1) + \ldots + \deg(s_n)
$$
이 되도록 하는 등급을 갖는다.
%% P03-STACKS-E960
이 상황에서는 모든 등급 $R$-대수 $A$에 대해 사상
$$
\Mor_{\text{graded }R\text{-alg}}(R\langle S \rangle, A) \to
\text{Map}_{\text{graded sets}}(S, A),\quad
\varphi \longmapsto (s \mapsto \varphi(s))
$$
이 전단사이다.

\medskip\noindent
$A$가 등급 $R$-대수이고 $S$가 등급 집합이면, 마찬가지로
$A\langle S \rangle$를 만들 수 있다. $A\langle S \rangle$의
원소들은 다음 꼴 원소들의 합이다.
$$
a_0 s_1 a_1 s_2 \ldots a_{n - 1} s_n a_n
$$
여기서 $a_i \in A$이고, 이 식들이 $(a_0, \ldots, a_n)$에 대해
$R$-다중선형이라는 관계들로 나눈다. 따라서 $S$의 원소들로
이루어진 각 열 $s_1, \ldots, s_n$에 대해 포함
$$
A \otimes_R \ldots \otimes_R A \subset A\langle S \rangle
$$
이 존재하며, 이 대수는 이들의 직합이다. 이 정의로부터 독자는
$\varphi$를 $(\varphi|_A, (s \mapsto \varphi(s)))$로 보내는 사상
$$
\Mor_{\text{graded }R\text{-alg}}(A\langle S \rangle, B) \to
\Mor_{\text{graded }R\text{-alg}}(A, B) \times
\text{Map}_{\text{graded sets}}(S, B),
$$
%% P03-STACKS-E588
이 사상은 모든 등급 $R$-대수 $A$에 대해 전단사임을 보일 수 있다.
$A$가 자유 등급 $R$-대수였다면 $A\langle S \rangle$도
그러함을 관찰한다.

\medskip\noindent
$A$가 미분 등급 $R$-대수이고 $S$가 등급 집합이라고 하자.
또한 각 $s \in S$에 대해 동차 원소 $f_s \in A$가 주어져 있고,
$\deg(f_s) = \deg(s) + 1$이며 $\text{d}f_s = 0$이라고 하자.
그러면 $\text{d}(s) = f_s$가 되게 하는
$A\langle S \rangle$ 위의 미분 등급 대수 구조가 유일하게
존재한다. 예를 들어 $a, b, c \in A$와 $s, t \in S$가 주어지면
다음과 같이 정의한다.
\begin{align*}
\text{d}(asbtc)
& =
\text{d}(a)sbtc + (-1)^{\deg(a)}a f_s b t c +
(-1)^{\deg(a) + \deg(s)} as\text{d}(b)tc \\
& + (-1)^{\deg(a) + \deg(s) + \deg(b)} asb f_t c +
(-1)^{\deg(a) + \deg(s) + \deg(b) + \deg(t)} asbt\text{d}(c)
\end{align*}
세부사항은 생략한다.

\begin{lemma}
\label{dga-lemma-K-flat-resolution}
$R$을 환이라 하자. $(B, \text{d})$를 미분 등급 $R$-대수라 하자.
다음 성질을 갖는 미분 등급 $R$-대수의 유사동형
$(A, \text{d}) \to (B, \text{d})$가 존재한다.
\begin{enumerate}
\item $A$는 $R$-가군의 복합체로서 K-평탄이다.
\item $A$는 자유 등급 $R$-대수이다.
\end{enumerate}
\end{lemma}

\begin{proof}
먼저 (1)과 (2)를 가지며 코호몰로지 위의 전사를 유도하는
$(A_0, \text{d}) \to (B, \text{d})$를 찾을 수 있다고 주장한다.
실제로 등급 집합 $S$를 택하고, 각 $s \in S$에 대해 차수가
$\deg(s)$인 동차 원소
$b_s \in \Ker(d : B \to B)$를 택하여, $H^*(B)$ 안의 동치류
$\overline{b}_s$들이 $H^*(B)$를 $R$-가군으로서 생성하게 한다.
이제 $A_0 = R\langle S \rangle$에 영 미분을 주고,
$s$를 $b_s$로 보내는 사상 $A_0 \to B$를 정의하면 된다.

\medskip\noindent
코호몰로지 위의 전사를 유도하는 $A_0 \to B$가 주어졌다고 하자.
귀납적으로 열
%% P03-STACKS-E589
$$
A_0 \to A_1 \to A_2 \to \ldots B
$$
을 구성한다. $A_n \to B$가 주어졌을 때 $S_n$을 등급 집합으로
잡고, 각 $s \in S_n$에 대해 차수가 $\deg(s) + 1$인 동차 원소
$a_s \in \Ker(\text{d} : A_n \to A_n)$를 택한다.
그 상 $\overline{a}_s$는 $H^*(A_n)$ 안의 동치류이고
$H^*(B)$에서 영으로 간다. 원소 $\overline{a}_s$들이
$\Ker(H^*(A_n) \to H^*(B))$를 $R$-가군으로서 생성하도록
$S_n$을 충분히 크게 택한다. 그런 다음
$$
A_{n + 1} = A_n\langle S_n \rangle
$$
으로 놓고, 위 논의에서처럼 미분을
$\text{d}(s) = a_s$로 정의한다.
%% P03-STACKS-E596
그러면 각 $(A_n, \text{d})$는 (1)과 (2)를 만족한다.
세부사항은 생략한다. $H^*(A_n) \to H^*(B)$는 $n = 0$일 때
전사였으므로 전사이다.

\medskip\noindent
$A = \colim A_n$이 자유 등급 $R$-대수임은 자명하다.
More on Algebra, Lemma \ref{more-algebra-lemma-colimit-K-flat}에 의해
이는 K-평탄이다. 사상 $H^*(A) \to H^*(B)$는 전사이면서 단사이므로
동형이다. 실제로 $H^*(A)$의 모든 원소는 어떤 $n$에 대한
$H^*(A_n)$의 원소에서 오며, 그것이 $H^*(B)$에서 사라지면
$H^*(A_{n + 1})$에서 사라지고 따라서 $H^*(A)$에서도 사라진다.
\end{proof}

\noindent
응용으로 Lemma \ref{dga-lemma-compose-tensor-functors-general-algebra}의
“올바른” 판을 증명한다.

\begin{lemma}
\label{dga-lemma-compose-tensor-functors-tor}
$R$을 환이라 하자. $(A, \text{d})$, $(B, \text{d})$,
$(C, \text{d})$를 미분 등급 $R$-대수라 하자.
$D(R)$에서 $A \otimes_R C$가 $A \otimes^\mathbf{L}_R C$를
나타낸다고 가정하자. $N$을 미분 등급 $(A, B)$-쌍가군,
$N'$을 미분 등급 $(B, C)$-쌍가군이라 하자. 그러면 합성
$$
\xymatrix{
D(A, \text{d}) \ar[rr]^{- \otimes_A^\mathbf{L} N} & &
D(B, \text{d}) \ar[rr]^{- \otimes_B^\mathbf{L} N'} & &
D(C, \text{d})
}
$$
은 어떤 미분 등급 $(A, C)$-쌍가군 $N''$에 대한
$- \otimes_A^\mathbf{L} N''$과 동형이다.
\end{lemma}

\begin{proof}
Lemma \ref{dga-lemma-K-flat-resolution}을 사용하여 $R$-가군의
복합체로서 K-평탄인 $B'$과 유사동형
$(B', \text{d}) \to (B, \text{d})$를 택한다.
Lemma \ref{dga-lemma-qis-equivalence}에 의해 함자
$-\otimes^\mathbf{L}_{B'} B : D(B', \text{d}) \to D(B, \text{d})$는
동치이고, 그 준역함자는 제한으로 주어진다. 제한은 $B$를
$(B, B')$-쌍가군으로 보았을 때의 함자
$- \otimes^\mathbf{L}_B B : D(B, \text{d}) \to D(B', \text{d})$와
표준적으로 동형이다. 따라서 다음 합성들에 대해 보조정리를
증명하면 충분하다.
$$
D(A) \to D(B) \to D(B'),\quad
D(B') \to D(B) \to D(C),\quad
D(A) \to D(B') \to D(C).
$$
첫 번째는 $B'$이 $R$-가군의 복합체로서 K-평탄이므로
Lemma \ref{dga-lemma-compose-tensor-functors}이다.
두 번째는
$B \otimes_B^\mathbf{L} N' = N' = B \otimes_B N'$이고 따라서
Lemma \ref{dga-lemma-compose-tensor-functors-general}을 적용할 수
있으므로 참이다. 이로써 $B$가 $R$-가군의 복합체로서
K-평탄인 경우로 환원된다.

\medskip\noindent
$B$가 $R$-가군의 복합체로서 K-평탄이라고 가정하자.
(\ref{dga-equation-plain-versus-derived-algebras})이 동형임을 보이는
것으로 충분하다. Lemma \ref{dga-lemma-compose-tensor-functors-general-algebra}를
보라. $L$이 성질 (P)를 갖는 미분 등급
$R$-가군인 유사동형 $L \to A$를 택하자.
그러면 $P = L \otimes_R B$가 미분 등급 $B$-가군으로서
성질 (P)를 가짐은 자명하다. 따라서 $P \to A \otimes_R B$가
유사동형
$$
P \otimes_B (B \otimes_R C)
\longrightarrow
(A \otimes_R B) \otimes_B (B \otimes_R C)
$$
을 유도함을 보여야 한다. 이를 다음과 같이 다시 쓸 수 있다.
%% P03-STACKS-E590
$$
P \otimes_R B \otimes_R C \longrightarrow A \otimes_R B \otimes_R C
$$
$B$가 $R$-가군의 복합체로서 K-평탄이므로, More on Algebra,
Lemma \ref{more-algebra-lemma-K-flat-quasi-isomorphism}에 의해
$$
P \otimes_R C \to A \otimes_R C
$$
가 유사동형임을 보이는 것으로 충분하다. 이는 정확히 우리의
가정이다.
\end{proof}

\noindent
다음 보조정리는 실제로 이 절에 속하지 않지만, 이를 둘 만한
자연스럽고 좋은 위치가 없는 듯하다.

\begin{lemma}
\label{dga-lemma-countable}
$(A, \text{d})$를 각 $i$에 대해 $H^i(A)$가 가산인
미분 등급 대수라 하자. $M$을 $D(A, \text{d})$의 객체라 하자.
다음은 서로 동치이다.
\begin{enumerate}
\item $D(A, \text{d})$에서 $E_n$이 콤팩트인
$M = \text{hocolim} E_n$이다.
\item 각 $i$에 대해 $H^i(M)$은 가산이다.
\end{enumerate}
\end{lemma}

\begin{proof}
(1)이 성립한다고 가정하자. Derived Categories, Lemma
\ref{derived-lemma-cohomology-of-hocolim}에 의해
$H^i(M) = \colim H^i(E_n)$이다. 따라서 각 $n$에 대해
$H^i(E_n)$이 가산임을 증명하면 충분하다. Proposition
\ref{dga-proposition-compact}에 의해 $E_n$은 $D(A, \text{d})$에서,
미분 등급 부분가군들에 의한 유한 필트레이션 $F_\bullet$을 갖고
$F_jP/F_{j - 1}P$가 $A$의 시프트들의 유한 직합인 미분 등급
가군 $P$의 직합인자와 동형이다. 가정에 의해 군
$H^i(F_jP/F_{j - 1}P)$는 가산이다. 필트레이션의 길이에 관해
귀납하고 코호몰로지 긴 완전열을 사용하면 (2)가 성립한다.
흥미로운 함의는 그 역방향이다.

\medskip\noindent
포함 사상 $A' \to A$가 코호몰로지 위의 동형을 정의하도록 하는
가산 미분 등급 부분대수 $A' \subset A$가 존재한다고 주장한다.
$A'$를 구성하기 위해 가산 미분 등급 부분대수들
$$
A_1 \subset A_2 \subset A_3 \subset \ldots
$$
을 다음과 같이 택한다. (a) $H^i(A_1) \to H^i(A)$는 전사이고,
(b) $n > 1$이면 사상 $H^i(A_{n - 1}) \to H^i(A_n)$의 핵은
사상 $H^i(A_{n - 1}) \to H^i(A)$의 핵과 같다.
$A_1$을 구성하려면 $A$의 코호몰로지를 (환으로서 또는 등급
아벨군으로서) 생성하는 가산 쌍대사슬 모음 $S \subset A$를
아무거나 택하고, $S$가 생성하는 $A$의 미분 등급 부분대수를
$A_1$이라 한다.
%% P03-STACKS-E591
$A_{n - 1}$에서 $A_n$을 구성하려면, $H^i(A)$에서 영으로 가는
각 쌍대사슬 $a \in A_{n - 1}^i$에 대해
$\text{d}(s_a) = a$인 $s_a \in A^{i - 1}$를 택하고,
$A_{n - 1}$과 원소들 $s_a$가 생성하는 $A$의 미분 등급
부분대수를 $A_n$이라 한다. 마지막으로 $A' = \bigcup A_n$으로
놓는다.

\medskip\noindent
Lemma \ref{dga-lemma-qis-equivalence}에 의해 제한 사상
$D(A, \text{d}) \to D(A', \text{d})$,
$M \mapsto M_{A'}$는 동치이다. $M$과 $M_{A'}$의 코호몰로지
군들이 같으므로, $(A', \text{d})$에 대해 함의
(2) $\Rightarrow$ (1)을 증명하면 충분하다.

\medskip\noindent
$A$가 가산이라고 가정하자. 위와 정확히 같은 유형의 논증으로,
$D(A, \text{d})$의 $M$에 대해 다음 두 조건이 서로 동치임을
알 수 있다. 각 $i$에 대해 $H^i(M)$이 가산이다. 그리고 $M$은
가산 미분 등급 가군으로 나타낼 수 있다. 따라서 함의
(2) $\Rightarrow$ (1)을 증명하기 위해 다음 문단의 상황으로
환원된다.

\medskip\noindent
$A$가 가산이고 $M$이 $A$ 위의 가산 미분 등급 가군이라고
가정하자. 다음을 만족하는 미분 등급 $A$-가군의 준동형
$P \to M$이 존재한다고 주장한다.
\begin{enumerate}
\item $P \to M$은 유사동형이다.
\item $P$는 성질 (P)를 갖는다.
\item $P$는 가산이다.
\end{enumerate}
Lemma \ref{dga-lemma-resolve}의 P-분해 구성 증명을 살펴보면,
가산 미분 등급 가군의 상황에서 Lemma
\ref{dga-lemma-good-quotient}를 증명할 수 있음을 보이는 것으로
충분하다. 이는 그 증명에서 즉시 따른다.

\medskip\noindent
%% P03-STACKS-E592
$A$가 가산이고 $M$이 성질 (P)를 갖는 가산 미분 등급
가군이라고 가정하자. 다음과 같은 미분 등급 부분가군들에 의한
필트레이션을 택하자.
$$
0 = F_{-1}P \subset F_0P \subset F_1P \subset \ldots \subset P
$$
이 필트레이션은 다음을 만족한다.
\begin{enumerate}
\item $P = \bigcup F_pP$이다.
\item $F_iP \to F_{i + 1}P$는 허용 단사사상이다.
\item 어떤 집합들 $J_i$와 정수들 $k_j$에 대한 미분 등급
가군의 동형
$F_iP/F_{i - 1}P \to \bigoplus_{j \in J_i} A[k_j]$가 존재한다.
\end{enumerate}
물론 각 $i$에 대해 $J_i$는 가산이다. 각 $i$와 $j \in J_i$에
대해, $F_iP/F_{i - 1}P$ 안에서 $j$에 대응하는 직합인자를
생성하는 상을 갖는 차수 $k_j$의 원소 $x_{i, j} \in F_iP$를 택한다.

\medskip\noindent
%% P03-STACKS-E593
%% P03-STACKS-E594
주장: $n$과 유한 부분집합 $S_i \subset J_i$,
$i = 1, \ldots, n$이 주어지면, $P$의 미분 등급 부분가군이 되는
$P' = \bigoplus_{i \leq n} \bigoplus_{j \in T_i} Ax_{i, j}$가
존재하도록 하는 유한 부분집합 $S_i \subset T_i \subset J_i$,
$i = 1, \ldots, n$이 존재한다. 이는 Lemma
\ref{dga-lemma-factor-through-nicer}의 증명에서 보였지만 직접 보이기도
쉽다. 원소들 $x_{i, j}$는 오른쪽 $A$-가군으로서 $P$를 자유롭게
생성한다. $P$의 구조로부터
%% P03-STACKS-E595
$$
\text{d}(x_{i, j}) = \sum\nolimits_{i' < i} x_{i', j'}a_{i', j'}
$$
이고, 물론 이 합은 유한하다.
따라서 $S_0, \ldots, S_n$이 주어지면, 먼저 모든 $j \in S_n$에 대해
$\text{d}(x_{n, j}) \in \bigoplus_{i' < n, j' \in S'_{i'}} x_{i', j'}A$가
되도록 $S_0 \subset S'_0, \ldots, S_{n - 1} \subset S'_{n - 1}$을
택할 수 있다. 그런 다음 $n$에 관한 귀납법으로,
$\bigoplus_{i' < n, j' \in T_{i'}} x_{i', j'}A$가 미분 등급
$A$-부분가군이 되도록
$S'_0 \subset T_0, \ldots, S'_{n - 1} \subset T_{n - 1}$을
택할 수 있다. $T_n = S_n$으로 놓으면
$P' = \bigoplus_{i \leq n, j \in T_i} x_{i, j}A$가 원하는
가군임을 알 수 있다.

\medskip\noindent
이 주장으로부터 $P = \bigcup P'_n$이 위와 같은 $P'_n$들의
가산 증가 합집합임은 자명하다. 구성상 각 $P'_n$은 성질 (P)를
갖는 미분 등급 가군이고, 필트레이션은 유한하며 연이은 몫들은
$A$의 시프트들의 유한 직합이다. 따라서 $P'_n$은
$D(A, \text{d})$의 콤팩트 객체를 정의한다. 예를 들어
Proposition \ref{dga-proposition-compact}을 보라.
Lemma \ref{dga-lemma-homotopy-colimit}에 의해 $D(A, \text{d})$에서
$P = \text{hocolim} P'_n$이므로 함의 (2) $\Rightarrow$ (1)의
증명이 끝난다.
\end{proof}
